% Môn: Vật lý
% Lớp: 12
% Chương: Chương III. Từ trường
% Bài: L12C3 Bài 15. Lực từ tác dụng lên dây dẫn mang dòng điện. Cảm ứng từ
% Số câu: 328
% App đọc file này trực tiếp — sửa rồi Commit trên GitHub, bấm Đồng bộ GitHub .tex

% L12C3 Bài 15. Lực từ tác dụng lên dây dẫn mang dòng điện. Cảm ứng từ

% ===== Câu 1 =====
% ID: L12C3B15-01-DS
% Mức: TH
\dangbt{Bài toán tổng hợp lực từ gắn với đồ thị, thí nghiệm và thực tiễn}
\begin{ex}
% ID: L12C3B15-01-DS
% Mức: TH
Xét Bài toán tổng hợp lực từ gắn với đồ thị, thí nghiệm và thực tiễn (Tình huống 1). Hãy xác định tính đúng, sai của bốn phát biểu sau.
\choiceTF
{\True Với $B$ và l không đổi, đồ thị $F$ theo $I$ là đường thẳng qua gốc có hệ số góc Bl sin(alpha).}
{Đồ thị $F$ theo $I$ luôn là đường cong và không phụ thuộc B.}
{\True Khi giải cần xác định đúng phương, chiều, điều kiện áp dụng và đơn vị.}
{Kết quả không phụ thuộc dữ kiện và có thể bỏ qua điều kiện.}
\loigiai{
\begin{itemchoice}
\item (Đúng) Với $B$ và l không đổi, đồ thị $F$ theo $I$ là đường thẳng qua gốc có hệ số góc Bl sin(alpha). (Vì): lpha
\item (Sai) Đồ thị $F$ theo $I$ luôn là đường cong và không phụ thuộc B.
\item (Đúng) Khi giải cần xác định đúng phương, chiều, điều kiện áp dụng và đơn vị.
\item (Sai) Kết quả không phụ thuộc dữ kiện và có thể bỏ qua điều kiện.
\end{itemchoice}
}
\end{ex}

% ===== Câu 2 =====
% ID: L12C3B15-01-TL
% Mức: TH
\begin{bt}
% ID: L12C3B15-01-TL
% Mức: TH
Trình bày kiến thức cốt lõi của Bài toán tổng hợp lực từ gắn với đồ thị, thí nghiệm và thực tiễn và nêu điều kiện áp dụng.
\loigiai{
Cần nêu được: Với $B$ và l không đổi, đồ thị $F$ theo $I$ là đường thẳng qua gốc có hệ số góc Bl sin(alpha). Khi vận dụng phải xác định đúng dữ kiện, phương chiều nếu có, điều kiện áp dụng, hệ đơn vị và kiểm tra tính hợp lí. Nhận định “Đồ thị $F$ theo $I$ luôn là đường cong và không phụ thuộc B.” sai vì trái với kiến thức nêu trên.
}
\end{bt}

% ===== Câu 3 =====
% ID: L12C3B15-01-TLN
% Mức: TH
\begin{ex}
% ID: L12C3B15-01-TLN
% Mức: TH
Trong bốn phát biểu về Bài toán tổng hợp lực từ gắn với đồ thị, thí nghiệm và thực tiễn, có bao nhiêu phát biểu đúng? (Chỉ ghi một số nguyên.) (1) Với $B$ và l không đổi, đồ thị $F$ theo $I$ là đường thẳng qua gốc có hệ số góc Bl sin(alpha). (2) Cần kiểm tra điều kiện áp dụng. (3) Đồ thị $F$ theo $I$ luôn là đường cong và không phụ thuộc B. (4) Có thể bỏ qua đơn vị của kết quả.
\shortans{2}
\loigiai{
Đối chiếu từng phát biểu với kiến thức: Với $B$ và l không đổi, đồ thị $F$ theo $I$ là đường thẳng qua gốc có hệ số góc Bl sin(alpha). Số phát biểu đúng là 2.
}
\end{ex}

% ===== Câu 4 =====
% ID: L12C3B15-01-TN
% Mức: NB
\begin{ex}
% ID: L12C3B15-01-TN
% Mức: NB
Phát biểu nào sau đây đúng về Bài toán tổng hợp lực từ gắn với đồ thị, thí nghiệm và thực tiễn?
\choice
{Đồ thị $F$ theo $I$ luôn là đường cong và không phụ thuộc B.}
{Có thể bỏ qua phương, chiều và đơn vị của các đại lượng trong mọi trường hợp.}
{Kết quả không phụ thuộc dữ kiện và điều kiện áp dụng.}
{\True Với $B$ và l không đổi, đồ thị $F$ theo $I$ là đường thẳng qua gốc có hệ số góc Bl sin(alpha).}
\loigiai{
Với $B$ và l không đổi, đồ thị $F$ theo $I$ là đường thẳng qua gốc có hệ số góc Bl sin(alpha). Vì vậy phương án $D$ đúng; các phương án còn lại trái với kiến thức hoặc điều kiện áp dụng.
}
\end{ex}

% ===== Câu 5 =====
% ID: L12C3B15-02-DS
% Mức: TH
\begin{ex}
% ID: L12C3B15-02-DS
% Mức: TH
Xét Bài toán tổng hợp lực từ gắn với đồ thị, thí nghiệm và thực tiễn (Tình huống 2). Hãy xác định tính đúng, sai của bốn phát biểu sau.
\choiceTF
{Đồ thị $F$ theo $I$ luôn là đường cong và không phụ thuộc B.}
{\True Với $B$ và l không đổi, đồ thị $F$ theo $I$ là đường thẳng qua gốc có hệ số góc Bl sin(alpha).}
{Kết quả không phụ thuộc dữ kiện và có thể bỏ qua điều kiện.}
{\True Khi giải cần xác định đúng phương, chiều, điều kiện áp dụng và đơn vị.}
\loigiai{
\begin{itemchoice}
\item (Sai) Đồ thị $F$ theo $I$ luôn là đường cong và không phụ thuộc B. (Vì): lpha
\item (Đúng) Với $B$ và l không đổi, đồ thị $F$ theo $I$ là đường thẳng qua gốc có hệ số góc Bl sin(alpha).
\item (Sai) Kết quả không phụ thuộc dữ kiện và có thể bỏ qua điều kiện.
\item (Đúng) Khi giải cần xác định đúng phương, chiều, điều kiện áp dụng và đơn vị.
\end{itemchoice}
}
\end{ex}

% ===== Câu 6 =====
% ID: L12C3B15-02-TL
% Mức: TH
\begin{bt}
% ID: L12C3B15-02-TL
% Mức: TH
Giải thích vì sao nhận định sau đúng: “Với $B$ và l không đổi, đồ thị $F$ theo $I$ là đường thẳng qua gốc có hệ số góc Bl sin(alpha).”
\loigiai{
Cần nêu được: Với $B$ và l không đổi, đồ thị $F$ theo $I$ là đường thẳng qua gốc có hệ số góc Bl sin(alpha). Khi vận dụng phải xác định đúng dữ kiện, phương chiều nếu có, điều kiện áp dụng, hệ đơn vị và kiểm tra tính hợp lí. Nhận định “Đồ thị $F$ theo $I$ luôn là đường cong và không phụ thuộc B.” sai vì trái với kiến thức nêu trên.
}
\end{bt}

% ===== Câu 7 =====
% ID: L12C3B15-02-TLN
% Mức: TH
\begin{ex}
% ID: L12C3B15-02-TLN
% Mức: TH
Trong bốn phát biểu về Bài toán tổng hợp lực từ gắn với đồ thị, thí nghiệm và thực tiễn, có bao nhiêu phát biểu đúng? (Chỉ ghi một số nguyên.) (1) Đồ thị $F$ theo $I$ luôn là đường cong và không phụ thuộc B. (2) Với $B$ và l không đổi, đồ thị $F$ theo $I$ là đường thẳng qua gốc có hệ số góc Bl sin(alpha). (3) Đồ thị $F$ theo $I$ luôn là đường cong và không phụ thuộc B. (4) Có thể bỏ qua đơn vị của kết quả.
\shortans{1}
\loigiai{
Đối chiếu từng phát biểu với kiến thức: Với $B$ và l không đổi, đồ thị $F$ theo $I$ là đường thẳng qua gốc có hệ số góc Bl sin(alpha). Số phát biểu đúng là 1.
}
\end{ex}

% ===== Câu 8 =====
% ID: L12C3B15-02-TN
% Mức: NB
\begin{ex}
% ID: L12C3B15-02-TN
% Mức: NB
Trong nội dung Bài toán tổng hợp lực từ gắn với đồ thị, thí nghiệm và thực tiễn?
\choice
{\True Với $B$ và l không đổi, đồ thị $F$ theo $I$ là đường thẳng qua gốc có hệ số góc Bl sin(alpha).}
{Đồ thị $F$ theo $I$ luôn là đường cong và không phụ thuộc B.}
{Có thể bỏ qua phương, chiều và đơn vị của các đại lượng trong mọi trường hợp.}
{Kết quả không phụ thuộc dữ kiện và điều kiện áp dụng.}
\loigiai{
Với $B$ và l không đổi, đồ thị $F$ theo $I$ là đường thẳng qua gốc có hệ số góc Bl sin(alpha). Vì vậy phương án $A$ đúng; các phương án còn lại trái với kiến thức hoặc điều kiện áp dụng.
}
\end{ex}

% ===== Câu 9 =====
% ID: L12C3B15-03-DS
% Mức: VD
\begin{ex}
% ID: L12C3B15-03-DS
% Mức: VD
Xét Bài toán tổng hợp lực từ gắn với đồ thị, thí nghiệm và thực tiễn (Tình huống 3). Hãy xác định tính đúng, sai của bốn phát biểu sau.
\choiceTF
{\True Khi giải cần xác định đúng phương, chiều, điều kiện áp dụng và đơn vị.}
{Kết quả không phụ thuộc dữ kiện và có thể bỏ qua điều kiện.}
{\True Với $B$ và l không đổi, đồ thị $F$ theo $I$ là đường thẳng qua gốc có hệ số góc Bl sin(alpha).}
{Đồ thị $F$ theo $I$ luôn là đường cong và không phụ thuộc B.}
\loigiai{
\begin{itemchoice}
\item (Đúng) Khi giải cần xác định đúng phương, chiều, điều kiện áp dụng và đơn vị. (Vì): lpha
\item (Sai) Kết quả không phụ thuộc dữ kiện và có thể bỏ qua điều kiện.
\item (Đúng) Với $B$ và l không đổi, đồ thị $F$ theo $I$ là đường thẳng qua gốc có hệ số góc Bl sin(alpha).
\item (Sai) Đồ thị $F$ theo $I$ luôn là đường cong và không phụ thuộc B.
\end{itemchoice}
}
\end{ex}

% ===== Câu 10 =====
% ID: L12C3B15-03-TL
% Mức: VD
\begin{bt}
% ID: L12C3B15-03-TL
% Mức: VD
Phân tích sai lầm trong nhận định: “Đồ thị $F$ theo $I$ luôn là đường cong và không phụ thuộc B.”
\loigiai{
Cần nêu được: Với $B$ và l không đổi, đồ thị $F$ theo $I$ là đường thẳng qua gốc có hệ số góc Bl sin(alpha). Khi vận dụng phải xác định đúng dữ kiện, phương chiều nếu có, điều kiện áp dụng, hệ đơn vị và kiểm tra tính hợp lí. Nhận định “Đồ thị $F$ theo $I$ luôn là đường cong và không phụ thuộc B.” sai vì trái với kiến thức nêu trên.
}
\end{bt}

% ===== Câu 11 =====
% ID: L12C3B15-03-TLN
% Mức: TH
\begin{ex}
% ID: L12C3B15-03-TLN
% Mức: TH
Trong bốn phát biểu về Bài toán tổng hợp lực từ gắn với đồ thị, thí nghiệm và thực tiễn, có bao nhiêu phát biểu đúng? (Chỉ ghi một số nguyên.) (1) Với $B$ và l không đổi, đồ thị $F$ theo $I$ là đường thẳng qua gốc có hệ số góc Bl sin(alpha). (2) Cần kiểm tra điều kiện áp dụng. (3) Đồ thị $F$ theo $I$ luôn là đường cong và không phụ thuộc B. (4) Có thể bỏ qua đơn vị của kết quả.
\shortans{2}
\loigiai{
Đối chiếu từng phát biểu với kiến thức: Với $B$ và l không đổi, đồ thị $F$ theo $I$ là đường thẳng qua gốc có hệ số góc Bl sin(alpha). Số phát biểu đúng là 2.
}
\end{ex}

% ===== Câu 12 =====
% ID: L12C3B15-03-TN
% Mức: NB
\begin{ex}
% ID: L12C3B15-03-TN
% Mức: NB
Tình huống 3: chọn kết luận chính xác về Bài toán tổng hợp lực từ gắn với đồ thị, thí nghiệm và thực tiễn?
\choice
{Kết quả không phụ thuộc dữ kiện và điều kiện áp dụng.}
{\True Với $B$ và l không đổi, đồ thị $F$ theo $I$ là đường thẳng qua gốc có hệ số góc Bl sin(alpha).}
{Đồ thị $F$ theo $I$ luôn là đường cong và không phụ thuộc B.}
{Có thể bỏ qua phương, chiều và đơn vị của các đại lượng trong mọi trường hợp.}
\loigiai{
Với $B$ và l không đổi, đồ thị $F$ theo $I$ là đường thẳng qua gốc có hệ số góc Bl sin(alpha). Vì vậy phương án $B$ đúng; các phương án còn lại trái với kiến thức hoặc điều kiện áp dụng.
}
\end{ex}

% ===== Câu 13 =====
% ID: L12C3B15-04-DS
% Mức: VD
\begin{ex}
% ID: L12C3B15-04-DS
% Mức: VD
Xét Bài toán tổng hợp lực từ gắn với đồ thị, thí nghiệm và thực tiễn (Tình huống 4). Hãy xác định tính đúng, sai của bốn phát biểu sau.
\choiceTF
{Kết quả không phụ thuộc dữ kiện và có thể bỏ qua điều kiện.}
{\True Khi giải cần xác định đúng phương, chiều, điều kiện áp dụng và đơn vị.}
{Đồ thị $F$ theo $I$ luôn là đường cong và không phụ thuộc B.}
{\True Với $B$ và l không đổi, đồ thị $F$ theo $I$ là đường thẳng qua gốc có hệ số góc Bl sin(alpha).}
\loigiai{
\begin{itemchoice}
\item (Sai) Kết quả không phụ thuộc dữ kiện và có thể bỏ qua điều kiện. (Vì): lpha
\item (Đúng) Khi giải cần xác định đúng phương, chiều, điều kiện áp dụng và đơn vị.
\item (Sai) Đồ thị $F$ theo $I$ luôn là đường cong và không phụ thuộc B.
\item (Đúng) Với $B$ và l không đổi, đồ thị $F$ theo $I$ là đường thẳng qua gốc có hệ số góc Bl sin(alpha).
\end{itemchoice}
}
\end{ex}

% ===== Câu 14 =====
% ID: L12C3B15-04-TL
% Mức: VD
\begin{bt}
% ID: L12C3B15-04-TL
% Mức: VD
Đề xuất các bước giải một bài toán thuộc dạng Bài toán tổng hợp lực từ gắn với đồ thị, thí nghiệm và thực tiễn.
\loigiai{
Cần nêu được: Với $B$ và l không đổi, đồ thị $F$ theo $I$ là đường thẳng qua gốc có hệ số góc Bl sin(alpha). Khi vận dụng phải xác định đúng dữ kiện, phương chiều nếu có, điều kiện áp dụng, hệ đơn vị và kiểm tra tính hợp lí. Nhận định “Đồ thị $F$ theo $I$ luôn là đường cong và không phụ thuộc B.” sai vì trái với kiến thức nêu trên.
}
\end{bt}

% ===== Câu 15 =====
% ID: L12C3B15-04-TLN
% Mức: VD
\begin{ex}
% ID: L12C3B15-04-TLN
% Mức: VD
Trong bốn phát biểu về Bài toán tổng hợp lực từ gắn với đồ thị, thí nghiệm và thực tiễn, có bao nhiêu phát biểu đúng? (Chỉ ghi một số nguyên.) (1) Đồ thị $F$ theo $I$ luôn là đường cong và không phụ thuộc B. (2) Với $B$ và l không đổi, đồ thị $F$ theo $I$ là đường thẳng qua gốc có hệ số góc Bl sin(alpha). (3) Đồ thị $F$ theo $I$ luôn là đường cong và không phụ thuộc B. (4) Có thể bỏ qua đơn vị của kết quả.
\shortans{2}
\loigiai{
Đối chiếu từng phát biểu với kiến thức: Với $B$ và l không đổi, đồ thị $F$ theo $I$ là đường thẳng qua gốc có hệ số góc Bl sin(alpha). Số phát biểu đúng là 2.
}
\end{ex}

% ===== Câu 16 =====
% ID: L12C3B15-04-TN
% Mức: TH
\begin{ex}
% ID: L12C3B15-04-TN
% Mức: TH
Nội dung nào mô tả đúng nhất Bài toán tổng hợp lực từ gắn với đồ thị, thí nghiệm và thực tiễn?
\choice
{Có thể bỏ qua phương, chiều và đơn vị của các đại lượng trong mọi trường hợp.}
{Kết quả không phụ thuộc dữ kiện và điều kiện áp dụng.}
{\True Với $B$ và l không đổi, đồ thị $F$ theo $I$ là đường thẳng qua gốc có hệ số góc Bl sin(alpha).}
{Đồ thị $F$ theo $I$ luôn là đường cong và không phụ thuộc B.}
\loigiai{
Với $B$ và l không đổi, đồ thị $F$ theo $I$ là đường thẳng qua gốc có hệ số góc Bl sin(alpha). Vì vậy phương án $C$ đúng; các phương án còn lại trái với kiến thức hoặc điều kiện áp dụng.
}
\end{ex}

% ===== Câu 17 =====
% ID: L12C3B15-05-DS
% Mức: VD
\begin{ex}
% ID: L12C3B15-05-DS
% Mức: VD
Xét Bài toán tổng hợp lực từ gắn với đồ thị, thí nghiệm và thực tiễn (Tình huống 5). Hãy xác định tính đúng, sai của bốn phát biểu sau.
\choiceTF
{\True Với $B$ và l không đổi, đồ thị $F$ theo $I$ là đường thẳng qua gốc có hệ số góc Bl sin(alpha).}
{Kết quả không phụ thuộc dữ kiện và có thể bỏ qua điều kiện.}
{Đồ thị $F$ theo $I$ luôn là đường cong và không phụ thuộc B.}
{\True Khi giải cần xác định đúng phương, chiều, điều kiện áp dụng và đơn vị.}
\loigiai{
\begin{itemchoice}
\item (Đúng) Với $B$ và l không đổi, đồ thị $F$ theo $I$ là đường thẳng qua gốc có hệ số góc Bl sin(alpha). (Vì): lpha
\item (Sai) Kết quả không phụ thuộc dữ kiện và có thể bỏ qua điều kiện.
\item (Sai) Đồ thị $F$ theo $I$ luôn là đường cong và không phụ thuộc B.
\item (Đúng) Khi giải cần xác định đúng phương, chiều, điều kiện áp dụng và đơn vị.
\end{itemchoice}
}
\end{ex}

% ===== Câu 18 =====
% ID: L12C3B15-05-TL
% Mức: VD
\begin{bt}
% ID: L12C3B15-05-TL
% Mức: VD
Nêu một tình huống thực tiễn liên quan đến Bài toán tổng hợp lực từ gắn với đồ thị, thí nghiệm và thực tiễn và giải thích bằng kiến thức Vật lí.
\loigiai{
Cần nêu được: Với $B$ và l không đổi, đồ thị $F$ theo $I$ là đường thẳng qua gốc có hệ số góc Bl sin(alpha). Khi vận dụng phải xác định đúng dữ kiện, phương chiều nếu có, điều kiện áp dụng, hệ đơn vị và kiểm tra tính hợp lí. Nhận định “Đồ thị $F$ theo $I$ luôn là đường cong và không phụ thuộc B.” sai vì trái với kiến thức nêu trên.
}
\end{bt}

% ===== Câu 19 =====
% ID: L12C3B15-05-TLN
% Mức: VD
\begin{ex}
% ID: L12C3B15-05-TLN
% Mức: VD
Trong bốn phát biểu về Bài toán tổng hợp lực từ gắn với đồ thị, thí nghiệm và thực tiễn, có bao nhiêu phát biểu đúng? (Chỉ ghi một số nguyên.) (1) Với $B$ và l không đổi, đồ thị $F$ theo $I$ là đường thẳng qua gốc có hệ số góc Bl sin(alpha). (2) Cần kiểm tra điều kiện áp dụng. (3) Đồ thị $F$ theo $I$ luôn là đường cong và không phụ thuộc B. (4) Có thể bỏ qua đơn vị của kết quả.
\shortans{2}
\loigiai{
Đối chiếu từng phát biểu với kiến thức: Với $B$ và l không đổi, đồ thị $F$ theo $I$ là đường thẳng qua gốc có hệ số góc Bl sin(alpha). Số phát biểu đúng là 2.
}
\end{ex}

% ===== Câu 20 =====
% ID: L12C3B15-05-TN
% Mức: TH
\begin{ex}
% ID: L12C3B15-05-TN
% Mức: TH
Khi phân tích Bài toán tổng hợp lực từ gắn với đồ thị, thí nghiệm và thực tiễn?
\choice
{Đồ thị $F$ theo $I$ luôn là đường cong và không phụ thuộc B.}
{Có thể bỏ qua phương, chiều và đơn vị của các đại lượng trong mọi trường hợp.}
{Kết quả không phụ thuộc dữ kiện và điều kiện áp dụng.}
{\True Với $B$ và l không đổi, đồ thị $F$ theo $I$ là đường thẳng qua gốc có hệ số góc Bl sin(alpha).}
\loigiai{
Với $B$ và l không đổi, đồ thị $F$ theo $I$ là đường thẳng qua gốc có hệ số góc Bl sin(alpha). Vì vậy phương án $D$ đúng; các phương án còn lại trái với kiến thức hoặc điều kiện áp dụng.
}
\end{ex}

% ===== Câu 21 =====
% ID: L12C3B15-06-DS
% Mức: TH
\dangbt{Cân bằng và chuyển động của dây dẫn mang dòng điện trong từ trường}
\begin{ex}
% ID: L12C3B15-06-DS
% Mức: TH
Xét Cân bằng và chuyển động của dây dẫn mang dòng điện trong từ trường (Tình huống 1). Hãy xác định tính đúng, sai của bốn phát biểu sau.
\choiceTF
{\True Dây cân bằng khi tổng vectơ các lực, gồm lực từ và các lực cơ học, bằng vectơ không.}
{Dây cân bằng chỉ cần lực từ khác không, không cần xét trọng lực hay lực căng.}
{\True Khi giải cần xác định đúng phương, chiều, điều kiện áp dụng và đơn vị.}
{Kết quả không phụ thuộc dữ kiện và có thể bỏ qua điều kiện.}
\loigiai{
\begin{itemchoice}
\item (Đúng) Dây cân bằng khi tổng vectơ các lực, gồm lực từ và các lực cơ học, bằng vectơ không. (Vì): ây cân bằng khi tổng vectơ các lực, gồm lực từ và các lực cơ học, bằng vectơ không. b) Sai: Dây cân bằng chỉ cần lực từ khác không, không cần xét trọng lực hay lực căng. c) Đúng: Khi giải cần xác định đúng phương, chiều, điều kiện áp dụng và đơn vị. d) Sai: Kết quả không phụ thuộc dữ kiện và có thể bỏ qua điều kiện. Kiến thức cốt lõi: Dây cân bằng khi tổng vectơ các lực, gồm lực từ và các lực cơ học, bằng vectơ không
\item (Sai) Dây cân bằng chỉ cần lực từ khác không, không cần xét trọng lực hay lực căng.
\item (Đúng) Khi giải cần xác định đúng phương, chiều, điều kiện áp dụng và đơn vị.
\item (Sai) Kết quả không phụ thuộc dữ kiện và có thể bỏ qua điều kiện.
\end{itemchoice}
}
\end{ex}

% ===== Câu 22 =====
% ID: L12C3B15-06-TL
% Mức: VDC
\dangbt{Bài toán tổng hợp lực từ gắn với đồ thị, thí nghiệm và thực tiễn}
\begin{bt}
% ID: L12C3B15-06-TL
% Mức: VDC
So sánh nhận định đúng “Với $B$ và l không đổi, đồ thị $F$ theo $I$ là đường thẳng qua gốc có hệ số góc Bl sin(alpha).” với nhận định sai “Đồ thị $F$ theo $I$ luôn là đường cong và không phụ thuộc B.”; chỉ rõ căn cứ Vật lí.
\loigiai{
Cần nêu được: Với $B$ và l không đổi, đồ thị $F$ theo $I$ là đường thẳng qua gốc có hệ số góc Bl sin(alpha). Khi vận dụng phải xác định đúng dữ kiện, phương chiều nếu có, điều kiện áp dụng, hệ đơn vị và kiểm tra tính hợp lí. Nhận định “Đồ thị $F$ theo $I$ luôn là đường cong và không phụ thuộc B.” sai vì trái với kiến thức nêu trên.
}
\end{bt}

% ===== Câu 23 =====
% ID: L12C3B15-06-TLN
% Mức: VDC
\begin{ex}
% ID: L12C3B15-06-TLN
% Mức: VDC
Trong bốn phát biểu về Bài toán tổng hợp lực từ gắn với đồ thị, thí nghiệm và thực tiễn, có bao nhiêu phát biểu đúng? (Chỉ ghi một số nguyên.) (1) Đồ thị $F$ theo $I$ luôn là đường cong và không phụ thuộc B. (2) Với $B$ và l không đổi, đồ thị $F$ theo $I$ là đường thẳng qua gốc có hệ số góc Bl sin(alpha). (3) Phải dùng đúng định luật cho tình huống. (4) Với $B$ và l không đổi, đồ thị $F$ theo $I$ là đường thẳng qua gốc có hệ số góc Bl sin(alpha).
\shortans{3}
\loigiai{
Đối chiếu từng phát biểu với kiến thức: Với $B$ và l không đổi, đồ thị $F$ theo $I$ là đường thẳng qua gốc có hệ số góc Bl sin(alpha). Số phát biểu đúng là 3.
}
\end{ex}

% ===== Câu 24 =====
% ID: L12C3B15-06-TN
% Mức: TH
\begin{ex}
% ID: L12C3B15-06-TN
% Mức: TH
Một học sinh ôn tập Bài toán tổng hợp lực từ gắn với đồ thị, thí nghiệm và thực tiễn?
\choice
{\True Với $B$ và l không đổi, đồ thị $F$ theo $I$ là đường thẳng qua gốc có hệ số góc Bl sin(alpha).}
{Đồ thị $F$ theo $I$ luôn là đường cong và không phụ thuộc B.}
{Có thể bỏ qua phương, chiều và đơn vị của các đại lượng trong mọi trường hợp.}
{Kết quả không phụ thuộc dữ kiện và điều kiện áp dụng.}
\loigiai{
Với $B$ và l không đổi, đồ thị $F$ theo $I$ là đường thẳng qua gốc có hệ số góc Bl sin(alpha). Vì vậy phương án $A$ đúng; các phương án còn lại trái với kiến thức hoặc điều kiện áp dụng.
}
\end{ex}

% ===== Câu 25 =====
% ID: L12C3B15-07-DS
% Mức: TH
\dangbt{Cân bằng và chuyển động của dây dẫn mang dòng điện trong từ trường}
\begin{ex}
% ID: L12C3B15-07-DS
% Mức: TH
Xét Cân bằng và chuyển động của dây dẫn mang dòng điện trong từ trường (Tình huống 2). Hãy xác định tính đúng, sai của bốn phát biểu sau.
\choiceTF
{Dây cân bằng chỉ cần lực từ khác không, không cần xét trọng lực hay lực căng.}
{\True Dây cân bằng khi tổng vectơ các lực, gồm lực từ và các lực cơ học, bằng vectơ không.}
{Kết quả không phụ thuộc dữ kiện và có thể bỏ qua điều kiện.}
{\True Khi giải cần xác định đúng phương, chiều, điều kiện áp dụng và đơn vị.}
\loigiai{
\begin{itemchoice}
\item (Sai) Dây cân bằng chỉ cần lực từ khác không, không cần xét trọng lực hay lực căng. (Vì): Dây cân bằng chỉ cần lực từ khác không, không cần xét trọng lực hay lực căng. b) Đúng: Dây cân bằng khi tổng vectơ các lực, gồm lực từ và các lực cơ học, bằng vectơ không. c) Sai: Kết quả không phụ thuộc dữ kiện và có thể bỏ qua điều kiện. d) Đúng: Khi giải cần xác định đúng phương, chiều, điều kiện áp dụng và đơn vị. Kiến thức cốt lõi: Dây cân bằng khi tổng vectơ các lực, gồm lực từ và các lực cơ học, bằng vectơ không
\item (Đúng) Dây cân bằng khi tổng vectơ các lực, gồm lực từ và các lực cơ học, bằng vectơ không.
\item (Sai) Kết quả không phụ thuộc dữ kiện và có thể bỏ qua điều kiện.
\item (Đúng) Khi giải cần xác định đúng phương, chiều, điều kiện áp dụng và đơn vị.
\end{itemchoice}
}
\end{ex}

% ===== Câu 26 =====
% ID: L12C3B15-07-TL
% Mức: TH
\begin{bt}
% ID: L12C3B15-07-TL
% Mức: TH
Trình bày kiến thức cốt lõi của Cân bằng và chuyển động của dây dẫn mang dòng điện trong từ trường và nêu điều kiện áp dụng.
\loigiai{
Cần nêu được: Dây cân bằng khi tổng vectơ các lực, gồm lực từ và các lực cơ học, bằng vectơ không. Khi vận dụng phải xác định đúng dữ kiện, phương chiều nếu có, điều kiện áp dụng, hệ đơn vị và kiểm tra tính hợp lí. Nhận định “Dây cân bằng chỉ cần lực từ khác không, không cần xét trọng lực hay lực căng.” sai vì trái với kiến thức nêu trên.
}
\end{bt}

% ===== Câu 27 =====
% ID: L12C3B15-07-TLN
% Mức: TH
\begin{ex}
% ID: L12C3B15-07-TLN
% Mức: TH
Trong bốn phát biểu về Cân bằng và chuyển động của dây dẫn mang dòng điện trong từ trường, có bao nhiêu phát biểu đúng? (Chỉ ghi một số nguyên.) (1) Dây cân bằng khi tổng vectơ các lực, gồm lực từ và các lực cơ học, bằng vectơ không. (2) Cần kiểm tra điều kiện áp dụng. (3) Dây cân bằng chỉ cần lực từ khác không, không cần xét trọng lực hay lực căng. (4) Có thể bỏ qua đơn vị của kết quả.
\shortans{2}
\loigiai{
Đối chiếu từng phát biểu với kiến thức: Dây cân bằng khi tổng vectơ các lực, gồm lực từ và các lực cơ học, bằng vectơ không. Số phát biểu đúng là 2.
}
\end{ex}

% ===== Câu 28 =====
% ID: L12C3B15-07-TN
% Mức: TH
\dangbt{Bài toán tổng hợp lực từ gắn với đồ thị, thí nghiệm và thực tiễn}
\begin{ex}
% ID: L12C3B15-07-TN
% Mức: TH
Chọn phát biểu không trái với kiến thức về Bài toán tổng hợp lực từ gắn với đồ thị, thí nghiệm và thực tiễn?
\choice
{Kết quả không phụ thuộc dữ kiện và điều kiện áp dụng.}
{\True Với $B$ và l không đổi, đồ thị $F$ theo $I$ là đường thẳng qua gốc có hệ số góc Bl sin(alpha).}
{Đồ thị $F$ theo $I$ luôn là đường cong và không phụ thuộc B.}
{Có thể bỏ qua phương, chiều và đơn vị của các đại lượng trong mọi trường hợp.}
\loigiai{
Với $B$ và l không đổi, đồ thị $F$ theo $I$ là đường thẳng qua gốc có hệ số góc Bl sin(alpha). Vì vậy phương án $B$ đúng; các phương án còn lại trái với kiến thức hoặc điều kiện áp dụng.
}
\end{ex}

% ===== Câu 29 =====
% ID: L12C3B15-08-DS
% Mức: VD
\dangbt{Cân bằng và chuyển động của dây dẫn mang dòng điện trong từ trường}
\begin{ex}
% ID: L12C3B15-08-DS
% Mức: VD
Xét Cân bằng và chuyển động của dây dẫn mang dòng điện trong từ trường (Tình huống 3). Hãy xác định tính đúng, sai của bốn phát biểu sau.
\choiceTF
{\True Khi giải cần xác định đúng phương, chiều, điều kiện áp dụng và đơn vị.}
{Kết quả không phụ thuộc dữ kiện và có thể bỏ qua điều kiện.}
{\True Dây cân bằng khi tổng vectơ các lực, gồm lực từ và các lực cơ học, bằng vectơ không.}
{Dây cân bằng chỉ cần lực từ khác không, không cần xét trọng lực hay lực căng.}
\loigiai{
\begin{itemchoice}
\item (Đúng) Khi giải cần xác định đúng phương, chiều, điều kiện áp dụng và đơn vị. (Vì): Khi giải cần xác định đúng phương, chiều, điều kiện áp dụng và đơn vị. b) Sai: Kết quả không phụ thuộc dữ kiện và có thể bỏ qua điều kiện. c) Đúng: Dây cân bằng khi tổng vectơ các lực, gồm lực từ và các lực cơ học, bằng vectơ không. d) Sai: Dây cân bằng chỉ cần lực từ khác không, không cần xét trọng lực hay lực căng. Kiến thức cốt lõi: Dây cân bằng khi tổng vectơ các lực, gồm lực từ và các lực cơ học, bằng vectơ không
\item (Sai) Kết quả không phụ thuộc dữ kiện và có thể bỏ qua điều kiện.
\item (Đúng) Dây cân bằng khi tổng vectơ các lực, gồm lực từ và các lực cơ học, bằng vectơ không.
\item (Sai) Dây cân bằng chỉ cần lực từ khác không, không cần xét trọng lực hay lực căng.
\end{itemchoice}
}
\end{ex}

% ===== Câu 30 =====
% ID: L12C3B15-08-TL
% Mức: TH
\begin{bt}
% ID: L12C3B15-08-TL
% Mức: TH
Giải thích vì sao nhận định sau đúng: “Dây cân bằng khi tổng vectơ các lực, gồm lực từ và các lực cơ học, bằng vectơ không.”
\loigiai{
Cần nêu được: Dây cân bằng khi tổng vectơ các lực, gồm lực từ và các lực cơ học, bằng vectơ không. Khi vận dụng phải xác định đúng dữ kiện, phương chiều nếu có, điều kiện áp dụng, hệ đơn vị và kiểm tra tính hợp lí. Nhận định “Dây cân bằng chỉ cần lực từ khác không, không cần xét trọng lực hay lực căng.” sai vì trái với kiến thức nêu trên.
}
\end{bt}

% ===== Câu 31 =====
% ID: L12C3B15-08-TLN
% Mức: TH
\begin{ex}
% ID: L12C3B15-08-TLN
% Mức: TH
Trong bốn phát biểu về Cân bằng và chuyển động của dây dẫn mang dòng điện trong từ trường, có bao nhiêu phát biểu đúng? (Chỉ ghi một số nguyên.) (1) Dây cân bằng chỉ cần lực từ khác không, không cần xét trọng lực hay lực căng. (2) Dây cân bằng khi tổng vectơ các lực, gồm lực từ và các lực cơ học, bằng vectơ không. (3) Dây cân bằng chỉ cần lực từ khác không, không cần xét trọng lực hay lực căng. (4) Có thể bỏ qua đơn vị của kết quả.
\shortans{1}
\loigiai{
Đối chiếu từng phát biểu với kiến thức: Dây cân bằng khi tổng vectơ các lực, gồm lực từ và các lực cơ học, bằng vectơ không. Số phát biểu đúng là 1.
}
\end{ex}

% ===== Câu 32 =====
% ID: L12C3B15-08-TN
% Mức: VD
\dangbt{Bài toán tổng hợp lực từ gắn với đồ thị, thí nghiệm và thực tiễn}
\begin{ex}
% ID: L12C3B15-08-TN
% Mức: VD
Cơ sở Vật lí đúng dùng để giải Bài toán tổng hợp lực từ gắn với đồ thị, thí nghiệm và thực tiễn?
\choice
{Có thể bỏ qua phương, chiều và đơn vị của các đại lượng trong mọi trường hợp.}
{Kết quả không phụ thuộc dữ kiện và điều kiện áp dụng.}
{\True Với $B$ và l không đổi, đồ thị $F$ theo $I$ là đường thẳng qua gốc có hệ số góc Bl sin(alpha).}
{Đồ thị $F$ theo $I$ luôn là đường cong và không phụ thuộc B.}
\loigiai{
Với $B$ và l không đổi, đồ thị $F$ theo $I$ là đường thẳng qua gốc có hệ số góc Bl sin(alpha). Vì vậy phương án $C$ đúng; các phương án còn lại trái với kiến thức hoặc điều kiện áp dụng.
}
\end{ex}

% ===== Câu 33 =====
% ID: L12C3B15-09-DS
% Mức: VD
\dangbt{Cân bằng và chuyển động của dây dẫn mang dòng điện trong từ trường}
\begin{ex}
% ID: L12C3B15-09-DS
% Mức: VD
Xét Cân bằng và chuyển động của dây dẫn mang dòng điện trong từ trường (Tình huống 4). Hãy xác định tính đúng, sai của bốn phát biểu sau.
\choiceTF
{Kết quả không phụ thuộc dữ kiện và có thể bỏ qua điều kiện.}
{\True Khi giải cần xác định đúng phương, chiều, điều kiện áp dụng và đơn vị.}
{Dây cân bằng chỉ cần lực từ khác không, không cần xét trọng lực hay lực căng.}
{\True Dây cân bằng khi tổng vectơ các lực, gồm lực từ và các lực cơ học, bằng vectơ không.}
\loigiai{
\begin{itemchoice}
\item (Sai) Kết quả không phụ thuộc dữ kiện và có thể bỏ qua điều kiện. (Vì): Kết quả không phụ thuộc dữ kiện và có thể bỏ qua điều kiện. b) Đúng: Khi giải cần xác định đúng phương, chiều, điều kiện áp dụng và đơn vị. c) Sai: Dây cân bằng chỉ cần lực từ khác không, không cần xét trọng lực hay lực căng. d) Đúng: Dây cân bằng khi tổng vectơ các lực, gồm lực từ và các lực cơ học, bằng vectơ không. Kiến thức cốt lõi: Dây cân bằng khi tổng vectơ các lực, gồm lực từ và các lực cơ học, bằng vectơ không
\item (Đúng) Khi giải cần xác định đúng phương, chiều, điều kiện áp dụng và đơn vị.
\item (Sai) Dây cân bằng chỉ cần lực từ khác không, không cần xét trọng lực hay lực căng.
\item (Đúng) Dây cân bằng khi tổng vectơ các lực, gồm lực từ và các lực cơ học, bằng vectơ không.
\end{itemchoice}
}
\end{ex}

% ===== Câu 34 =====
% ID: L12C3B15-09-TL
% Mức: VD
\begin{bt}
% ID: L12C3B15-09-TL
% Mức: VD
Phân tích sai lầm trong nhận định: “Dây cân bằng chỉ cần lực từ khác không, không cần xét trọng lực hay lực căng.”
\loigiai{
Cần nêu được: Dây cân bằng khi tổng vectơ các lực, gồm lực từ và các lực cơ học, bằng vectơ không. Khi vận dụng phải xác định đúng dữ kiện, phương chiều nếu có, điều kiện áp dụng, hệ đơn vị và kiểm tra tính hợp lí. Nhận định “Dây cân bằng chỉ cần lực từ khác không, không cần xét trọng lực hay lực căng.” sai vì trái với kiến thức nêu trên.
}
\end{bt}

% ===== Câu 35 =====
% ID: L12C3B15-09-TLN
% Mức: TH
\begin{ex}
% ID: L12C3B15-09-TLN
% Mức: TH
Trong bốn phát biểu về Cân bằng và chuyển động của dây dẫn mang dòng điện trong từ trường, có bao nhiêu phát biểu đúng? (Chỉ ghi một số nguyên.) (1) Dây cân bằng khi tổng vectơ các lực, gồm lực từ và các lực cơ học, bằng vectơ không. (2) Cần kiểm tra điều kiện áp dụng. (3) Dây cân bằng chỉ cần lực từ khác không, không cần xét trọng lực hay lực căng. (4) Có thể bỏ qua đơn vị của kết quả.
\shortans{2}
\loigiai{
Đối chiếu từng phát biểu với kiến thức: Dây cân bằng khi tổng vectơ các lực, gồm lực từ và các lực cơ học, bằng vectơ không. Số phát biểu đúng là 2.
}
\end{ex}

% ===== Câu 36 =====
% ID: L12C3B15-09-TN
% Mức: VD
\dangbt{Bài toán tổng hợp lực từ gắn với đồ thị, thí nghiệm và thực tiễn}
\begin{ex}
% ID: L12C3B15-09-TN
% Mức: VD
Trong một bài thực hành về Bài toán tổng hợp lực từ gắn với đồ thị, thí nghiệm và thực tiễn?
\choice
{Đồ thị $F$ theo $I$ luôn là đường cong và không phụ thuộc B.}
{Có thể bỏ qua phương, chiều và đơn vị của các đại lượng trong mọi trường hợp.}
{Kết quả không phụ thuộc dữ kiện và điều kiện áp dụng.}
{\True Với $B$ và l không đổi, đồ thị $F$ theo $I$ là đường thẳng qua gốc có hệ số góc Bl sin(alpha).}
\loigiai{
Với $B$ và l không đổi, đồ thị $F$ theo $I$ là đường thẳng qua gốc có hệ số góc Bl sin(alpha). Vì vậy phương án $D$ đúng; các phương án còn lại trái với kiến thức hoặc điều kiện áp dụng.
}
\end{ex}

% ===== Câu 37 =====
% ID: L12C3B15-10-DS
% Mức: VD
\dangbt{Cân bằng và chuyển động của dây dẫn mang dòng điện trong từ trường}
\begin{ex}
% ID: L12C3B15-10-DS
% Mức: VD
Xét Cân bằng và chuyển động của dây dẫn mang dòng điện trong từ trường (Tình huống 5). Hãy xác định tính đúng, sai của bốn phát biểu sau.
\choiceTF
{\True Dây cân bằng khi tổng vectơ các lực, gồm lực từ và các lực cơ học, bằng vectơ không.}
{Kết quả không phụ thuộc dữ kiện và có thể bỏ qua điều kiện.}
{Dây cân bằng chỉ cần lực từ khác không, không cần xét trọng lực hay lực căng.}
{\True Khi giải cần xác định đúng phương, chiều, điều kiện áp dụng và đơn vị.}
\loigiai{
\begin{itemchoice}
\item (Đúng) Dây cân bằng khi tổng vectơ các lực, gồm lực từ và các lực cơ học, bằng vectơ không. (Vì): ây cân bằng khi tổng vectơ các lực, gồm lực từ và các lực cơ học, bằng vectơ không. b) Sai: Kết quả không phụ thuộc dữ kiện và có thể bỏ qua điều kiện. c) Sai: Dây cân bằng chỉ cần lực từ khác không, không cần xét trọng lực hay lực căng. d) Đúng: Khi giải cần xác định đúng phương, chiều, điều kiện áp dụng và đơn vị. Kiến thức cốt lõi: Dây cân bằng khi tổng vectơ các lực, gồm lực từ và các lực cơ học, bằng vectơ không
\item (Sai) Kết quả không phụ thuộc dữ kiện và có thể bỏ qua điều kiện.
\item (Sai) Dây cân bằng chỉ cần lực từ khác không, không cần xét trọng lực hay lực căng.
\item (Đúng) Khi giải cần xác định đúng phương, chiều, điều kiện áp dụng và đơn vị.
\end{itemchoice}
}
\end{ex}

% ===== Câu 38 =====
% ID: L12C3B15-10-TL
% Mức: VD
\begin{bt}
% ID: L12C3B15-10-TL
% Mức: VD
Đề xuất các bước giải một bài toán thuộc dạng Cân bằng và chuyển động của dây dẫn mang dòng điện trong từ trường.
\loigiai{
Cần nêu được: Dây cân bằng khi tổng vectơ các lực, gồm lực từ và các lực cơ học, bằng vectơ không. Khi vận dụng phải xác định đúng dữ kiện, phương chiều nếu có, điều kiện áp dụng, hệ đơn vị và kiểm tra tính hợp lí. Nhận định “Dây cân bằng chỉ cần lực từ khác không, không cần xét trọng lực hay lực căng.” sai vì trái với kiến thức nêu trên.
}
\end{bt}

% ===== Câu 39 =====
% ID: L12C3B15-10-TLN
% Mức: VD
\begin{ex}
% ID: L12C3B15-10-TLN
% Mức: VD
Trong bốn phát biểu về Cân bằng và chuyển động của dây dẫn mang dòng điện trong từ trường, có bao nhiêu phát biểu đúng? (Chỉ ghi một số nguyên.) (1) Dây cân bằng chỉ cần lực từ khác không, không cần xét trọng lực hay lực căng. (2) Dây cân bằng khi tổng vectơ các lực, gồm lực từ và các lực cơ học, bằng vectơ không. (3) Dây cân bằng chỉ cần lực từ khác không, không cần xét trọng lực hay lực căng. (4) Có thể bỏ qua đơn vị của kết quả.
\shortans{2}
\loigiai{
Đối chiếu từng phát biểu với kiến thức: Dây cân bằng khi tổng vectơ các lực, gồm lực từ và các lực cơ học, bằng vectơ không. Số phát biểu đúng là 2.
}
\end{ex}

% ===== Câu 40 =====
% ID: L12C3B15-10-TN
% Mức: VD
\dangbt{Bài toán tổng hợp lực từ gắn với đồ thị, thí nghiệm và thực tiễn}
\begin{ex}
% ID: L12C3B15-10-TN
% Mức: VD
Kết luận đúng khi vận dụng Bài toán tổng hợp lực từ gắn với đồ thị, thí nghiệm và thực tiễn?
\choice
{\True Với $B$ và l không đổi, đồ thị $F$ theo $I$ là đường thẳng qua gốc có hệ số góc Bl sin(alpha).}
{Đồ thị $F$ theo $I$ luôn là đường cong và không phụ thuộc B.}
{Có thể bỏ qua phương, chiều và đơn vị của các đại lượng trong mọi trường hợp.}
{Kết quả không phụ thuộc dữ kiện và điều kiện áp dụng.}
\loigiai{
Với $B$ và l không đổi, đồ thị $F$ theo $I$ là đường thẳng qua gốc có hệ số góc Bl sin(alpha). Vì vậy phương án $A$ đúng; các phương án còn lại trái với kiến thức hoặc điều kiện áp dụng.
}
\end{ex}

% ===== Câu 41 =====
% ID: L12C3B15-100-TN
% Mức: VD
\dangbt{Xác định cảm ứng từ từ phép đo lực từ}
\begin{ex}
% ID: L12C3B15-100-TN
% Mức: VD
Trong một bài thực hành về Xác định cảm ứng từ từ phép đo lực từ?
\choice
{Khi dây vuông góc từ trường, cảm ứng từ bằng FIl.}
{Có thể bỏ qua phương, chiều và đơn vị của các đại lượng trong mọi trường hợp.}
{Kết quả không phụ thuộc dữ kiện và điều kiện áp dụng.}
{\True Khi dây vuông góc từ trường, cảm ứng từ được xác định bởi $B$ = F/(Il).}
\loigiai{
Khi dây vuông góc từ trường, cảm ứng từ được xác định bởi $B$ = F/(Il). Vì vậy phương án $D$ đúng; các phương án còn lại trái với kiến thức hoặc điều kiện áp dụng.
}
\end{ex}

% ===== Câu 42 =====
% ID: L12C3B15-101-TN
% Mức: VD
\begin{ex}
% ID: L12C3B15-101-TN
% Mức: VD
Kết luận đúng khi vận dụng Xác định cảm ứng từ từ phép đo lực từ?
\choice
{\True Khi dây vuông góc từ trường, cảm ứng từ được xác định bởi $B$ = F/(Il).}
{Khi dây vuông góc từ trường, cảm ứng từ bằng FIl.}
{Có thể bỏ qua phương, chiều và đơn vị của các đại lượng trong mọi trường hợp.}
{Kết quả không phụ thuộc dữ kiện và điều kiện áp dụng.}
\loigiai{
Khi dây vuông góc từ trường, cảm ứng từ được xác định bởi $B$ = F/(Il). Vì vậy phương án $A$ đúng; các phương án còn lại trái với kiến thức hoặc điều kiện áp dụng.
}
\end{ex}

% ===== Câu 43 =====
% ID: L12C3B15-102-TN
% Mức: NB
\dangbt{Cân bằng và chuyển động của dây dẫn mang dòng điện trong từ trường}
\begin{ex}
% ID: L12C3B15-102-TN
% Mức: NB
Phát biểu nào sau đây đúng về Cân bằng và chuyển động của dây dẫn mang dòng điện trong từ trường?
\choice
{Dây cân bằng chỉ cần lực từ khác không, không cần xét trọng lực hay lực căng.}
{Có thể bỏ qua phương, chiều và đơn vị của các đại lượng trong mọi trường hợp.}
{Kết quả không phụ thuộc dữ kiện và điều kiện áp dụng.}
{\True Dây cân bằng khi tổng vectơ các lực, gồm lực từ và các lực cơ học, bằng vectơ không.}
\loigiai{
Dây cân bằng khi tổng vectơ các lực, gồm lực từ và các lực cơ học, bằng vectơ không. Vì vậy phương án $D$ đúng; các phương án còn lại trái với kiến thức hoặc điều kiện áp dụng.
}
\end{ex}

% ===== Câu 44 =====
% ID: L12C3B15-103-TN
% Mức: NB
\begin{ex}
% ID: L12C3B15-103-TN
% Mức: NB
Trong nội dung Cân bằng và chuyển động của dây dẫn mang dòng điện trong từ trường?
\choice
{\True Dây cân bằng khi tổng vectơ các lực, gồm lực từ và các lực cơ học, bằng vectơ không.}
{Dây cân bằng chỉ cần lực từ khác không, không cần xét trọng lực hay lực căng.}
{Có thể bỏ qua phương, chiều và đơn vị của các đại lượng trong mọi trường hợp.}
{Kết quả không phụ thuộc dữ kiện và điều kiện áp dụng.}
\loigiai{
Dây cân bằng khi tổng vectơ các lực, gồm lực từ và các lực cơ học, bằng vectơ không. Vì vậy phương án $A$ đúng; các phương án còn lại trái với kiến thức hoặc điều kiện áp dụng.
}
\end{ex}

% ===== Câu 45 =====
% ID: L12C3B15-104-TN
% Mức: NB
\begin{ex}
% ID: L12C3B15-104-TN
% Mức: NB
Tình huống 3: chọn kết luận chính xác về Cân bằng và chuyển động của dây dẫn mang dòng điện trong từ trường?
\choice
{Kết quả không phụ thuộc dữ kiện và điều kiện áp dụng.}
{\True Dây cân bằng khi tổng vectơ các lực, gồm lực từ và các lực cơ học, bằng vectơ không.}
{Dây cân bằng chỉ cần lực từ khác không, không cần xét trọng lực hay lực căng.}
{Có thể bỏ qua phương, chiều và đơn vị của các đại lượng trong mọi trường hợp.}
\loigiai{
Dây cân bằng khi tổng vectơ các lực, gồm lực từ và các lực cơ học, bằng vectơ không. Vì vậy phương án $B$ đúng; các phương án còn lại trái với kiến thức hoặc điều kiện áp dụng.
}
\end{ex}

% ===== Câu 46 =====
% ID: L12C3B15-105-TN
% Mức: TH
\begin{ex}
% ID: L12C3B15-105-TN
% Mức: TH
Nội dung nào mô tả đúng nhất Cân bằng và chuyển động của dây dẫn mang dòng điện trong từ trường?
\choice
{Có thể bỏ qua phương, chiều và đơn vị của các đại lượng trong mọi trường hợp.}
{Kết quả không phụ thuộc dữ kiện và điều kiện áp dụng.}
{\True Dây cân bằng khi tổng vectơ các lực, gồm lực từ và các lực cơ học, bằng vectơ không.}
{Dây cân bằng chỉ cần lực từ khác không, không cần xét trọng lực hay lực căng.}
\loigiai{
Dây cân bằng khi tổng vectơ các lực, gồm lực từ và các lực cơ học, bằng vectơ không. Vì vậy phương án $C$ đúng; các phương án còn lại trái với kiến thức hoặc điều kiện áp dụng.
}
\end{ex}

% ===== Câu 47 =====
% ID: L12C3B15-106-TN
% Mức: TH
\begin{ex}
% ID: L12C3B15-106-TN
% Mức: TH
Khi phân tích Cân bằng và chuyển động của dây dẫn mang dòng điện trong từ trường?
\choice
{Dây cân bằng chỉ cần lực từ khác không, không cần xét trọng lực hay lực căng.}
{Có thể bỏ qua phương, chiều và đơn vị của các đại lượng trong mọi trường hợp.}
{Kết quả không phụ thuộc dữ kiện và điều kiện áp dụng.}
{\True Dây cân bằng khi tổng vectơ các lực, gồm lực từ và các lực cơ học, bằng vectơ không.}
\loigiai{
Dây cân bằng khi tổng vectơ các lực, gồm lực từ và các lực cơ học, bằng vectơ không. Vì vậy phương án $D$ đúng; các phương án còn lại trái với kiến thức hoặc điều kiện áp dụng.
}
\end{ex}

% ===== Câu 48 =====
% ID: L12C3B15-107-TN
% Mức: TH
\begin{ex}
% ID: L12C3B15-107-TN
% Mức: TH
Một học sinh ôn tập Cân bằng và chuyển động của dây dẫn mang dòng điện trong từ trường?
\choice
{\True Dây cân bằng khi tổng vectơ các lực, gồm lực từ và các lực cơ học, bằng vectơ không.}
{Dây cân bằng chỉ cần lực từ khác không, không cần xét trọng lực hay lực căng.}
{Có thể bỏ qua phương, chiều và đơn vị của các đại lượng trong mọi trường hợp.}
{Kết quả không phụ thuộc dữ kiện và điều kiện áp dụng.}
\loigiai{
Dây cân bằng khi tổng vectơ các lực, gồm lực từ và các lực cơ học, bằng vectơ không. Vì vậy phương án $A$ đúng; các phương án còn lại trái với kiến thức hoặc điều kiện áp dụng.
}
\end{ex}

% ===== Câu 49 =====
% ID: L12C3B15-108-TN
% Mức: TH
\begin{ex}
% ID: L12C3B15-108-TN
% Mức: TH
Chọn phát biểu không trái với kiến thức về Cân bằng và chuyển động của dây dẫn mang dòng điện trong từ trường?
\choice
{Kết quả không phụ thuộc dữ kiện và điều kiện áp dụng.}
{\True Dây cân bằng khi tổng vectơ các lực, gồm lực từ và các lực cơ học, bằng vectơ không.}
{Dây cân bằng chỉ cần lực từ khác không, không cần xét trọng lực hay lực căng.}
{Có thể bỏ qua phương, chiều và đơn vị của các đại lượng trong mọi trường hợp.}
\loigiai{
Dây cân bằng khi tổng vectơ các lực, gồm lực từ và các lực cơ học, bằng vectơ không. Vì vậy phương án $B$ đúng; các phương án còn lại trái với kiến thức hoặc điều kiện áp dụng.
}
\end{ex}

% ===== Câu 50 =====
% ID: L12C3B15-109-TN
% Mức: VD
\begin{ex}
% ID: L12C3B15-109-TN
% Mức: VD
Cơ sở Vật lí đúng dùng để giải Cân bằng và chuyển động của dây dẫn mang dòng điện trong từ trường?
\choice
{Có thể bỏ qua phương, chiều và đơn vị của các đại lượng trong mọi trường hợp.}
{Kết quả không phụ thuộc dữ kiện và điều kiện áp dụng.}
{\True Dây cân bằng khi tổng vectơ các lực, gồm lực từ và các lực cơ học, bằng vectơ không.}
{Dây cân bằng chỉ cần lực từ khác không, không cần xét trọng lực hay lực căng.}
\loigiai{
Dây cân bằng khi tổng vectơ các lực, gồm lực từ và các lực cơ học, bằng vectơ không. Vì vậy phương án $C$ đúng; các phương án còn lại trái với kiến thức hoặc điều kiện áp dụng.
}
\end{ex}

% ===== Câu 51 =====
% ID: L12C3B15-11-DS
% Mức: VD
\begin{ex}
% ID: L12C3B15-11-DS
% Mức: VD
Một thanh kim loại có khối lượng m= $50\,\text{g}$ có thể trượt với ma sát không đáng kể trên hai thanh ray song song nằm ngang cách nhau một khoảng d= $4\,\text{cm}$. Đường ray nằm trong một từ trường đều thẳng đứng có độ lớn B=0,3 $T$ và có hướng như hình bên. Tại thời điểm t= $0\,\text{s}$, điện kế $G$ được kết nối với thanh ray, tạo ra dòng điện không đổiI=2 $A$ (có chiều như hình) trong dây và thanh ray (kể cả khi dây chuyển động). Biết ban đầu thanh đứng yên
\choiceTF
{\True Lực tác dụng lên thanh là lực từ.}
{Từ trường do dòng điện tạo ra có hướng hướng theo chiều như từ trường bên ngoài.}
{\True Thanh kim loại chuyển động sang trái và tại lúc t= $1\,\text{s}$ vận tốc của thanh có độ lớn là $0,48\,\text{m}$ /s.}
{Quãng đường thanh đi được sau thời gian $2\,\text{s}$ kể từ lúc thiết bị $G$ được kết nối là $0,48\,\text{m}$.}
\loigiai{
\begin{itemchoice}
\item (Đúng) Lực tác dụng lên thanh là lực từ.
\item (Sai) Từ trường do dòng điện tạo ra có hướng hướng theo chiều như từ trường bên ngoài.
\item (Đúng) Thanh kim loại chuyển động sang trái và tại lúc t= $1\,\text{s}$ vận tốc của thanh có độ lớn là $0,48\,\text{m}$ /s. (Vì): ùng chiều với từ trường bên ngoài), nhưng ở phía bên phải của thanh, từ trường này hướng từ ngoài vào trong (ngược chiều với từ trường bên ngoài). Do đó, không thể nói chung rằng từ trường do dòng điện tạo ra có hướng cùng chiều với từ trường bên ngoài.
 
 
• \textbf{Đúng.}
 
 Thanh kim loại chuyển động sang trái và tại lúc $t=1\,\text{s}$, vận tốc của thanh có độ lớn là $0,48\,\text{m/s}$.
 
 Vì theo quy tắc bàn tay trái, lực từ tác dụng lên thanh hướng sang trái.
 
 Gia tốc của thanh là:
  
 a=\dfrac{F_{\text{từ}}}{m}
 =\dfrac{IB\ell}{m}
 =\dfrac{2\cdot 0,3\cdot 4\cdot 10^{-2}}{50\cdot 10^{-3}}
 =0,48 \text{m/s}^2.
  
 
 Tốc độ của thanh tại $t=1\,\text{s}$ là:
  
 v=at=0,48\cdot 1=0,48 \text{m/s}.
  
 
 
• $\textbf{Sai.}$
 
 Quãng đường thanh đi được sau thời gian $2\,\text{s}$ kể từ lúc thiết bị $(G)$ được kết nối là $0,48\,\text{m}$.
 
 Vì chuyển động của thanh là chuyển động nhanh dần đều với gia tốc:
  
 a=\dfrac{F_{\text{từ}}}{m}
 =\dfrac{IB\ell}{m}
 =\dfrac{2\cdot 0,3\cdot 4\cdot 10^{-2}}{50\cdot 10^{-3}}
 =0,48 \text{m/s}^2.
  
 
 Quãng đường thanh đi được sau $2\,\text{s}$ là:
  
 s=\dfrac{1}{2}at^2
 =\dfrac{1}{2}\cdot 0,48\cdot 2^2
 =0,96 \text{m}.
  
 
 Vậy quãng đường đúng là $0,96\,\text{m}$, không phải $0,48\,\text{m}$
\item (Sai) Quãng đường thanh đi được sau thời gian $2\,\text{s}$ kể từ lúc thiết bị $G$ được kết nối là $0,48\,\text{m}$.
\end{itemchoice}
}
\end{ex}

% ===== Câu 52 =====
% ID: L12C3B15-11-TL
% Mức: VD
\begin{bt}
% ID: L12C3B15-11-TL
% Mức: VD
Nêu một tình huống thực tiễn liên quan đến Cân bằng và chuyển động của dây dẫn mang dòng điện trong từ trường và giải thích bằng kiến thức Vật lí.
\loigiai{
Cần nêu được: Dây cân bằng khi tổng vectơ các lực, gồm lực từ và các lực cơ học, bằng vectơ không. Khi vận dụng phải xác định đúng dữ kiện, phương chiều nếu có, điều kiện áp dụng, hệ đơn vị và kiểm tra tính hợp lí. Nhận định “Dây cân bằng chỉ cần lực từ khác không, không cần xét trọng lực hay lực căng.” sai vì trái với kiến thức nêu trên.
}
\end{bt}

% ===== Câu 53 =====
% ID: L12C3B15-11-TLN
% Mức: VD
\begin{ex}
% ID: L12C3B15-11-TLN
% Mức: VD
Trong bốn phát biểu về Cân bằng và chuyển động của dây dẫn mang dòng điện trong từ trường, có bao nhiêu phát biểu đúng? (Chỉ ghi một số nguyên.) (1) Dây cân bằng khi tổng vectơ các lực, gồm lực từ và các lực cơ học, bằng vectơ không. (2) Cần kiểm tra điều kiện áp dụng. (3) Dây cân bằng chỉ cần lực từ khác không, không cần xét trọng lực hay lực căng. (4) Có thể bỏ qua đơn vị của kết quả.
\shortans{2}
\loigiai{
Đối chiếu từng phát biểu với kiến thức: Dây cân bằng khi tổng vectơ các lực, gồm lực từ và các lực cơ học, bằng vectơ không. Số phát biểu đúng là 2.
}
\end{ex}

% ===== Câu 54 =====
% ID: L12C3B15-11-TN
% Mức: NB
\begin{ex}
% ID: L12C3B15-11-TN
% Mức: NB
Phát biểu nào sau đây đúng về Cân bằng và chuyển động của dây dẫn mang dòng điện trong từ trường?
\choice
{Dây cân bằng chỉ cần lực từ khác không, không cần xét trọng lực hay lực căng.}
{Có thể bỏ qua phương, chiều và đơn vị của các đại lượng trong mọi trường hợp.}
{Kết quả không phụ thuộc dữ kiện và điều kiện áp dụng.}
{\True Dây cân bằng khi tổng vectơ các lực, gồm lực từ và các lực cơ học, bằng vectơ không.}
\loigiai{
Dây cân bằng khi tổng vectơ các lực, gồm lực từ và các lực cơ học, bằng vectơ không. Vì vậy phương án $D$ đúng; các phương án còn lại trái với kiến thức hoặc điều kiện áp dụng.
}
\end{ex}

% ===== Câu 55 =====
% ID: L12C3B15-110-TN
% Mức: VD
\begin{ex}
% ID: L12C3B15-110-TN
% Mức: VD
Trong một bài thực hành về Cân bằng và chuyển động của dây dẫn mang dòng điện trong từ trường?
\choice
{Dây cân bằng chỉ cần lực từ khác không, không cần xét trọng lực hay lực căng.}
{Có thể bỏ qua phương, chiều và đơn vị của các đại lượng trong mọi trường hợp.}
{Kết quả không phụ thuộc dữ kiện và điều kiện áp dụng.}
{\True Dây cân bằng khi tổng vectơ các lực, gồm lực từ và các lực cơ học, bằng vectơ không.}
\loigiai{
Dây cân bằng khi tổng vectơ các lực, gồm lực từ và các lực cơ học, bằng vectơ không. Vì vậy phương án $D$ đúng; các phương án còn lại trái với kiến thức hoặc điều kiện áp dụng.
}
\end{ex}

% ===== Câu 56 =====
% ID: L12C3B15-111-TN
% Mức: VD
\begin{ex}
% ID: L12C3B15-111-TN
% Mức: VD
Kết luận đúng khi vận dụng Cân bằng và chuyển động của dây dẫn mang dòng điện trong từ trường?
\choice
{\True Dây cân bằng khi tổng vectơ các lực, gồm lực từ và các lực cơ học, bằng vectơ không.}
{Dây cân bằng chỉ cần lực từ khác không, không cần xét trọng lực hay lực căng.}
{Có thể bỏ qua phương, chiều và đơn vị của các đại lượng trong mọi trường hợp.}
{Kết quả không phụ thuộc dữ kiện và điều kiện áp dụng.}
\loigiai{
Dây cân bằng khi tổng vectơ các lực, gồm lực từ và các lực cơ học, bằng vectơ không. Vì vậy phương án $A$ đúng; các phương án còn lại trái với kiến thức hoặc điều kiện áp dụng.
}
\end{ex}

% ===== Câu 57 =====
% ID: L12C3B15-112-TN
% Mức: NB
\dangbt{Bài toán tổng hợp lực từ gắn với đồ thị, thí nghiệm và thực tiễn}
\begin{ex}
% ID: L12C3B15-112-TN
% Mức: NB
Phát biểu nào sau đây đúng về Bài toán tổng hợp lực từ gắn với đồ thị, thí nghiệm và thực tiễn?
\choice
{Đồ thị $F$ theo $I$ luôn là đường cong và không phụ thuộc B.}
{Có thể bỏ qua phương, chiều và đơn vị của các đại lượng trong mọi trường hợp.}
{Kết quả không phụ thuộc dữ kiện và điều kiện áp dụng.}
{\True Với $B$ và l không đổi, đồ thị $F$ theo $I$ là đường thẳng qua gốc có hệ số góc Bl sin(alpha).}
\loigiai{
Với $B$ và l không đổi, đồ thị $F$ theo $I$ là đường thẳng qua gốc có hệ số góc Bl sin(alpha). Vì vậy phương án $D$ đúng; các phương án còn lại trái với kiến thức hoặc điều kiện áp dụng.
}
\end{ex}

% ===== Câu 58 =====
% ID: L12C3B15-113-TN
% Mức: NB
\begin{ex}
% ID: L12C3B15-113-TN
% Mức: NB
Trong nội dung Bài toán tổng hợp lực từ gắn với đồ thị, thí nghiệm và thực tiễn?
\choice
{\True Với $B$ và l không đổi, đồ thị $F$ theo $I$ là đường thẳng qua gốc có hệ số góc Bl sin(alpha).}
{Đồ thị $F$ theo $I$ luôn là đường cong và không phụ thuộc B.}
{Có thể bỏ qua phương, chiều và đơn vị của các đại lượng trong mọi trường hợp.}
{Kết quả không phụ thuộc dữ kiện và điều kiện áp dụng.}
\loigiai{
Với $B$ và l không đổi, đồ thị $F$ theo $I$ là đường thẳng qua gốc có hệ số góc Bl sin(alpha). Vì vậy phương án $A$ đúng; các phương án còn lại trái với kiến thức hoặc điều kiện áp dụng.
}
\end{ex}

% ===== Câu 59 =====
% ID: L12C3B15-114-TN
% Mức: NB
\begin{ex}
% ID: L12C3B15-114-TN
% Mức: NB
Tình huống 3: chọn kết luận chính xác về Bài toán tổng hợp lực từ gắn với đồ thị, thí nghiệm và thực tiễn?
\choice
{Kết quả không phụ thuộc dữ kiện và điều kiện áp dụng.}
{\True Với $B$ và l không đổi, đồ thị $F$ theo $I$ là đường thẳng qua gốc có hệ số góc Bl sin(alpha).}
{Đồ thị $F$ theo $I$ luôn là đường cong và không phụ thuộc B.}
{Có thể bỏ qua phương, chiều và đơn vị của các đại lượng trong mọi trường hợp.}
\loigiai{
Với $B$ và l không đổi, đồ thị $F$ theo $I$ là đường thẳng qua gốc có hệ số góc Bl sin(alpha). Vì vậy phương án $B$ đúng; các phương án còn lại trái với kiến thức hoặc điều kiện áp dụng.
}
\end{ex}

% ===== Câu 60 =====
% ID: L12C3B15-115-TN
% Mức: TH
\begin{ex}
% ID: L12C3B15-115-TN
% Mức: TH
Nội dung nào mô tả đúng nhất Bài toán tổng hợp lực từ gắn với đồ thị, thí nghiệm và thực tiễn?
\choice
{Có thể bỏ qua phương, chiều và đơn vị của các đại lượng trong mọi trường hợp.}
{Kết quả không phụ thuộc dữ kiện và điều kiện áp dụng.}
{\True Với $B$ và l không đổi, đồ thị $F$ theo $I$ là đường thẳng qua gốc có hệ số góc Bl sin(alpha).}
{Đồ thị $F$ theo $I$ luôn là đường cong và không phụ thuộc B.}
\loigiai{
Với $B$ và l không đổi, đồ thị $F$ theo $I$ là đường thẳng qua gốc có hệ số góc Bl sin(alpha). Vì vậy phương án $C$ đúng; các phương án còn lại trái với kiến thức hoặc điều kiện áp dụng.
}
\end{ex}

% ===== Câu 61 =====
% ID: L12C3B15-116-TN
% Mức: TH
\begin{ex}
% ID: L12C3B15-116-TN
% Mức: TH
Khi phân tích Bài toán tổng hợp lực từ gắn với đồ thị, thí nghiệm và thực tiễn?
\choice
{Đồ thị $F$ theo $I$ luôn là đường cong và không phụ thuộc B.}
{Có thể bỏ qua phương, chiều và đơn vị của các đại lượng trong mọi trường hợp.}
{Kết quả không phụ thuộc dữ kiện và điều kiện áp dụng.}
{\True Với $B$ và l không đổi, đồ thị $F$ theo $I$ là đường thẳng qua gốc có hệ số góc Bl sin(alpha).}
\loigiai{
Với $B$ và l không đổi, đồ thị $F$ theo $I$ là đường thẳng qua gốc có hệ số góc Bl sin(alpha). Vì vậy phương án $D$ đúng; các phương án còn lại trái với kiến thức hoặc điều kiện áp dụng.
}
\end{ex}

% ===== Câu 62 =====
% ID: L12C3B15-117-TN
% Mức: TH
\begin{ex}
% ID: L12C3B15-117-TN
% Mức: TH
Một học sinh ôn tập Bài toán tổng hợp lực từ gắn với đồ thị, thí nghiệm và thực tiễn?
\choice
{\True Với $B$ và l không đổi, đồ thị $F$ theo $I$ là đường thẳng qua gốc có hệ số góc Bl sin(alpha).}
{Đồ thị $F$ theo $I$ luôn là đường cong và không phụ thuộc B.}
{Có thể bỏ qua phương, chiều và đơn vị của các đại lượng trong mọi trường hợp.}
{Kết quả không phụ thuộc dữ kiện và điều kiện áp dụng.}
\loigiai{
Với $B$ và l không đổi, đồ thị $F$ theo $I$ là đường thẳng qua gốc có hệ số góc Bl sin(alpha). Vì vậy phương án $A$ đúng; các phương án còn lại trái với kiến thức hoặc điều kiện áp dụng.
}
\end{ex}

% ===== Câu 63 =====
% ID: L12C3B15-118-TN
% Mức: TH
\begin{ex}
% ID: L12C3B15-118-TN
% Mức: TH
Chọn phát biểu không trái với kiến thức về Bài toán tổng hợp lực từ gắn với đồ thị, thí nghiệm và thực tiễn?
\choice
{Kết quả không phụ thuộc dữ kiện và điều kiện áp dụng.}
{\True Với $B$ và l không đổi, đồ thị $F$ theo $I$ là đường thẳng qua gốc có hệ số góc Bl sin(alpha).}
{Đồ thị $F$ theo $I$ luôn là đường cong và không phụ thuộc B.}
{Có thể bỏ qua phương, chiều và đơn vị của các đại lượng trong mọi trường hợp.}
\loigiai{
Với $B$ và l không đổi, đồ thị $F$ theo $I$ là đường thẳng qua gốc có hệ số góc Bl sin(alpha). Vì vậy phương án $B$ đúng; các phương án còn lại trái với kiến thức hoặc điều kiện áp dụng.
}
\end{ex}

% ===== Câu 64 =====
% ID: L12C3B15-119-TN
% Mức: VD
\begin{ex}
% ID: L12C3B15-119-TN
% Mức: VD
Cơ sở Vật lí đúng dùng để giải Bài toán tổng hợp lực từ gắn với đồ thị, thí nghiệm và thực tiễn?
\choice
{Có thể bỏ qua phương, chiều và đơn vị của các đại lượng trong mọi trường hợp.}
{Kết quả không phụ thuộc dữ kiện và điều kiện áp dụng.}
{\True Với $B$ và l không đổi, đồ thị $F$ theo $I$ là đường thẳng qua gốc có hệ số góc Bl sin(alpha).}
{Đồ thị $F$ theo $I$ luôn là đường cong và không phụ thuộc B.}
\loigiai{
Với $B$ và l không đổi, đồ thị $F$ theo $I$ là đường thẳng qua gốc có hệ số góc Bl sin(alpha). Vì vậy phương án $C$ đúng; các phương án còn lại trái với kiến thức hoặc điều kiện áp dụng.
}
\end{ex}

% ===== Câu 65 =====
% ID: L12C3B15-12-DS
% Mức: TH
\dangbt{Nhận biết lực từ và đặc trưng của vectơ cảm ứng từ}
\begin{ex}
% ID: L12C3B15-12-DS
% Mức: TH
Xét Nhận biết lực từ và đặc trưng của vectơ cảm ứng từ (Tình huống 1). Hãy xác định tính đúng, sai của bốn phát biểu sau.
\choiceTF
{\True Vectơ cảm ứng từ tại một điểm có hướng trùng với hướng của từ trường tại điểm đó; đơn vị là tesla (T).}
{Cảm ứng từ là đại lượng vô hướng và có đơn vị niutơn.}
{\True Khi giải cần xác định đúng phương, chiều, điều kiện áp dụng và đơn vị.}
{Kết quả không phụ thuộc dữ kiện và có thể bỏ qua điều kiện.}
\loigiai{
\begin{itemchoice}
\item (Đúng) Vectơ cảm ứng từ tại một điểm có hướng trùng với hướng của từ trường tại điểm đó; đơn vị là tesla (T). (Vì): Vectơ cảm ứng từ tại một điểm có hướng trùng với hướng của từ trường tại điểm đó; đơn vị là tesla (T). b) Sai: Cảm ứng từ là đại lượng vô hướng và có đơn vị niutơn. c) Đúng: Khi giải cần xác định đúng phương, chiều, điều kiện áp dụng và đơn vị. d) Sai: Kết quả không phụ thuộc dữ kiện và có thể bỏ qua điều kiện. Kiến thức cốt lõi: Vectơ cảm ứng từ tại một điểm có hướng trùng với hướng của từ trường tại điểm đó; đơn vị là tesla (T
\item (Sai) Cảm ứng từ là đại lượng vô hướng và có đơn vị niutơn.
\item (Đúng) Khi giải cần xác định đúng phương, chiều, điều kiện áp dụng và đơn vị.
\item (Sai) Kết quả không phụ thuộc dữ kiện và có thể bỏ qua điều kiện.
\end{itemchoice}
}
\end{ex}

% ===== Câu 66 =====
% ID: L12C3B15-12-TL
% Mức: VDC
\dangbt{Cân bằng và chuyển động của dây dẫn mang dòng điện trong từ trường}
\begin{bt}
% ID: L12C3B15-12-TL
% Mức: VDC
So sánh nhận định đúng “Dây cân bằng khi tổng vectơ các lực, gồm lực từ và các lực cơ học, bằng vectơ không.” với nhận định sai “Dây cân bằng chỉ cần lực từ khác không, không cần xét trọng lực hay lực căng.”; chỉ rõ căn cứ Vật lí.
\loigiai{
Cần nêu được: Dây cân bằng khi tổng vectơ các lực, gồm lực từ và các lực cơ học, bằng vectơ không. Khi vận dụng phải xác định đúng dữ kiện, phương chiều nếu có, điều kiện áp dụng, hệ đơn vị và kiểm tra tính hợp lí. Nhận định “Dây cân bằng chỉ cần lực từ khác không, không cần xét trọng lực hay lực căng.” sai vì trái với kiến thức nêu trên.
}
\end{bt}

% ===== Câu 67 =====
% ID: L12C3B15-12-TLN
% Mức: VD
\begin{ex}
% ID: L12C3B15-12-TLN
% Mức: VD
Một dây dẫn thẳng, cứng, dài $\ell=0,10\mathrm{~m}$, có khối lượng $m=0,025\mathrm{~kg}$ được giữ nằm yên theo phương ngang trong một từ trường có độ lớn cảm ứng từ là $B=0,5\mathrm{~T}$ và có hướng nằm ngang, vuông góc với dây dẫn. Cường độ dòng điện chạy trong dây là bao nhiêu ampe để khi dây được thả ra thì nó vẫn nằm yên (làm tròn kết quả đến chữ số hàng phần mười)?
\shortans{}
\loigiai{
Để dây dẫn nằm yên khi được thả ra, lực từ tác dụng lên dây phải cân bằng với trọng lực của dây. Lực từ tác dụng lên dây dẫn có dòng điện chạy qua trong từ trường được tính bằng công thức: $F = B \cdot I \cdot \ell \cdot \sin\theta$, với $\theta$ là góc giữa dây dẫn và đường sức từ. Trong trường hợp này, $\theta = 90^\circ$ nên $\sin\theta = 1$. Trọng lực tác dụng lên dây là $P = m \cdot g$, với $g = 9,8 \, \mathrm{m/s^2}$. Để dây nằm yên, ta có $F = P$, do đó $B \cdot I \cdot \ell = m \cdot g$. Thay số vào công thức: $0,5 \cdot I \cdot 0,10 = 0,025 \cdot 9,8$. Giải phương trình này để tìm $I$.
 
  Lực từ $F = B \cdot I \cdot \ell = m \cdot g$ với $B = 0,5 \, \mathrm{T}$, $\ell = 0,10 \, \mathrm{m}$, $m = 0,025 \, \mathrm{kg}$, $g = 9,8 \, \mathrm{m/s^2}$. 
 
 Thay vào: $0,5 \cdot I \cdot 0,10 = 0,025 \cdot 9,8$. Giải ra $I = \frac{0,025 \cdot 9,8}{0,5 \cdot 0,10} = 4,9 \, \mathrm{A}$.
}
\end{ex}

% ===== Câu 68 =====
% ID: L12C3B15-12-TN
% Mức: NB
\begin{ex}
% ID: L12C3B15-12-TN
% Mức: NB
Trong nội dung Cân bằng và chuyển động của dây dẫn mang dòng điện trong từ trường?
\choice
{\True Dây cân bằng khi tổng vectơ các lực, gồm lực từ và các lực cơ học, bằng vectơ không.}
{Dây cân bằng chỉ cần lực từ khác không, không cần xét trọng lực hay lực căng.}
{Có thể bỏ qua phương, chiều và đơn vị của các đại lượng trong mọi trường hợp.}
{Kết quả không phụ thuộc dữ kiện và điều kiện áp dụng.}
\loigiai{
Dây cân bằng khi tổng vectơ các lực, gồm lực từ và các lực cơ học, bằng vectơ không. Vì vậy phương án $A$ đúng; các phương án còn lại trái với kiến thức hoặc điều kiện áp dụng.
}
\end{ex}

% ===== Câu 69 =====
% ID: L12C3B15-120-TN
% Mức: VD
\dangbt{Bài toán tổng hợp lực từ gắn với đồ thị, thí nghiệm và thực tiễn}
\begin{ex}
% ID: L12C3B15-120-TN
% Mức: VD
Trong một bài thực hành về Bài toán tổng hợp lực từ gắn với đồ thị, thí nghiệm và thực tiễn?
\choice
{Đồ thị $F$ theo $I$ luôn là đường cong và không phụ thuộc B.}
{Có thể bỏ qua phương, chiều và đơn vị của các đại lượng trong mọi trường hợp.}
{Kết quả không phụ thuộc dữ kiện và điều kiện áp dụng.}
{\True Với $B$ và l không đổi, đồ thị $F$ theo $I$ là đường thẳng qua gốc có hệ số góc Bl sin(alpha).}
\loigiai{
Với $B$ và l không đổi, đồ thị $F$ theo $I$ là đường thẳng qua gốc có hệ số góc Bl sin(alpha). Vì vậy phương án $D$ đúng; các phương án còn lại trái với kiến thức hoặc điều kiện áp dụng.
}
\end{ex}

% ===== Câu 70 =====
% ID: L12C3B15-121-TN
% Mức: VD
\begin{ex}
% ID: L12C3B15-121-TN
% Mức: VD
Kết luận đúng khi vận dụng Bài toán tổng hợp lực từ gắn với đồ thị, thí nghiệm và thực tiễn?
\choice
{\True Với $B$ và l không đổi, đồ thị $F$ theo $I$ là đường thẳng qua gốc có hệ số góc Bl sin(alpha).}
{Đồ thị $F$ theo $I$ luôn là đường cong và không phụ thuộc B.}
{Có thể bỏ qua phương, chiều và đơn vị của các đại lượng trong mọi trường hợp.}
{Kết quả không phụ thuộc dữ kiện và điều kiện áp dụng.}
\loigiai{
Với $B$ và l không đổi, đồ thị $F$ theo $I$ là đường thẳng qua gốc có hệ số góc Bl sin(alpha). Vì vậy phương án $A$ đúng; các phương án còn lại trái với kiến thức hoặc điều kiện áp dụng.
}
\end{ex}

% ===== Câu 71 =====
% ID: L12C3B15-13-DS
% Mức: TH
\dangbt{Nhận biết lực từ và đặc trưng của vectơ cảm ứng từ}
\begin{ex}
% ID: L12C3B15-13-DS
% Mức: TH
Xét Nhận biết lực từ và đặc trưng của vectơ cảm ứng từ (Tình huống 2). Hãy xác định tính đúng, sai của bốn phát biểu sau.
\choiceTF
{Cảm ứng từ là đại lượng vô hướng và có đơn vị niutơn.}
{\True Vectơ cảm ứng từ tại một điểm có hướng trùng với hướng của từ trường tại điểm đó; đơn vị là tesla (T).}
{Kết quả không phụ thuộc dữ kiện và có thể bỏ qua điều kiện.}
{\True Khi giải cần xác định đúng phương, chiều, điều kiện áp dụng và đơn vị.}
\loigiai{
\begin{itemchoice}
\item (Sai) Cảm ứng từ là đại lượng vô hướng và có đơn vị niutơn. (Vì): Cảm ứng từ là đại lượng vô hướng và có đơn vị niutơn. b) Đúng: Vectơ cảm ứng từ tại một điểm có hướng trùng với hướng của từ trường tại điểm đó; đơn vị là tesla (T). c) Sai: Kết quả không phụ thuộc dữ kiện và có thể bỏ qua điều kiện. d) Đúng: Khi giải cần xác định đúng phương, chiều, điều kiện áp dụng và đơn vị. Kiến thức cốt lõi: Vectơ cảm ứng từ tại một điểm có hướng trùng với hướng của từ trường tại điểm đó; đơn vị là tesla (T
\item (Đúng) Vectơ cảm ứng từ tại một điểm có hướng trùng với hướng của từ trường tại điểm đó; đơn vị là tesla (T).
\item (Sai) Kết quả không phụ thuộc dữ kiện và có thể bỏ qua điều kiện.
\item (Đúng) Khi giải cần xác định đúng phương, chiều, điều kiện áp dụng và đơn vị.
\end{itemchoice}
}
\end{ex}

% ===== Câu 72 =====
% ID: L12C3B15-13-TL
% Mức: TH
\begin{bt}
% ID: L12C3B15-13-TL
% Mức: TH
Trình bày kiến thức cốt lõi của Nhận biết lực từ và đặc trưng của vectơ cảm ứng từ và nêu điều kiện áp dụng.
\loigiai{
Cần nêu được: Vectơ cảm ứng từ tại một điểm có hướng trùng với hướng của từ trường tại điểm đó; đơn vị là tesla (T). Khi vận dụng phải xác định đúng dữ kiện, phương chiều nếu có, điều kiện áp dụng, hệ đơn vị và kiểm tra tính hợp lí. Nhận định “Cảm ứng từ là đại lượng vô hướng và có đơn vị niutơn.” sai vì trái với kiến thức nêu trên.
}
\end{bt}

% ===== Câu 73 =====
% ID: L12C3B15-13-TLN
% Mức: VDC
\dangbt{Cân bằng và chuyển động của dây dẫn mang dòng điện trong từ trường}
\begin{ex}
% ID: L12C3B15-13-TLN
% Mức: VDC
Trong bốn phát biểu về Cân bằng và chuyển động của dây dẫn mang dòng điện trong từ trường, có bao nhiêu phát biểu đúng? (Chỉ ghi một số nguyên.) (1) Dây cân bằng chỉ cần lực từ khác không, không cần xét trọng lực hay lực căng. (2) Dây cân bằng khi tổng vectơ các lực, gồm lực từ và các lực cơ học, bằng vectơ không. (3) Phải dùng đúng định luật cho tình huống. (4) Dây cân bằng khi tổng vectơ các lực, gồm lực từ và các lực cơ học, bằng vectơ không.
\shortans{3}
\loigiai{
Đối chiếu từng phát biểu với kiến thức: Dây cân bằng khi tổng vectơ các lực, gồm lực từ và các lực cơ học, bằng vectơ không. Số phát biểu đúng là 3.
}
\end{ex}

% ===== Câu 74 =====
% ID: L12C3B15-13-TN
% Mức: NB
\begin{ex}
% ID: L12C3B15-13-TN
% Mức: NB
Tình huống 3: chọn kết luận chính xác về Cân bằng và chuyển động của dây dẫn mang dòng điện trong từ trường?
\choice
{Kết quả không phụ thuộc dữ kiện và điều kiện áp dụng.}
{\True Dây cân bằng khi tổng vectơ các lực, gồm lực từ và các lực cơ học, bằng vectơ không.}
{Dây cân bằng chỉ cần lực từ khác không, không cần xét trọng lực hay lực căng.}
{Có thể bỏ qua phương, chiều và đơn vị của các đại lượng trong mọi trường hợp.}
\loigiai{
Dây cân bằng khi tổng vectơ các lực, gồm lực từ và các lực cơ học, bằng vectơ không. Vì vậy phương án $B$ đúng; các phương án còn lại trái với kiến thức hoặc điều kiện áp dụng.
}
\end{ex}

% ===== Câu 75 =====
% ID: L12C3B15-14-DS
% Mức: VD
\dangbt{Nhận biết lực từ và đặc trưng của vectơ cảm ứng từ}
\begin{ex}
% ID: L12C3B15-14-DS
% Mức: VD
Xét Nhận biết lực từ và đặc trưng của vectơ cảm ứng từ (Tình huống 3). Hãy xác định tính đúng, sai của bốn phát biểu sau.
\choiceTF
{\True Khi giải cần xác định đúng phương, chiều, điều kiện áp dụng và đơn vị.}
{Kết quả không phụ thuộc dữ kiện và có thể bỏ qua điều kiện.}
{\True Vectơ cảm ứng từ tại một điểm có hướng trùng với hướng của từ trường tại điểm đó; đơn vị là tesla (T).}
{Cảm ứng từ là đại lượng vô hướng và có đơn vị niutơn.}
\loigiai{
\begin{itemchoice}
\item (Đúng) Khi giải cần xác định đúng phương, chiều, điều kiện áp dụng và đơn vị. (Vì): Khi giải cần xác định đúng phương, chiều, điều kiện áp dụng và đơn vị. b) Sai: Kết quả không phụ thuộc dữ kiện và có thể bỏ qua điều kiện. c) Đúng: Vectơ cảm ứng từ tại một điểm có hướng trùng với hướng của từ trường tại điểm đó; đơn vị là tesla (T). d) Sai: Cảm ứng từ là đại lượng vô hướng và có đơn vị niutơn. Kiến thức cốt lõi: Vectơ cảm ứng từ tại một điểm có hướng trùng với hướng của từ trường tại điểm đó; đơn vị là tesla (T
\item (Sai) Kết quả không phụ thuộc dữ kiện và có thể bỏ qua điều kiện.
\item (Đúng) Vectơ cảm ứng từ tại một điểm có hướng trùng với hướng của từ trường tại điểm đó; đơn vị là tesla (T).
\item (Sai) Cảm ứng từ là đại lượng vô hướng và có đơn vị niutơn.
\end{itemchoice}
}
\end{ex}

% ===== Câu 76 =====
% ID: L12C3B15-14-TL
% Mức: TH
\begin{bt}
% ID: L12C3B15-14-TL
% Mức: TH
Giải thích vì sao nhận định sau đúng: “Vectơ cảm ứng từ tại một điểm có hướng trùng với hướng của từ trường tại điểm đó; đơn vị là tesla (T).”
\loigiai{
Cần nêu được: Vectơ cảm ứng từ tại một điểm có hướng trùng với hướng của từ trường tại điểm đó; đơn vị là tesla (T). Khi vận dụng phải xác định đúng dữ kiện, phương chiều nếu có, điều kiện áp dụng, hệ đơn vị và kiểm tra tính hợp lí. Nhận định “Cảm ứng từ là đại lượng vô hướng và có đơn vị niutơn.” sai vì trái với kiến thức nêu trên.
}
\end{bt}

% ===== Câu 77 =====
% ID: L12C3B15-14-TLN
% Mức: TH
\begin{ex}
% ID: L12C3B15-14-TLN
% Mức: TH
Trong bốn phát biểu về Nhận biết lực từ và đặc trưng của vectơ cảm ứng từ, có bao nhiêu phát biểu đúng? (Chỉ ghi một số nguyên.) (1) Vectơ cảm ứng từ tại một điểm có hướng trùng với hướng của từ trường tại điểm đó; đơn vị là tesla (T). (2) Cần kiểm tra điều kiện áp dụng. (3) Cảm ứng từ là đại lượng vô hướng và có đơn vị niutơn. (4) Có thể bỏ qua đơn vị của kết quả.
\shortans{2}
\loigiai{
Đối chiếu từng phát biểu với kiến thức: Vectơ cảm ứng từ tại một điểm có hướng trùng với hướng của từ trường tại điểm đó; đơn vị là tesla (T). Số phát biểu đúng là 2.
}
\end{ex}

% ===== Câu 78 =====
% ID: L12C3B15-14-TN
% Mức: TH
\dangbt{Cân bằng và chuyển động của dây dẫn mang dòng điện trong từ trường}
\begin{ex}
% ID: L12C3B15-14-TN
% Mức: TH
Nội dung nào mô tả đúng nhất Cân bằng và chuyển động của dây dẫn mang dòng điện trong từ trường?
\choice
{Có thể bỏ qua phương, chiều và đơn vị của các đại lượng trong mọi trường hợp.}
{Kết quả không phụ thuộc dữ kiện và điều kiện áp dụng.}
{\True Dây cân bằng khi tổng vectơ các lực, gồm lực từ và các lực cơ học, bằng vectơ không.}
{Dây cân bằng chỉ cần lực từ khác không, không cần xét trọng lực hay lực căng.}
\loigiai{
Dây cân bằng khi tổng vectơ các lực, gồm lực từ và các lực cơ học, bằng vectơ không. Vì vậy phương án $C$ đúng; các phương án còn lại trái với kiến thức hoặc điều kiện áp dụng.
}
\end{ex}

% ===== Câu 79 =====
% ID: L12C3B15-15-DS
% Mức: VD
\dangbt{Nhận biết lực từ và đặc trưng của vectơ cảm ứng từ}
\begin{ex}
% ID: L12C3B15-15-DS
% Mức: VD
Xét Nhận biết lực từ và đặc trưng của vectơ cảm ứng từ (Tình huống 4). Hãy xác định tính đúng, sai của bốn phát biểu sau.
\choiceTF
{Kết quả không phụ thuộc dữ kiện và có thể bỏ qua điều kiện.}
{\True Khi giải cần xác định đúng phương, chiều, điều kiện áp dụng và đơn vị.}
{Cảm ứng từ là đại lượng vô hướng và có đơn vị niutơn.}
{\True Vectơ cảm ứng từ tại một điểm có hướng trùng với hướng của từ trường tại điểm đó; đơn vị là tesla (T).}
\loigiai{
\begin{itemchoice}
\item (Sai) Kết quả không phụ thuộc dữ kiện và có thể bỏ qua điều kiện. (Vì): Kết quả không phụ thuộc dữ kiện và có thể bỏ qua điều kiện. b) Đúng: Khi giải cần xác định đúng phương, chiều, điều kiện áp dụng và đơn vị. c) Sai: Cảm ứng từ là đại lượng vô hướng và có đơn vị niutơn. d) Đúng: Vectơ cảm ứng từ tại một điểm có hướng trùng với hướng của từ trường tại điểm đó; đơn vị là tesla (T). Kiến thức cốt lõi: Vectơ cảm ứng từ tại một điểm có hướng trùng với hướng của từ trường tại điểm đó; đơn vị là tesla (T
\item (Đúng) Khi giải cần xác định đúng phương, chiều, điều kiện áp dụng và đơn vị.
\item (Sai) Cảm ứng từ là đại lượng vô hướng và có đơn vị niutơn.
\item (Đúng) Vectơ cảm ứng từ tại một điểm có hướng trùng với hướng của từ trường tại điểm đó; đơn vị là tesla (T).
\end{itemchoice}
}
\end{ex}

% ===== Câu 80 =====
% ID: L12C3B15-15-TL
% Mức: VD
\begin{bt}
% ID: L12C3B15-15-TL
% Mức: VD
Phân tích sai lầm trong nhận định: “Cảm ứng từ là đại lượng vô hướng và có đơn vị niutơn.”
\loigiai{
Cần nêu được: Vectơ cảm ứng từ tại một điểm có hướng trùng với hướng của từ trường tại điểm đó; đơn vị là tesla (T). Khi vận dụng phải xác định đúng dữ kiện, phương chiều nếu có, điều kiện áp dụng, hệ đơn vị và kiểm tra tính hợp lí. Nhận định “Cảm ứng từ là đại lượng vô hướng và có đơn vị niutơn.” sai vì trái với kiến thức nêu trên.
}
\end{bt}

% ===== Câu 81 =====
% ID: L12C3B15-15-TLN
% Mức: TH
\begin{ex}
% ID: L12C3B15-15-TLN
% Mức: TH
Trong bốn phát biểu về Nhận biết lực từ và đặc trưng của vectơ cảm ứng từ, có bao nhiêu phát biểu đúng? (Chỉ ghi một số nguyên.) (1) Cảm ứng từ là đại lượng vô hướng và có đơn vị niutơn. (2) Vectơ cảm ứng từ tại một điểm có hướng trùng với hướng của từ trường tại điểm đó; đơn vị là tesla (T). (3) Cảm ứng từ là đại lượng vô hướng và có đơn vị niutơn. (4) Có thể bỏ qua đơn vị của kết quả.
\shortans{1}
\loigiai{
Đối chiếu từng phát biểu với kiến thức: Vectơ cảm ứng từ tại một điểm có hướng trùng với hướng của từ trường tại điểm đó; đơn vị là tesla (T). Số phát biểu đúng là 1.
}
\end{ex}

% ===== Câu 82 =====
% ID: L12C3B15-15-TN
% Mức: TH
\dangbt{Cân bằng và chuyển động của dây dẫn mang dòng điện trong từ trường}
\begin{ex}
% ID: L12C3B15-15-TN
% Mức: TH
Khi phân tích Cân bằng và chuyển động của dây dẫn mang dòng điện trong từ trường?
\choice
{Dây cân bằng chỉ cần lực từ khác không, không cần xét trọng lực hay lực căng.}
{Có thể bỏ qua phương, chiều và đơn vị của các đại lượng trong mọi trường hợp.}
{Kết quả không phụ thuộc dữ kiện và điều kiện áp dụng.}
{\True Dây cân bằng khi tổng vectơ các lực, gồm lực từ và các lực cơ học, bằng vectơ không.}
\loigiai{
Dây cân bằng khi tổng vectơ các lực, gồm lực từ và các lực cơ học, bằng vectơ không. Vì vậy phương án $D$ đúng; các phương án còn lại trái với kiến thức hoặc điều kiện áp dụng.
}
\end{ex}

% ===== Câu 83 =====
% ID: L12C3B15-16-DS
% Mức: VD
\dangbt{Nhận biết lực từ và đặc trưng của vectơ cảm ứng từ}
\begin{ex}
% ID: L12C3B15-16-DS
% Mức: VD
Xét Nhận biết lực từ và đặc trưng của vectơ cảm ứng từ (Tình huống 5). Hãy xác định tính đúng, sai của bốn phát biểu sau.
\choiceTF
{\True Vectơ cảm ứng từ tại một điểm có hướng trùng với hướng của từ trường tại điểm đó; đơn vị là tesla (T).}
{Kết quả không phụ thuộc dữ kiện và có thể bỏ qua điều kiện.}
{Cảm ứng từ là đại lượng vô hướng và có đơn vị niutơn.}
{\True Khi giải cần xác định đúng phương, chiều, điều kiện áp dụng và đơn vị.}
\loigiai{
\begin{itemchoice}
\item (Đúng) Vectơ cảm ứng từ tại một điểm có hướng trùng với hướng của từ trường tại điểm đó; đơn vị là tesla (T). (Vì): Vectơ cảm ứng từ tại một điểm có hướng trùng với hướng của từ trường tại điểm đó; đơn vị là tesla (T). b) Sai: Kết quả không phụ thuộc dữ kiện và có thể bỏ qua điều kiện. c) Sai: Cảm ứng từ là đại lượng vô hướng và có đơn vị niutơn. d) Đúng: Khi giải cần xác định đúng phương, chiều, điều kiện áp dụng và đơn vị. Kiến thức cốt lõi: Vectơ cảm ứng từ tại một điểm có hướng trùng với hướng của từ trường tại điểm đó; đơn vị là tesla (T
\item (Sai) Kết quả không phụ thuộc dữ kiện và có thể bỏ qua điều kiện.
\item (Sai) Cảm ứng từ là đại lượng vô hướng và có đơn vị niutơn.
\item (Đúng) Khi giải cần xác định đúng phương, chiều, điều kiện áp dụng và đơn vị.
\end{itemchoice}
}
\end{ex}

% ===== Câu 84 =====
% ID: L12C3B15-16-TL
% Mức: VD
\begin{bt}
% ID: L12C3B15-16-TL
% Mức: VD
Đề xuất các bước giải một bài toán thuộc dạng Nhận biết lực từ và đặc trưng của vectơ cảm ứng từ.
\loigiai{
Cần nêu được: Vectơ cảm ứng từ tại một điểm có hướng trùng với hướng của từ trường tại điểm đó; đơn vị là tesla (T). Khi vận dụng phải xác định đúng dữ kiện, phương chiều nếu có, điều kiện áp dụng, hệ đơn vị và kiểm tra tính hợp lí. Nhận định “Cảm ứng từ là đại lượng vô hướng và có đơn vị niutơn.” sai vì trái với kiến thức nêu trên.
}
\end{bt}

% ===== Câu 85 =====
% ID: L12C3B15-16-TLN
% Mức: TH
\begin{ex}
% ID: L12C3B15-16-TLN
% Mức: TH
Trong bốn phát biểu về Nhận biết lực từ và đặc trưng của vectơ cảm ứng từ, có bao nhiêu phát biểu đúng? (Chỉ ghi một số nguyên.) (1) Vectơ cảm ứng từ tại một điểm có hướng trùng với hướng của từ trường tại điểm đó; đơn vị là tesla (T). (2) Cần kiểm tra điều kiện áp dụng. (3) Cảm ứng từ là đại lượng vô hướng và có đơn vị niutơn. (4) Có thể bỏ qua đơn vị của kết quả.
\shortans{2}
\loigiai{
Đối chiếu từng phát biểu với kiến thức: Vectơ cảm ứng từ tại một điểm có hướng trùng với hướng của từ trường tại điểm đó; đơn vị là tesla (T). Số phát biểu đúng là 2.
}
\end{ex}

% ===== Câu 86 =====
% ID: L12C3B15-16-TN
% Mức: TH
\dangbt{Cân bằng và chuyển động của dây dẫn mang dòng điện trong từ trường}
\begin{ex}
% ID: L12C3B15-16-TN
% Mức: TH
Một học sinh ôn tập Cân bằng và chuyển động của dây dẫn mang dòng điện trong từ trường?
\choice
{\True Dây cân bằng khi tổng vectơ các lực, gồm lực từ và các lực cơ học, bằng vectơ không.}
{Dây cân bằng chỉ cần lực từ khác không, không cần xét trọng lực hay lực căng.}
{Có thể bỏ qua phương, chiều và đơn vị của các đại lượng trong mọi trường hợp.}
{Kết quả không phụ thuộc dữ kiện và điều kiện áp dụng.}
\loigiai{
Dây cân bằng khi tổng vectơ các lực, gồm lực từ và các lực cơ học, bằng vectơ không. Vì vậy phương án $A$ đúng; các phương án còn lại trái với kiến thức hoặc điều kiện áp dụng.
}
\end{ex}

% ===== Câu 87 =====
% ID: L12C3B15-17-DS
% Mức: VDC
\dangbt{Nhận biết lực từ và đặc trưng của vectơ cảm ứng từ}
\begin{ex}
% ID: L12C3B15-17-DS
% Mức: VDC
Ở một mặt phẳng nằm ngang cố định trong một từ trường đều có các đường sức từ nằm ngang, đặt một đoạn dây dẫn thẳng có dòng điện không đổi sao cho chiều của dòng điện hợp với chiều của đường sức một góc $\alpha$. Khi thay đổi góc $\alpha$, đoạn dây vẫn nằm trong mặt phẳng cố định nói trên.
\choiceTF
{Khi thay đổi góc $\alpha$ thì phương của lực từ tác dụng lên đoạn dây cũng thay đổi.}
{\True Khi đặt đoạn dây song song với các đường sức từ, không có lực từ tác dụng lên đoạn dây.}
{\True Lực từ tác dụng lên đoạn dây có phương vuông góc với cả đoạn dây và đường sức từ.}
{\True Lực từ tác dụng lên đoạn dây có độ lớn cực đại khi đoạn dây vuông góc với đường sức từ.}
\loigiai{
\begin{itemchoice}
\item (Sai) Khi thay đổi góc $\alpha$ thì phương của lực từ tác dụng lên đoạn dây cũng thay đổi. (Vì): — — — Lực từ luôn vuông góc với mặt phẳng chứa $\vec{B}$ và đoạn dây. → Sai
\item (Đúng) Khi đặt đoạn dây song song với các đường sức từ, không có lực từ tác dụng lên đoạn dây. (Vì): — — — Khi đoạn dây song song với đường sức từ, lực từ tác dụng lên đoạn dây bằng không. → Đúng
\item (Đúng) Lực từ tác dụng lên đoạn dây có phương vuông góc với cả đoạn dây và đường sức từ. (Vì): — — — Lực từ tác dụng lên đoạn dây luôn vuông góc với mặt phẳng chứa $\vec{B}$ và đoạn dây. → Đúng
\item (Đúng) Lực từ tác dụng lên đoạn dây có độ lớn cực đại khi đoạn dây vuông góc với đường sức từ. (Vì): — — — Khi góc hợp bởi chiều dòng điện và chiều đường sức từ bằng $90^\circ$, lực từ có độ lớn cực đại. → Đúng
\end{itemchoice}
}
\end{ex}

% ===== Câu 88 =====
% ID: L12C3B15-17-TL
% Mức: VD
\begin{bt}
% ID: L12C3B15-17-TL
% Mức: VD
Nêu một tình huống thực tiễn liên quan đến Nhận biết lực từ và đặc trưng của vectơ cảm ứng từ và giải thích bằng kiến thức Vật lí.
\loigiai{
Cần nêu được: Vectơ cảm ứng từ tại một điểm có hướng trùng với hướng của từ trường tại điểm đó; đơn vị là tesla (T). Khi vận dụng phải xác định đúng dữ kiện, phương chiều nếu có, điều kiện áp dụng, hệ đơn vị và kiểm tra tính hợp lí. Nhận định “Cảm ứng từ là đại lượng vô hướng và có đơn vị niutơn.” sai vì trái với kiến thức nêu trên.
}
\end{bt}

% ===== Câu 89 =====
% ID: L12C3B15-17-TLN
% Mức: VD
\begin{ex}
% ID: L12C3B15-17-TLN
% Mức: VD
Trong bốn phát biểu về Nhận biết lực từ và đặc trưng của vectơ cảm ứng từ, có bao nhiêu phát biểu đúng? (Chỉ ghi một số nguyên.) (1) Cảm ứng từ là đại lượng vô hướng và có đơn vị niutơn. (2) Vectơ cảm ứng từ tại một điểm có hướng trùng với hướng của từ trường tại điểm đó; đơn vị là tesla (T). (3) Cảm ứng từ là đại lượng vô hướng và có đơn vị niutơn. (4) Có thể bỏ qua đơn vị của kết quả.
\shortans{2}
\loigiai{
Đối chiếu từng phát biểu với kiến thức: Vectơ cảm ứng từ tại một điểm có hướng trùng với hướng của từ trường tại điểm đó; đơn vị là tesla (T). Số phát biểu đúng là 2.
}
\end{ex}

% ===== Câu 90 =====
% ID: L12C3B15-17-TN
% Mức: TH
\dangbt{Cân bằng và chuyển động của dây dẫn mang dòng điện trong từ trường}
\begin{ex}
% ID: L12C3B15-17-TN
% Mức: TH
Chọn phát biểu không trái với kiến thức về Cân bằng và chuyển động của dây dẫn mang dòng điện trong từ trường?
\choice
{Kết quả không phụ thuộc dữ kiện và điều kiện áp dụng.}
{\True Dây cân bằng khi tổng vectơ các lực, gồm lực từ và các lực cơ học, bằng vectơ không.}
{Dây cân bằng chỉ cần lực từ khác không, không cần xét trọng lực hay lực căng.}
{Có thể bỏ qua phương, chiều và đơn vị của các đại lượng trong mọi trường hợp.}
\loigiai{
Dây cân bằng khi tổng vectơ các lực, gồm lực từ và các lực cơ học, bằng vectơ không. Vì vậy phương án $B$ đúng; các phương án còn lại trái với kiến thức hoặc điều kiện áp dụng.
}
\end{ex}

% ===== Câu 91 =====
% ID: L12C3B15-18-DS
% Mức: TH
\dangbt{Tính lực từ tác dụng lên đoạn dây dẫn mang dòng điện}
\begin{ex}
% ID: L12C3B15-18-DS
% Mức: TH
Xét Tính lực từ tác dụng lên đoạn dây dẫn mang dòng điện (Tình huống 1). Hãy xác định tính đúng, sai của bốn phát biểu sau.
\choiceTF
{\True Độ lớn lực từ lên đoạn dây thẳng dài l mang dòng $I$ trong từ trường đều là $F$ = BIl sin(alpha).}
{Lực từ được tính bằng $F$ = BI/l và không phụ thuộc góc.}
{\True Khi giải cần xác định đúng phương, chiều, điều kiện áp dụng và đơn vị.}
{Kết quả không phụ thuộc dữ kiện và có thể bỏ qua điều kiện.}
\loigiai{
\begin{itemchoice}
\item (Đúng) Độ lớn lực từ lên đoạn dây thẳng dài l mang dòng $I$ trong từ trường đều là $F$ = BIl sin(alpha). (Vì): lpha
\item (Sai) Lực từ được tính bằng $F$ = BI/l và không phụ thuộc góc.
\item (Đúng) Khi giải cần xác định đúng phương, chiều, điều kiện áp dụng và đơn vị.
\item (Sai) Kết quả không phụ thuộc dữ kiện và có thể bỏ qua điều kiện.
\end{itemchoice}
}
\end{ex}

% ===== Câu 92 =====
% ID: L12C3B15-18-TL
% Mức: VDC
\dangbt{Nhận biết lực từ và đặc trưng của vectơ cảm ứng từ}
\begin{bt}
% ID: L12C3B15-18-TL
% Mức: VDC
So sánh nhận định đúng “Vectơ cảm ứng từ tại một điểm có hướng trùng với hướng của từ trường tại điểm đó; đơn vị là tesla (T).” với nhận định sai “Cảm ứng từ là đại lượng vô hướng và có đơn vị niutơn.”; chỉ rõ căn cứ Vật lí.
\loigiai{
Cần nêu được: Vectơ cảm ứng từ tại một điểm có hướng trùng với hướng của từ trường tại điểm đó; đơn vị là tesla (T). Khi vận dụng phải xác định đúng dữ kiện, phương chiều nếu có, điều kiện áp dụng, hệ đơn vị và kiểm tra tính hợp lí. Nhận định “Cảm ứng từ là đại lượng vô hướng và có đơn vị niutơn.” sai vì trái với kiến thức nêu trên.
}
\end{bt}

% ===== Câu 93 =====
% ID: L12C3B15-18-TLN
% Mức: VD
\begin{ex}
% ID: L12C3B15-18-TLN
% Mức: VD
Trong bốn phát biểu về Nhận biết lực từ và đặc trưng của vectơ cảm ứng từ, có bao nhiêu phát biểu đúng? (Chỉ ghi một số nguyên.) (1) Vectơ cảm ứng từ tại một điểm có hướng trùng với hướng của từ trường tại điểm đó; đơn vị là tesla (T). (2) Cần kiểm tra điều kiện áp dụng. (3) Cảm ứng từ là đại lượng vô hướng và có đơn vị niutơn. (4) Có thể bỏ qua đơn vị của kết quả.
\shortans{2}
\loigiai{
Đối chiếu từng phát biểu với kiến thức: Vectơ cảm ứng từ tại một điểm có hướng trùng với hướng của từ trường tại điểm đó; đơn vị là tesla (T). Số phát biểu đúng là 2.
}
\end{ex}

% ===== Câu 94 =====
% ID: L12C3B15-18-TN
% Mức: VD
\dangbt{Cân bằng và chuyển động của dây dẫn mang dòng điện trong từ trường}
\begin{ex}
% ID: L12C3B15-18-TN
% Mức: VD
Cơ sở Vật lí đúng dùng để giải Cân bằng và chuyển động của dây dẫn mang dòng điện trong từ trường?
\choice
{Có thể bỏ qua phương, chiều và đơn vị của các đại lượng trong mọi trường hợp.}
{Kết quả không phụ thuộc dữ kiện và điều kiện áp dụng.}
{\True Dây cân bằng khi tổng vectơ các lực, gồm lực từ và các lực cơ học, bằng vectơ không.}
{Dây cân bằng chỉ cần lực từ khác không, không cần xét trọng lực hay lực căng.}
\loigiai{
Dây cân bằng khi tổng vectơ các lực, gồm lực từ và các lực cơ học, bằng vectơ không. Vì vậy phương án $C$ đúng; các phương án còn lại trái với kiến thức hoặc điều kiện áp dụng.
}
\end{ex}

% ===== Câu 95 =====
% ID: L12C3B15-19-DS
% Mức: TH
\dangbt{Tính lực từ tác dụng lên đoạn dây dẫn mang dòng điện}
\begin{ex}
% ID: L12C3B15-19-DS
% Mức: TH
Xét Tính lực từ tác dụng lên đoạn dây dẫn mang dòng điện (Tình huống 2). Hãy xác định tính đúng, sai của bốn phát biểu sau.
\choiceTF
{Lực từ được tính bằng $F$ = BI/l và không phụ thuộc góc.}
{\True Độ lớn lực từ lên đoạn dây thẳng dài l mang dòng $I$ trong từ trường đều là $F$ = BIl sin(alpha).}
{Kết quả không phụ thuộc dữ kiện và có thể bỏ qua điều kiện.}
{\True Khi giải cần xác định đúng phương, chiều, điều kiện áp dụng và đơn vị.}
\loigiai{
\begin{itemchoice}
\item (Sai) Lực từ được tính bằng $F$ = BI/l và không phụ thuộc góc. (Vì): lpha
\item (Đúng) Độ lớn lực từ lên đoạn dây thẳng dài l mang dòng $I$ trong từ trường đều là $F$ = BIl sin(alpha).
\item (Sai) Kết quả không phụ thuộc dữ kiện và có thể bỏ qua điều kiện.
\item (Đúng) Khi giải cần xác định đúng phương, chiều, điều kiện áp dụng và đơn vị.
\end{itemchoice}
}
\end{ex}

% ===== Câu 96 =====
% ID: L12C3B15-19-TL
% Mức: TH
\begin{bt}
% ID: L12C3B15-19-TL
% Mức: TH
Trình bày kiến thức cốt lõi của Tính lực từ tác dụng lên đoạn dây dẫn mang dòng điện và nêu điều kiện áp dụng.
\loigiai{
Cần nêu được: Độ lớn lực từ lên đoạn dây thẳng dài l mang dòng $I$ trong từ trường đều là $F$ = BIl sin(alpha). Khi vận dụng phải xác định đúng dữ kiện, phương chiều nếu có, điều kiện áp dụng, hệ đơn vị và kiểm tra tính hợp lí. Nhận định “Lực từ được tính bằng $F$ = BI/l và không phụ thuộc góc.” sai vì trái với kiến thức nêu trên.
}
\end{bt}

% ===== Câu 97 =====
% ID: L12C3B15-19-TLN
% Mức: VDC
\dangbt{Nhận biết lực từ và đặc trưng của vectơ cảm ứng từ}
\begin{ex}
% ID: L12C3B15-19-TLN
% Mức: VDC
Trong bốn phát biểu về Nhận biết lực từ và đặc trưng của vectơ cảm ứng từ, có bao nhiêu phát biểu đúng? (Chỉ ghi một số nguyên.) (1) Cảm ứng từ là đại lượng vô hướng và có đơn vị niutơn. (2) Vectơ cảm ứng từ tại một điểm có hướng trùng với hướng của từ trường tại điểm đó; đơn vị là tesla (T). (3) Phải dùng đúng định luật cho tình huống. (4) Vectơ cảm ứng từ tại một điểm có hướng trùng với hướng của từ trường tại điểm đó; đơn vị là tesla (T).
\shortans{3}
\loigiai{
Đối chiếu từng phát biểu với kiến thức: Vectơ cảm ứng từ tại một điểm có hướng trùng với hướng của từ trường tại điểm đó; đơn vị là tesla (T). Số phát biểu đúng là 3.
}
\end{ex}

% ===== Câu 98 =====
% ID: L12C3B15-19-TN
% Mức: VD
\dangbt{Cân bằng và chuyển động của dây dẫn mang dòng điện trong từ trường}
\begin{ex}
% ID: L12C3B15-19-TN
% Mức: VD
Trong một bài thực hành về Cân bằng và chuyển động của dây dẫn mang dòng điện trong từ trường?
\choice
{Dây cân bằng chỉ cần lực từ khác không, không cần xét trọng lực hay lực căng.}
{Có thể bỏ qua phương, chiều và đơn vị của các đại lượng trong mọi trường hợp.}
{Kết quả không phụ thuộc dữ kiện và điều kiện áp dụng.}
{\True Dây cân bằng khi tổng vectơ các lực, gồm lực từ và các lực cơ học, bằng vectơ không.}
\loigiai{
Dây cân bằng khi tổng vectơ các lực, gồm lực từ và các lực cơ học, bằng vectơ không. Vì vậy phương án $D$ đúng; các phương án còn lại trái với kiến thức hoặc điều kiện áp dụng.
}
\end{ex}

% ===== Câu 99 =====
% ID: L12C3B15-20-DS
% Mức: VD
\dangbt{Tính lực từ tác dụng lên đoạn dây dẫn mang dòng điện}
\begin{ex}
% ID: L12C3B15-20-DS
% Mức: VD
Xét Tính lực từ tác dụng lên đoạn dây dẫn mang dòng điện (Tình huống 3). Hãy xác định tính đúng, sai của bốn phát biểu sau.
\choiceTF
{\True Khi giải cần xác định đúng phương, chiều, điều kiện áp dụng và đơn vị.}
{Kết quả không phụ thuộc dữ kiện và có thể bỏ qua điều kiện.}
{\True Độ lớn lực từ lên đoạn dây thẳng dài l mang dòng $I$ trong từ trường đều là $F$ = BIl sin(alpha).}
{Lực từ được tính bằng $F$ = BI/l và không phụ thuộc góc.}
\loigiai{
\begin{itemchoice}
\item (Đúng) Khi giải cần xác định đúng phương, chiều, điều kiện áp dụng và đơn vị. (Vì): lpha
\item (Sai) Kết quả không phụ thuộc dữ kiện và có thể bỏ qua điều kiện.
\item (Đúng) Độ lớn lực từ lên đoạn dây thẳng dài l mang dòng $I$ trong từ trường đều là $F$ = BIl sin(alpha).
\item (Sai) Lực từ được tính bằng $F$ = BI/l và không phụ thuộc góc.
\end{itemchoice}
}
\end{ex}

% ===== Câu 100 =====
% ID: L12C3B15-20-TL
% Mức: TH
\begin{bt}
% ID: L12C3B15-20-TL
% Mức: TH
Giải thích vì sao nhận định sau đúng: “Độ lớn lực từ lên đoạn dây thẳng dài l mang dòng $I$ trong từ trường đều là $F$ = BIl sin(alpha).”
\loigiai{
Cần nêu được: Độ lớn lực từ lên đoạn dây thẳng dài l mang dòng $I$ trong từ trường đều là $F$ = BIl sin(alpha). Khi vận dụng phải xác định đúng dữ kiện, phương chiều nếu có, điều kiện áp dụng, hệ đơn vị và kiểm tra tính hợp lí. Nhận định “Lực từ được tính bằng $F$ = BI/l và không phụ thuộc góc.” sai vì trái với kiến thức nêu trên.
}
\end{bt}

% ===== Câu 101 =====
% ID: L12C3B15-20-TLN
% Mức: TH
\begin{ex}
% ID: L12C3B15-20-TLN
% Mức: TH
Trong bốn phát biểu về Tính lực từ tác dụng lên đoạn dây dẫn mang dòng điện, có bao nhiêu phát biểu đúng? (Chỉ ghi một số nguyên.) (1) Độ lớn lực từ lên đoạn dây thẳng dài l mang dòng $I$ trong từ trường đều là $F$ = BIl sin(alpha). (2) Cần kiểm tra điều kiện áp dụng. (3) Lực từ được tính bằng $F$ = BI/l và không phụ thuộc góc. (4) Có thể bỏ qua đơn vị của kết quả.
\shortans{2}
\loigiai{
Đối chiếu từng phát biểu với kiến thức: Độ lớn lực từ lên đoạn dây thẳng dài l mang dòng $I$ trong từ trường đều là $F$ = BIl sin(alpha). Số phát biểu đúng là 2.
}
\end{ex}

% ===== Câu 102 =====
% ID: L12C3B15-20-TN
% Mức: VD
\dangbt{Cân bằng và chuyển động của dây dẫn mang dòng điện trong từ trường}
\begin{ex}
% ID: L12C3B15-20-TN
% Mức: VD
Kết luận đúng khi vận dụng Cân bằng và chuyển động của dây dẫn mang dòng điện trong từ trường?
\choice
{\True Dây cân bằng khi tổng vectơ các lực, gồm lực từ và các lực cơ học, bằng vectơ không.}
{Dây cân bằng chỉ cần lực từ khác không, không cần xét trọng lực hay lực căng.}
{Có thể bỏ qua phương, chiều và đơn vị của các đại lượng trong mọi trường hợp.}
{Kết quả không phụ thuộc dữ kiện và điều kiện áp dụng.}
\loigiai{
Dây cân bằng khi tổng vectơ các lực, gồm lực từ và các lực cơ học, bằng vectơ không. Vì vậy phương án $A$ đúng; các phương án còn lại trái với kiến thức hoặc điều kiện áp dụng.
}
\end{ex}

% ===== Câu 103 =====
% ID: L12C3B15-21-DS
% Mức: VD
\dangbt{Tính lực từ tác dụng lên đoạn dây dẫn mang dòng điện}
\begin{ex}
% ID: L12C3B15-21-DS
% Mức: VD
Xét Tính lực từ tác dụng lên đoạn dây dẫn mang dòng điện (Tình huống 4). Hãy xác định tính đúng, sai của bốn phát biểu sau.
\choiceTF
{Kết quả không phụ thuộc dữ kiện và có thể bỏ qua điều kiện.}
{\True Khi giải cần xác định đúng phương, chiều, điều kiện áp dụng và đơn vị.}
{Lực từ được tính bằng $F$ = BI/l và không phụ thuộc góc.}
{\True Độ lớn lực từ lên đoạn dây thẳng dài l mang dòng $I$ trong từ trường đều là $F$ = BIl sin(alpha).}
\loigiai{
\begin{itemchoice}
\item (Sai) Kết quả không phụ thuộc dữ kiện và có thể bỏ qua điều kiện. (Vì): lpha
\item (Đúng) Khi giải cần xác định đúng phương, chiều, điều kiện áp dụng và đơn vị.
\item (Sai) Lực từ được tính bằng $F$ = BI/l và không phụ thuộc góc.
\item (Đúng) Độ lớn lực từ lên đoạn dây thẳng dài l mang dòng $I$ trong từ trường đều là $F$ = BIl sin(alpha).
\end{itemchoice}
}
\end{ex}

% ===== Câu 104 =====
% ID: L12C3B15-21-TL
% Mức: VD
\begin{bt}
% ID: L12C3B15-21-TL
% Mức: VD
Phân tích sai lầm trong nhận định: “Lực từ được tính bằng $F$ = BI/l và không phụ thuộc góc.”
\loigiai{
Cần nêu được: Độ lớn lực từ lên đoạn dây thẳng dài l mang dòng $I$ trong từ trường đều là $F$ = BIl sin(alpha). Khi vận dụng phải xác định đúng dữ kiện, phương chiều nếu có, điều kiện áp dụng, hệ đơn vị và kiểm tra tính hợp lí. Nhận định “Lực từ được tính bằng $F$ = BI/l và không phụ thuộc góc.” sai vì trái với kiến thức nêu trên.
}
\end{bt}

% ===== Câu 105 =====
% ID: L12C3B15-21-TLN
% Mức: TH
\begin{ex}
% ID: L12C3B15-21-TLN
% Mức: TH
Trong bốn phát biểu về Tính lực từ tác dụng lên đoạn dây dẫn mang dòng điện, có bao nhiêu phát biểu đúng? (Chỉ ghi một số nguyên.) (1) Lực từ được tính bằng $F$ = BI/l và không phụ thuộc góc. (2) Độ lớn lực từ lên đoạn dây thẳng dài l mang dòng $I$ trong từ trường đều là $F$ = BIl sin(alpha). (3) Lực từ được tính bằng $F$ = BI/l và không phụ thuộc góc. (4) Có thể bỏ qua đơn vị của kết quả.
\shortans{1}
\loigiai{
Đối chiếu từng phát biểu với kiến thức: Độ lớn lực từ lên đoạn dây thẳng dài l mang dòng $I$ trong từ trường đều là $F$ = BIl sin(alpha). Số phát biểu đúng là 1.
}
\end{ex}

% ===== Câu 106 =====
% ID: L12C3B15-21-TN
% Mức: NB
\dangbt{Nhận biết lực từ và đặc trưng của vectơ cảm ứng từ}
\begin{ex}
% ID: L12C3B15-21-TN
% Mức: NB
Phát biểu nào sau đây đúng về Nhận biết lực từ và đặc trưng của vectơ cảm ứng từ?
\choice
{Cảm ứng từ là đại lượng vô hướng và có đơn vị niutơn.}
{Có thể bỏ qua phương, chiều và đơn vị của các đại lượng trong mọi trường hợp.}
{Kết quả không phụ thuộc dữ kiện và điều kiện áp dụng.}
{\True Vectơ cảm ứng từ tại một điểm có hướng trùng với hướng của từ trường tại điểm đó; đơn vị là tesla (T).}
\loigiai{
Vectơ cảm ứng từ tại một điểm có hướng trùng với hướng của từ trường tại điểm đó; đơn vị là tesla (T). Vì vậy phương án $D$ đúng; các phương án còn lại trái với kiến thức hoặc điều kiện áp dụng.
}
\end{ex}

% ===== Câu 107 =====
% ID: L12C3B15-22-DS
% Mức: VD
\dangbt{Tính lực từ tác dụng lên đoạn dây dẫn mang dòng điện}
\begin{ex}
% ID: L12C3B15-22-DS
% Mức: VD
Xét Tính lực từ tác dụng lên đoạn dây dẫn mang dòng điện (Tình huống 5). Hãy xác định tính đúng, sai của bốn phát biểu sau.
\choiceTF
{\True Độ lớn lực từ lên đoạn dây thẳng dài l mang dòng $I$ trong từ trường đều là $F$ = BIl sin(alpha).}
{Kết quả không phụ thuộc dữ kiện và có thể bỏ qua điều kiện.}
{Lực từ được tính bằng $F$ = BI/l và không phụ thuộc góc.}
{\True Khi giải cần xác định đúng phương, chiều, điều kiện áp dụng và đơn vị.}
\loigiai{
\begin{itemchoice}
\item (Đúng) Độ lớn lực từ lên đoạn dây thẳng dài l mang dòng $I$ trong từ trường đều là $F$ = BIl sin(alpha). (Vì): lpha
\item (Sai) Kết quả không phụ thuộc dữ kiện và có thể bỏ qua điều kiện.
\item (Sai) Lực từ được tính bằng $F$ = BI/l và không phụ thuộc góc.
\item (Đúng) Khi giải cần xác định đúng phương, chiều, điều kiện áp dụng và đơn vị.
\end{itemchoice}
}
\end{ex}

% ===== Câu 108 =====
% ID: L12C3B15-22-TL
% Mức: VD
\begin{bt}
% ID: L12C3B15-22-TL
% Mức: VD
Đề xuất các bước giải một bài toán thuộc dạng Tính lực từ tác dụng lên đoạn dây dẫn mang dòng điện.
\loigiai{
Cần nêu được: Độ lớn lực từ lên đoạn dây thẳng dài l mang dòng $I$ trong từ trường đều là $F$ = BIl sin(alpha). Khi vận dụng phải xác định đúng dữ kiện, phương chiều nếu có, điều kiện áp dụng, hệ đơn vị và kiểm tra tính hợp lí. Nhận định “Lực từ được tính bằng $F$ = BI/l và không phụ thuộc góc.” sai vì trái với kiến thức nêu trên.
}
\end{bt}

% ===== Câu 109 =====
% ID: L12C3B15-22-TLN
% Mức: TH
\begin{ex}
% ID: L12C3B15-22-TLN
% Mức: TH
Trong bốn phát biểu về Tính lực từ tác dụng lên đoạn dây dẫn mang dòng điện, có bao nhiêu phát biểu đúng? (Chỉ ghi một số nguyên.) (1) Độ lớn lực từ lên đoạn dây thẳng dài l mang dòng $I$ trong từ trường đều là $F$ = BIl sin(alpha). (2) Cần kiểm tra điều kiện áp dụng. (3) Lực từ được tính bằng $F$ = BI/l và không phụ thuộc góc. (4) Có thể bỏ qua đơn vị của kết quả.
\shortans{2}
\loigiai{
Đối chiếu từng phát biểu với kiến thức: Độ lớn lực từ lên đoạn dây thẳng dài l mang dòng $I$ trong từ trường đều là $F$ = BIl sin(alpha). Số phát biểu đúng là 2.
}
\end{ex}

% ===== Câu 110 =====
% ID: L12C3B15-22-TN
% Mức: NB
\dangbt{Nhận biết lực từ và đặc trưng của vectơ cảm ứng từ}
\begin{ex}
% ID: L12C3B15-22-TN
% Mức: NB
Trong nội dung Nhận biết lực từ và đặc trưng của vectơ cảm ứng từ?
\choice
{\True Vectơ cảm ứng từ tại một điểm có hướng trùng với hướng của từ trường tại điểm đó; đơn vị là tesla (T).}
{Cảm ứng từ là đại lượng vô hướng và có đơn vị niutơn.}
{Có thể bỏ qua phương, chiều và đơn vị của các đại lượng trong mọi trường hợp.}
{Kết quả không phụ thuộc dữ kiện và điều kiện áp dụng.}
\loigiai{
Vectơ cảm ứng từ tại một điểm có hướng trùng với hướng của từ trường tại điểm đó; đơn vị là tesla (T). Vì vậy phương án $A$ đúng; các phương án còn lại trái với kiến thức hoặc điều kiện áp dụng.
}
\end{ex}

% ===== Câu 111 =====
% ID: L12C3B15-23-DS
% Mức: TH
\dangbt{Xác định cảm ứng từ từ phép đo lực từ}
\begin{ex}
% ID: L12C3B15-23-DS
% Mức: TH
Xét Xác định cảm ứng từ từ phép đo lực từ (Tình huống 1). Hãy xác định tính đúng, sai của bốn phát biểu sau.
\choiceTF
{\True Khi dây vuông góc từ trường, cảm ứng từ được xác định bởi $B$ = F/(Il).}
{Khi dây vuông góc từ trường, cảm ứng từ bằng FIl.}
{\True Khi giải cần xác định đúng phương, chiều, điều kiện áp dụng và đơn vị.}
{Kết quả không phụ thuộc dữ kiện và có thể bỏ qua điều kiện.}
\loigiai{
\begin{itemchoice}
\item (Đúng) Khi dây vuông góc từ trường, cảm ứng từ được xác định bởi $B$ = F/(Il). (Vì): Khi dây vuông góc từ trường, cảm ứng từ được xác định bởi $B$ = F/(Il). b) Sai: Khi dây vuông góc từ trường, cảm ứng từ bằng FIl. c) Đúng: Khi giải cần xác định đúng phương, chiều, điều kiện áp dụng và đơn vị. d) Sai: Kết quả không phụ thuộc dữ kiện và có thể bỏ qua điều kiện. Kiến thức cốt lõi: Khi dây vuông góc từ trường, cảm ứng từ được xác định bởi $B$ = F/(Il
\item (Sai) Khi dây vuông góc từ trường, cảm ứng từ bằng FIl.
\item (Đúng) Khi giải cần xác định đúng phương, chiều, điều kiện áp dụng và đơn vị.
\item (Sai) Kết quả không phụ thuộc dữ kiện và có thể bỏ qua điều kiện.
\end{itemchoice}
}
\end{ex}

% ===== Câu 112 =====
% ID: L12C3B15-23-TL
% Mức: VD
\dangbt{Tính lực từ tác dụng lên đoạn dây dẫn mang dòng điện}
\begin{bt}
% ID: L12C3B15-23-TL
% Mức: VD
Nêu một tình huống thực tiễn liên quan đến Tính lực từ tác dụng lên đoạn dây dẫn mang dòng điện và giải thích bằng kiến thức Vật lí.
\loigiai{
Cần nêu được: Độ lớn lực từ lên đoạn dây thẳng dài l mang dòng $I$ trong từ trường đều là $F$ = BIl sin(alpha). Khi vận dụng phải xác định đúng dữ kiện, phương chiều nếu có, điều kiện áp dụng, hệ đơn vị và kiểm tra tính hợp lí. Nhận định “Lực từ được tính bằng $F$ = BI/l và không phụ thuộc góc.” sai vì trái với kiến thức nêu trên.
}
\end{bt}

% ===== Câu 113 =====
% ID: L12C3B15-23-TLN
% Mức: VD
\begin{ex}
% ID: L12C3B15-23-TLN
% Mức: VD
Trong bốn phát biểu về Tính lực từ tác dụng lên đoạn dây dẫn mang dòng điện, có bao nhiêu phát biểu đúng? (Chỉ ghi một số nguyên.) (1) Lực từ được tính bằng $F$ = BI/l và không phụ thuộc góc. (2) Độ lớn lực từ lên đoạn dây thẳng dài l mang dòng $I$ trong từ trường đều là $F$ = BIl sin(alpha). (3) Lực từ được tính bằng $F$ = BI/l và không phụ thuộc góc. (4) Có thể bỏ qua đơn vị của kết quả.
\shortans{2}
\loigiai{
Đối chiếu từng phát biểu với kiến thức: Độ lớn lực từ lên đoạn dây thẳng dài l mang dòng $I$ trong từ trường đều là $F$ = BIl sin(alpha). Số phát biểu đúng là 2.
}
\end{ex}

% ===== Câu 114 =====
% ID: L12C3B15-23-TN
% Mức: NB
\dangbt{Nhận biết lực từ và đặc trưng của vectơ cảm ứng từ}
\begin{ex}
% ID: L12C3B15-23-TN
% Mức: NB
Tình huống 3: chọn kết luận chính xác về Nhận biết lực từ và đặc trưng của vectơ cảm ứng từ?
\choice
{Kết quả không phụ thuộc dữ kiện và điều kiện áp dụng.}
{\True Vectơ cảm ứng từ tại một điểm có hướng trùng với hướng của từ trường tại điểm đó; đơn vị là tesla (T).}
{Cảm ứng từ là đại lượng vô hướng và có đơn vị niutơn.}
{Có thể bỏ qua phương, chiều và đơn vị của các đại lượng trong mọi trường hợp.}
\loigiai{
Vectơ cảm ứng từ tại một điểm có hướng trùng với hướng của từ trường tại điểm đó; đơn vị là tesla (T). Vì vậy phương án $B$ đúng; các phương án còn lại trái với kiến thức hoặc điều kiện áp dụng.
}
\end{ex}

% ===== Câu 115 =====
% ID: L12C3B15-24-DS
% Mức: TH
\dangbt{Xác định cảm ứng từ từ phép đo lực từ}
\begin{ex}
% ID: L12C3B15-24-DS
% Mức: TH
Xét Xác định cảm ứng từ từ phép đo lực từ (Tình huống 2). Hãy xác định tính đúng, sai của bốn phát biểu sau.
\choiceTF
{Khi dây vuông góc từ trường, cảm ứng từ bằng FIl.}
{\True Khi dây vuông góc từ trường, cảm ứng từ được xác định bởi $B$ = F/(Il).}
{Kết quả không phụ thuộc dữ kiện và có thể bỏ qua điều kiện.}
{\True Khi giải cần xác định đúng phương, chiều, điều kiện áp dụng và đơn vị.}
\loigiai{
\begin{itemchoice}
\item (Sai) Khi dây vuông góc từ trường, cảm ứng từ bằng FIl. (Vì): Khi dây vuông góc từ trường, cảm ứng từ bằng FIl. b) Đúng: Khi dây vuông góc từ trường, cảm ứng từ được xác định bởi $B$ = F/(Il). c) Sai: Kết quả không phụ thuộc dữ kiện và có thể bỏ qua điều kiện. d) Đúng: Khi giải cần xác định đúng phương, chiều, điều kiện áp dụng và đơn vị. Kiến thức cốt lõi: Khi dây vuông góc từ trường, cảm ứng từ được xác định bởi $B$ = F/(Il
\item (Đúng) Khi dây vuông góc từ trường, cảm ứng từ được xác định bởi $B$ = F/(Il).
\item (Sai) Kết quả không phụ thuộc dữ kiện và có thể bỏ qua điều kiện.
\item (Đúng) Khi giải cần xác định đúng phương, chiều, điều kiện áp dụng và đơn vị.
\end{itemchoice}
}
\end{ex}

% ===== Câu 116 =====
% ID: L12C3B15-24-TL
% Mức: VDC
\dangbt{Tính lực từ tác dụng lên đoạn dây dẫn mang dòng điện}
\begin{bt}
% ID: L12C3B15-24-TL
% Mức: VDC
So sánh nhận định đúng “Độ lớn lực từ lên đoạn dây thẳng dài l mang dòng $I$ trong từ trường đều là $F$ = BIl sin(alpha).” với nhận định sai “Lực từ được tính bằng $F$ = BI/l và không phụ thuộc góc.”; chỉ rõ căn cứ Vật lí.
\loigiai{
Cần nêu được: Độ lớn lực từ lên đoạn dây thẳng dài l mang dòng $I$ trong từ trường đều là $F$ = BIl sin(alpha). Khi vận dụng phải xác định đúng dữ kiện, phương chiều nếu có, điều kiện áp dụng, hệ đơn vị và kiểm tra tính hợp lí. Nhận định “Lực từ được tính bằng $F$ = BI/l và không phụ thuộc góc.” sai vì trái với kiến thức nêu trên.
}
\end{bt}

% ===== Câu 117 =====
% ID: L12C3B15-24-TLN
% Mức: VD
\begin{ex}
% ID: L12C3B15-24-TLN
% Mức: VD
Trong bốn phát biểu về Tính lực từ tác dụng lên đoạn dây dẫn mang dòng điện, có bao nhiêu phát biểu đúng? (Chỉ ghi một số nguyên.) (1) Độ lớn lực từ lên đoạn dây thẳng dài l mang dòng $I$ trong từ trường đều là $F$ = BIl sin(alpha). (2) Cần kiểm tra điều kiện áp dụng. (3) Lực từ được tính bằng $F$ = BI/l và không phụ thuộc góc. (4) Có thể bỏ qua đơn vị của kết quả.
\shortans{2}
\loigiai{
Đối chiếu từng phát biểu với kiến thức: Độ lớn lực từ lên đoạn dây thẳng dài l mang dòng $I$ trong từ trường đều là $F$ = BIl sin(alpha). Số phát biểu đúng là 2.
}
\end{ex}

% ===== Câu 118 =====
% ID: L12C3B15-24-TN
% Mức: TH
\dangbt{Nhận biết lực từ và đặc trưng của vectơ cảm ứng từ}
\begin{ex}
% ID: L12C3B15-24-TN
% Mức: TH
Nội dung nào mô tả đúng nhất Nhận biết lực từ và đặc trưng của vectơ cảm ứng từ?
\choice
{Có thể bỏ qua phương, chiều và đơn vị của các đại lượng trong mọi trường hợp.}
{Kết quả không phụ thuộc dữ kiện và điều kiện áp dụng.}
{\True Vectơ cảm ứng từ tại một điểm có hướng trùng với hướng của từ trường tại điểm đó; đơn vị là tesla (T).}
{Cảm ứng từ là đại lượng vô hướng và có đơn vị niutơn.}
\loigiai{
Vectơ cảm ứng từ tại một điểm có hướng trùng với hướng của từ trường tại điểm đó; đơn vị là tesla (T). Vì vậy phương án $C$ đúng; các phương án còn lại trái với kiến thức hoặc điều kiện áp dụng.
}
\end{ex}

% ===== Câu 119 =====
% ID: L12C3B15-25-DS
% Mức: VD
\dangbt{Xác định cảm ứng từ từ phép đo lực từ}
\begin{ex}
% ID: L12C3B15-25-DS
% Mức: VD
Xét Xác định cảm ứng từ từ phép đo lực từ (Tình huống 3). Hãy xác định tính đúng, sai của bốn phát biểu sau.
\choiceTF
{\True Khi giải cần xác định đúng phương, chiều, điều kiện áp dụng và đơn vị.}
{Kết quả không phụ thuộc dữ kiện và có thể bỏ qua điều kiện.}
{\True Khi dây vuông góc từ trường, cảm ứng từ được xác định bởi $B$ = F/(Il).}
{Khi dây vuông góc từ trường, cảm ứng từ bằng FIl.}
\loigiai{
\begin{itemchoice}
\item (Đúng) Khi giải cần xác định đúng phương, chiều, điều kiện áp dụng và đơn vị. (Vì): Khi giải cần xác định đúng phương, chiều, điều kiện áp dụng và đơn vị. b) Sai: Kết quả không phụ thuộc dữ kiện và có thể bỏ qua điều kiện. c) Đúng: Khi dây vuông góc từ trường, cảm ứng từ được xác định bởi $B$ = F/(Il). d) Sai: Khi dây vuông góc từ trường, cảm ứng từ bằng FIl. Kiến thức cốt lõi: Khi dây vuông góc từ trường, cảm ứng từ được xác định bởi $B$ = F/(Il
\item (Sai) Kết quả không phụ thuộc dữ kiện và có thể bỏ qua điều kiện.
\item (Đúng) Khi dây vuông góc từ trường, cảm ứng từ được xác định bởi $B$ = F/(Il).
\item (Sai) Khi dây vuông góc từ trường, cảm ứng từ bằng FIl.
\end{itemchoice}
}
\end{ex}

% ===== Câu 120 =====
% ID: L12C3B15-25-TL
% Mức: TH
\begin{bt}
% ID: L12C3B15-25-TL
% Mức: TH
Trình bày kiến thức cốt lõi của Xác định cảm ứng từ từ phép đo lực từ và nêu điều kiện áp dụng.
\loigiai{
Cần nêu được: Khi dây vuông góc từ trường, cảm ứng từ được xác định bởi $B$ = F/(Il). Khi vận dụng phải xác định đúng dữ kiện, phương chiều nếu có, điều kiện áp dụng, hệ đơn vị và kiểm tra tính hợp lí. Nhận định “Khi dây vuông góc từ trường, cảm ứng từ bằng FIl.” sai vì trái với kiến thức nêu trên.
}
\end{bt}

% ===== Câu 121 =====
% ID: L12C3B15-25-TLN
% Mức: VDC
\dangbt{Tính lực từ tác dụng lên đoạn dây dẫn mang dòng điện}
\begin{ex}
% ID: L12C3B15-25-TLN
% Mức: VDC
Trong bốn phát biểu về Tính lực từ tác dụng lên đoạn dây dẫn mang dòng điện, có bao nhiêu phát biểu đúng? (Chỉ ghi một số nguyên.) (1) Lực từ được tính bằng $F$ = BI/l và không phụ thuộc góc. (2) Độ lớn lực từ lên đoạn dây thẳng dài l mang dòng $I$ trong từ trường đều là $F$ = BIl sin(alpha). (3) Phải dùng đúng định luật cho tình huống. (4) Độ lớn lực từ lên đoạn dây thẳng dài l mang dòng $I$ trong từ trường đều là $F$ = BIl sin(alpha).
\shortans{3}
\loigiai{
Đối chiếu từng phát biểu với kiến thức: Độ lớn lực từ lên đoạn dây thẳng dài l mang dòng $I$ trong từ trường đều là $F$ = BIl sin(alpha). Số phát biểu đúng là 3.
}
\end{ex}

% ===== Câu 122 =====
% ID: L12C3B15-25-TN
% Mức: TH
\dangbt{Nhận biết lực từ và đặc trưng của vectơ cảm ứng từ}
\begin{ex}
% ID: L12C3B15-25-TN
% Mức: TH
Khi phân tích Nhận biết lực từ và đặc trưng của vectơ cảm ứng từ?
\choice
{Cảm ứng từ là đại lượng vô hướng và có đơn vị niutơn.}
{Có thể bỏ qua phương, chiều và đơn vị của các đại lượng trong mọi trường hợp.}
{Kết quả không phụ thuộc dữ kiện và điều kiện áp dụng.}
{\True Vectơ cảm ứng từ tại một điểm có hướng trùng với hướng của từ trường tại điểm đó; đơn vị là tesla (T).}
\loigiai{
Vectơ cảm ứng từ tại một điểm có hướng trùng với hướng của từ trường tại điểm đó; đơn vị là tesla (T). Vì vậy phương án $D$ đúng; các phương án còn lại trái với kiến thức hoặc điều kiện áp dụng.
}
\end{ex}

% ===== Câu 123 =====
% ID: L12C3B15-26-DS
% Mức: VD
\dangbt{Xác định cảm ứng từ từ phép đo lực từ}
\begin{ex}
% ID: L12C3B15-26-DS
% Mức: VD
Xét Xác định cảm ứng từ từ phép đo lực từ (Tình huống 4). Hãy xác định tính đúng, sai của bốn phát biểu sau.
\choiceTF
{Kết quả không phụ thuộc dữ kiện và có thể bỏ qua điều kiện.}
{\True Khi giải cần xác định đúng phương, chiều, điều kiện áp dụng và đơn vị.}
{Khi dây vuông góc từ trường, cảm ứng từ bằng FIl.}
{\True Khi dây vuông góc từ trường, cảm ứng từ được xác định bởi $B$ = F/(Il).}
\loigiai{
\begin{itemchoice}
\item (Sai) Kết quả không phụ thuộc dữ kiện và có thể bỏ qua điều kiện. (Vì): Kết quả không phụ thuộc dữ kiện và có thể bỏ qua điều kiện. b) Đúng: Khi giải cần xác định đúng phương, chiều, điều kiện áp dụng và đơn vị. c) Sai: Khi dây vuông góc từ trường, cảm ứng từ bằng FIl. d) Đúng: Khi dây vuông góc từ trường, cảm ứng từ được xác định bởi $B$ = F/(Il). Kiến thức cốt lõi: Khi dây vuông góc từ trường, cảm ứng từ được xác định bởi $B$ = F/(Il
\item (Đúng) Khi giải cần xác định đúng phương, chiều, điều kiện áp dụng và đơn vị.
\item (Sai) Khi dây vuông góc từ trường, cảm ứng từ bằng FIl.
\item (Đúng) Khi dây vuông góc từ trường, cảm ứng từ được xác định bởi $B$ = F/(Il).
\end{itemchoice}
}
\end{ex}

% ===== Câu 124 =====
% ID: L12C3B15-26-TL
% Mức: TH
\begin{bt}
% ID: L12C3B15-26-TL
% Mức: TH
Giải thích vì sao nhận định sau đúng: “Khi dây vuông góc từ trường, cảm ứng từ được xác định bởi $B$ = F/(Il).”
\loigiai{
Cần nêu được: Khi dây vuông góc từ trường, cảm ứng từ được xác định bởi $B$ = F/(Il). Khi vận dụng phải xác định đúng dữ kiện, phương chiều nếu có, điều kiện áp dụng, hệ đơn vị và kiểm tra tính hợp lí. Nhận định “Khi dây vuông góc từ trường, cảm ứng từ bằng FIl.” sai vì trái với kiến thức nêu trên.
}
\end{bt}

% ===== Câu 125 =====
% ID: L12C3B15-26-TLN
% Mức: TH
\begin{ex}
% ID: L12C3B15-26-TLN
% Mức: TH
Trong bốn phát biểu về Xác định cảm ứng từ từ phép đo lực từ, có bao nhiêu phát biểu đúng? (Chỉ ghi một số nguyên.) (1) Khi dây vuông góc từ trường, cảm ứng từ được xác định bởi $B$ = F/(Il). (2) Cần kiểm tra điều kiện áp dụng. (3) Khi dây vuông góc từ trường, cảm ứng từ bằng FIl. (4) Có thể bỏ qua đơn vị của kết quả.
\shortans{2}
\loigiai{
Đối chiếu từng phát biểu với kiến thức: Khi dây vuông góc từ trường, cảm ứng từ được xác định bởi $B$ = F/(Il). Số phát biểu đúng là 2.
}
\end{ex}

% ===== Câu 126 =====
% ID: L12C3B15-26-TN
% Mức: TH
\dangbt{Nhận biết lực từ và đặc trưng của vectơ cảm ứng từ}
\begin{ex}
% ID: L12C3B15-26-TN
% Mức: TH
Một học sinh ôn tập Nhận biết lực từ và đặc trưng của vectơ cảm ứng từ?
\choice
{\True Vectơ cảm ứng từ tại một điểm có hướng trùng với hướng của từ trường tại điểm đó; đơn vị là tesla (T).}
{Cảm ứng từ là đại lượng vô hướng và có đơn vị niutơn.}
{Có thể bỏ qua phương, chiều và đơn vị của các đại lượng trong mọi trường hợp.}
{Kết quả không phụ thuộc dữ kiện và điều kiện áp dụng.}
\loigiai{
Vectơ cảm ứng từ tại một điểm có hướng trùng với hướng của từ trường tại điểm đó; đơn vị là tesla (T). Vì vậy phương án $A$ đúng; các phương án còn lại trái với kiến thức hoặc điều kiện áp dụng.
}
\end{ex}

% ===== Câu 127 =====
% ID: L12C3B15-27-DS
% Mức: VD
\dangbt{Xác định cảm ứng từ từ phép đo lực từ}
\begin{ex}
% ID: L12C3B15-27-DS
% Mức: VD
Xét Xác định cảm ứng từ từ phép đo lực từ (Tình huống 5). Hãy xác định tính đúng, sai của bốn phát biểu sau.
\choiceTF
{\True Khi dây vuông góc từ trường, cảm ứng từ được xác định bởi $B$ = F/(Il).}
{Kết quả không phụ thuộc dữ kiện và có thể bỏ qua điều kiện.}
{Khi dây vuông góc từ trường, cảm ứng từ bằng FIl.}
{\True Khi giải cần xác định đúng phương, chiều, điều kiện áp dụng và đơn vị.}
\loigiai{
\begin{itemchoice}
\item (Đúng) Khi dây vuông góc từ trường, cảm ứng từ được xác định bởi $B$ = F/(Il). (Vì): Khi dây vuông góc từ trường, cảm ứng từ được xác định bởi $B$ = F/(Il). b) Sai: Kết quả không phụ thuộc dữ kiện và có thể bỏ qua điều kiện. c) Sai: Khi dây vuông góc từ trường, cảm ứng từ bằng FIl. d) Đúng: Khi giải cần xác định đúng phương, chiều, điều kiện áp dụng và đơn vị. Kiến thức cốt lõi: Khi dây vuông góc từ trường, cảm ứng từ được xác định bởi $B$ = F/(Il
\item (Sai) Kết quả không phụ thuộc dữ kiện và có thể bỏ qua điều kiện.
\item (Sai) Khi dây vuông góc từ trường, cảm ứng từ bằng FIl.
\item (Đúng) Khi giải cần xác định đúng phương, chiều, điều kiện áp dụng và đơn vị.
\end{itemchoice}
}
\end{ex}

% ===== Câu 128 =====
% ID: L12C3B15-27-TL
% Mức: VD
\begin{bt}
% ID: L12C3B15-27-TL
% Mức: VD
Phân tích sai lầm trong nhận định: “Khi dây vuông góc từ trường, cảm ứng từ bằng FIl.”
\loigiai{
Cần nêu được: Khi dây vuông góc từ trường, cảm ứng từ được xác định bởi $B$ = F/(Il). Khi vận dụng phải xác định đúng dữ kiện, phương chiều nếu có, điều kiện áp dụng, hệ đơn vị và kiểm tra tính hợp lí. Nhận định “Khi dây vuông góc từ trường, cảm ứng từ bằng FIl.” sai vì trái với kiến thức nêu trên.
}
\end{bt}

% ===== Câu 129 =====
% ID: L12C3B15-27-TLN
% Mức: TH
\begin{ex}
% ID: L12C3B15-27-TLN
% Mức: TH
Trong bốn phát biểu về Xác định cảm ứng từ từ phép đo lực từ, có bao nhiêu phát biểu đúng? (Chỉ ghi một số nguyên.) (1) Khi dây vuông góc từ trường, cảm ứng từ bằng FIl. (2) Khi dây vuông góc từ trường, cảm ứng từ được xác định bởi $B$ = F/(Il). (3) Khi dây vuông góc từ trường, cảm ứng từ bằng FIl. (4) Có thể bỏ qua đơn vị của kết quả.
\shortans{1}
\loigiai{
Đối chiếu từng phát biểu với kiến thức: Khi dây vuông góc từ trường, cảm ứng từ được xác định bởi $B$ = F/(Il). Số phát biểu đúng là 1.
}
\end{ex}

% ===== Câu 130 =====
% ID: L12C3B15-27-TN
% Mức: TH
\dangbt{Nhận biết lực từ và đặc trưng của vectơ cảm ứng từ}
\begin{ex}
% ID: L12C3B15-27-TN
% Mức: TH
Chọn phát biểu không trái với kiến thức về Nhận biết lực từ và đặc trưng của vectơ cảm ứng từ?
\choice
{Kết quả không phụ thuộc dữ kiện và điều kiện áp dụng.}
{\True Vectơ cảm ứng từ tại một điểm có hướng trùng với hướng của từ trường tại điểm đó; đơn vị là tesla (T).}
{Cảm ứng từ là đại lượng vô hướng và có đơn vị niutơn.}
{Có thể bỏ qua phương, chiều và đơn vị của các đại lượng trong mọi trường hợp.}
\loigiai{
Vectơ cảm ứng từ tại một điểm có hướng trùng với hướng của từ trường tại điểm đó; đơn vị là tesla (T). Vì vậy phương án $B$ đúng; các phương án còn lại trái với kiến thức hoặc điều kiện áp dụng.
}
\end{ex}

% ===== Câu 131 =====
% ID: L12C3B15-28-DS
% Mức: TH
\dangbt{Xác định phương, chiều lực từ bằng quy tắc bàn tay trái}
\begin{ex}
% ID: L12C3B15-28-DS
% Mức: TH
Xét Xác định phương, chiều lực từ bằng quy tắc bàn tay trái (Tình huống 1). Hãy xác định tính đúng, sai của bốn phát biểu sau.
\choiceTF
{\True Lực từ tác dụng lên đoạn dây có phương vuông góc với cả dòng điện và vectơ cảm ứng từ; chiều xác định bằng quy tắc bàn tay trái.}
{Lực từ luôn cùng chiều dòng điện trong dây.}
{\True Khi giải cần xác định đúng phương, chiều, điều kiện áp dụng và đơn vị.}
{Kết quả không phụ thuộc dữ kiện và có thể bỏ qua điều kiện.}
\loigiai{
\begin{itemchoice}
\item (Đúng) Lực từ tác dụng lên đoạn dây có phương vuông góc với cả dòng điện và vectơ cảm ứng từ; chiều xác định bằng quy tắc bàn tay trái. (Vì): Lực từ tác dụng lên đoạn dây có phương vuông góc với cả dòng điện và vectơ cảm ứng từ; chiều xác định bằng quy tắc bàn tay trái. b) Sai: Lực từ luôn cùng chiều dòng điện trong dây. c) Đúng: Khi giải cần xác định đúng phương, chiều, điều kiện áp dụng và đơn vị. d) Sai: Kết quả không phụ thuộc dữ kiện và có thể bỏ qua điều kiện. Kiến thức cốt lõi: Lực từ tác dụng lên đoạn dây có phương vuông góc với cả dòng điện và vectơ cảm ứng từ; chiều xác định bằng quy tắc bàn tay trái
\item (Sai) Lực từ luôn cùng chiều dòng điện trong dây.
\item (Đúng) Khi giải cần xác định đúng phương, chiều, điều kiện áp dụng và đơn vị.
\item (Sai) Kết quả không phụ thuộc dữ kiện và có thể bỏ qua điều kiện.
\end{itemchoice}
}
\end{ex}

% ===== Câu 132 =====
% ID: L12C3B15-28-TL
% Mức: VD
\dangbt{Xác định cảm ứng từ từ phép đo lực từ}
\begin{bt}
% ID: L12C3B15-28-TL
% Mức: VD
Đề xuất các bước giải một bài toán thuộc dạng Xác định cảm ứng từ từ phép đo lực từ.
\loigiai{
Cần nêu được: Khi dây vuông góc từ trường, cảm ứng từ được xác định bởi $B$ = F/(Il). Khi vận dụng phải xác định đúng dữ kiện, phương chiều nếu có, điều kiện áp dụng, hệ đơn vị và kiểm tra tính hợp lí. Nhận định “Khi dây vuông góc từ trường, cảm ứng từ bằng FIl.” sai vì trái với kiến thức nêu trên.
}
\end{bt}

% ===== Câu 133 =====
% ID: L12C3B15-28-TLN
% Mức: TH
\begin{ex}
% ID: L12C3B15-28-TLN
% Mức: TH
Trong bốn phát biểu về Xác định cảm ứng từ từ phép đo lực từ, có bao nhiêu phát biểu đúng? (Chỉ ghi một số nguyên.) (1) Khi dây vuông góc từ trường, cảm ứng từ được xác định bởi $B$ = F/(Il). (2) Cần kiểm tra điều kiện áp dụng. (3) Khi dây vuông góc từ trường, cảm ứng từ bằng FIl. (4) Có thể bỏ qua đơn vị của kết quả.
\shortans{2}
\loigiai{
Đối chiếu từng phát biểu với kiến thức: Khi dây vuông góc từ trường, cảm ứng từ được xác định bởi $B$ = F/(Il). Số phát biểu đúng là 2.
}
\end{ex}

% ===== Câu 134 =====
% ID: L12C3B15-28-TN
% Mức: VD
\dangbt{Nhận biết lực từ và đặc trưng của vectơ cảm ứng từ}
\begin{ex}
% ID: L12C3B15-28-TN
% Mức: VD
Cơ sở Vật lí đúng dùng để giải Nhận biết lực từ và đặc trưng của vectơ cảm ứng từ?
\choice
{Có thể bỏ qua phương, chiều và đơn vị của các đại lượng trong mọi trường hợp.}
{Kết quả không phụ thuộc dữ kiện và điều kiện áp dụng.}
{\True Vectơ cảm ứng từ tại một điểm có hướng trùng với hướng của từ trường tại điểm đó; đơn vị là tesla (T).}
{Cảm ứng từ là đại lượng vô hướng và có đơn vị niutơn.}
\loigiai{
Vectơ cảm ứng từ tại một điểm có hướng trùng với hướng của từ trường tại điểm đó; đơn vị là tesla (T). Vì vậy phương án $C$ đúng; các phương án còn lại trái với kiến thức hoặc điều kiện áp dụng.
}
\end{ex}

% ===== Câu 135 =====
% ID: L12C3B15-29-DS
% Mức: TH
\dangbt{Xác định phương, chiều lực từ bằng quy tắc bàn tay trái}
\begin{ex}
% ID: L12C3B15-29-DS
% Mức: TH
Xét Xác định phương, chiều lực từ bằng quy tắc bàn tay trái (Tình huống 2). Hãy xác định tính đúng, sai của bốn phát biểu sau.
\choiceTF
{Lực từ luôn cùng chiều dòng điện trong dây.}
{\True Lực từ tác dụng lên đoạn dây có phương vuông góc với cả dòng điện và vectơ cảm ứng từ; chiều xác định bằng quy tắc bàn tay trái.}
{Kết quả không phụ thuộc dữ kiện và có thể bỏ qua điều kiện.}
{\True Khi giải cần xác định đúng phương, chiều, điều kiện áp dụng và đơn vị.}
\loigiai{
\begin{itemchoice}
\item (Sai) Lực từ luôn cùng chiều dòng điện trong dây. (Vì): Lực từ luôn cùng chiều dòng điện trong dây. b) Đúng: Lực từ tác dụng lên đoạn dây có phương vuông góc với cả dòng điện và vectơ cảm ứng từ; chiều xác định bằng quy tắc bàn tay trái. c) Sai: Kết quả không phụ thuộc dữ kiện và có thể bỏ qua điều kiện. d) Đúng: Khi giải cần xác định đúng phương, chiều, điều kiện áp dụng và đơn vị. Kiến thức cốt lõi: Lực từ tác dụng lên đoạn dây có phương vuông góc với cả dòng điện và vectơ cảm ứng từ; chiều xác định bằng quy tắc bàn tay trái
\item (Đúng) Lực từ tác dụng lên đoạn dây có phương vuông góc với cả dòng điện và vectơ cảm ứng từ; chiều xác định bằng quy tắc bàn tay trái.
\item (Sai) Kết quả không phụ thuộc dữ kiện và có thể bỏ qua điều kiện.
\item (Đúng) Khi giải cần xác định đúng phương, chiều, điều kiện áp dụng và đơn vị.
\end{itemchoice}
}
\end{ex}

% ===== Câu 136 =====
% ID: L12C3B15-29-TL
% Mức: VD
\dangbt{Xác định cảm ứng từ từ phép đo lực từ}
\begin{bt}
% ID: L12C3B15-29-TL
% Mức: VD
Nêu một tình huống thực tiễn liên quan đến Xác định cảm ứng từ từ phép đo lực từ và giải thích bằng kiến thức Vật lí.
\loigiai{
Cần nêu được: Khi dây vuông góc từ trường, cảm ứng từ được xác định bởi $B$ = F/(Il). Khi vận dụng phải xác định đúng dữ kiện, phương chiều nếu có, điều kiện áp dụng, hệ đơn vị và kiểm tra tính hợp lí. Nhận định “Khi dây vuông góc từ trường, cảm ứng từ bằng FIl.” sai vì trái với kiến thức nêu trên.
}
\end{bt}

% ===== Câu 137 =====
% ID: L12C3B15-29-TLN
% Mức: VD
\begin{ex}
% ID: L12C3B15-29-TLN
% Mức: VD
Trong bốn phát biểu về Xác định cảm ứng từ từ phép đo lực từ, có bao nhiêu phát biểu đúng? (Chỉ ghi một số nguyên.) (1) Khi dây vuông góc từ trường, cảm ứng từ bằng FIl. (2) Khi dây vuông góc từ trường, cảm ứng từ được xác định bởi $B$ = F/(Il). (3) Khi dây vuông góc từ trường, cảm ứng từ bằng FIl. (4) Có thể bỏ qua đơn vị của kết quả.
\shortans{2}
\loigiai{
Đối chiếu từng phát biểu với kiến thức: Khi dây vuông góc từ trường, cảm ứng từ được xác định bởi $B$ = F/(Il). Số phát biểu đúng là 2.
}
\end{ex}

% ===== Câu 138 =====
% ID: L12C3B15-29-TN
% Mức: VD
\dangbt{Nhận biết lực từ và đặc trưng của vectơ cảm ứng từ}
\begin{ex}
% ID: L12C3B15-29-TN
% Mức: VD
Trong một bài thực hành về Nhận biết lực từ và đặc trưng của vectơ cảm ứng từ?
\choice
{Cảm ứng từ là đại lượng vô hướng và có đơn vị niutơn.}
{Có thể bỏ qua phương, chiều và đơn vị của các đại lượng trong mọi trường hợp.}
{Kết quả không phụ thuộc dữ kiện và điều kiện áp dụng.}
{\True Vectơ cảm ứng từ tại một điểm có hướng trùng với hướng của từ trường tại điểm đó; đơn vị là tesla (T).}
\loigiai{
Vectơ cảm ứng từ tại một điểm có hướng trùng với hướng của từ trường tại điểm đó; đơn vị là tesla (T). Vì vậy phương án $D$ đúng; các phương án còn lại trái với kiến thức hoặc điều kiện áp dụng.
}
\end{ex}

% ===== Câu 139 =====
% ID: L12C3B15-30-DS
% Mức: VD
\dangbt{Xác định phương, chiều lực từ bằng quy tắc bàn tay trái}
\begin{ex}
% ID: L12C3B15-30-DS
% Mức: VD
Xét Xác định phương, chiều lực từ bằng quy tắc bàn tay trái (Tình huống 3). Hãy xác định tính đúng, sai của bốn phát biểu sau.
\choiceTF
{\True Khi giải cần xác định đúng phương, chiều, điều kiện áp dụng và đơn vị.}
{Kết quả không phụ thuộc dữ kiện và có thể bỏ qua điều kiện.}
{\True Lực từ tác dụng lên đoạn dây có phương vuông góc với cả dòng điện và vectơ cảm ứng từ; chiều xác định bằng quy tắc bàn tay trái.}
{Lực từ luôn cùng chiều dòng điện trong dây.}
\loigiai{
\begin{itemchoice}
\item (Đúng) Khi giải cần xác định đúng phương, chiều, điều kiện áp dụng và đơn vị. (Vì): Khi giải cần xác định đúng phương, chiều, điều kiện áp dụng và đơn vị. b) Sai: Kết quả không phụ thuộc dữ kiện và có thể bỏ qua điều kiện. c) Đúng: Lực từ tác dụng lên đoạn dây có phương vuông góc với cả dòng điện và vectơ cảm ứng từ; chiều xác định bằng quy tắc bàn tay trái. d) Sai: Lực từ luôn cùng chiều dòng điện trong dây. Kiến thức cốt lõi: Lực từ tác dụng lên đoạn dây có phương vuông góc với cả dòng điện và vectơ cảm ứng từ; chiều xác định bằng quy tắc bàn tay trái
\item (Sai) Kết quả không phụ thuộc dữ kiện và có thể bỏ qua điều kiện.
\item (Đúng) Lực từ tác dụng lên đoạn dây có phương vuông góc với cả dòng điện và vectơ cảm ứng từ; chiều xác định bằng quy tắc bàn tay trái.
\item (Sai) Lực từ luôn cùng chiều dòng điện trong dây.
\end{itemchoice}
}
\end{ex}

% ===== Câu 140 =====
% ID: L12C3B15-30-TL
% Mức: VDC
\dangbt{Xác định cảm ứng từ từ phép đo lực từ}
\begin{bt}
% ID: L12C3B15-30-TL
% Mức: VDC
So sánh nhận định đúng “Khi dây vuông góc từ trường, cảm ứng từ được xác định bởi $B$ = F/(Il).” với nhận định sai “Khi dây vuông góc từ trường, cảm ứng từ bằng FIl.”; chỉ rõ căn cứ Vật lí.
\loigiai{
Cần nêu được: Khi dây vuông góc từ trường, cảm ứng từ được xác định bởi $B$ = F/(Il). Khi vận dụng phải xác định đúng dữ kiện, phương chiều nếu có, điều kiện áp dụng, hệ đơn vị và kiểm tra tính hợp lí. Nhận định “Khi dây vuông góc từ trường, cảm ứng từ bằng FIl.” sai vì trái với kiến thức nêu trên.
}
\end{bt}

% ===== Câu 141 =====
% ID: L12C3B15-30-TLN
% Mức: VD
\begin{ex}
% ID: L12C3B15-30-TLN
% Mức: VD
Trong bốn phát biểu về Xác định cảm ứng từ từ phép đo lực từ, có bao nhiêu phát biểu đúng? (Chỉ ghi một số nguyên.) (1) Khi dây vuông góc từ trường, cảm ứng từ được xác định bởi $B$ = F/(Il). (2) Cần kiểm tra điều kiện áp dụng. (3) Khi dây vuông góc từ trường, cảm ứng từ bằng FIl. (4) Có thể bỏ qua đơn vị của kết quả.
\shortans{2}
\loigiai{
Đối chiếu từng phát biểu với kiến thức: Khi dây vuông góc từ trường, cảm ứng từ được xác định bởi $B$ = F/(Il). Số phát biểu đúng là 2.
}
\end{ex}

% ===== Câu 142 =====
% ID: L12C3B15-30-TN
% Mức: VD
\dangbt{Nhận biết lực từ và đặc trưng của vectơ cảm ứng từ}
\begin{ex}
% ID: L12C3B15-30-TN
% Mức: VD
Kết luận đúng khi vận dụng Nhận biết lực từ và đặc trưng của vectơ cảm ứng từ?
\choice
{\True Vectơ cảm ứng từ tại một điểm có hướng trùng với hướng của từ trường tại điểm đó; đơn vị là tesla (T).}
{Cảm ứng từ là đại lượng vô hướng và có đơn vị niutơn.}
{Có thể bỏ qua phương, chiều và đơn vị của các đại lượng trong mọi trường hợp.}
{Kết quả không phụ thuộc dữ kiện và điều kiện áp dụng.}
\loigiai{
Vectơ cảm ứng từ tại một điểm có hướng trùng với hướng của từ trường tại điểm đó; đơn vị là tesla (T). Vì vậy phương án $A$ đúng; các phương án còn lại trái với kiến thức hoặc điều kiện áp dụng.
}
\end{ex}

% ===== Câu 143 =====
% ID: L12C3B15-31-DS
% Mức: VD
\dangbt{Xác định phương, chiều lực từ bằng quy tắc bàn tay trái}
\begin{ex}
% ID: L12C3B15-31-DS
% Mức: VD
Xét Xác định phương, chiều lực từ bằng quy tắc bàn tay trái (Tình huống 4). Hãy xác định tính đúng, sai của bốn phát biểu sau.
\choiceTF
{Kết quả không phụ thuộc dữ kiện và có thể bỏ qua điều kiện.}
{\True Khi giải cần xác định đúng phương, chiều, điều kiện áp dụng và đơn vị.}
{Lực từ luôn cùng chiều dòng điện trong dây.}
{\True Lực từ tác dụng lên đoạn dây có phương vuông góc với cả dòng điện và vectơ cảm ứng từ; chiều xác định bằng quy tắc bàn tay trái.}
\loigiai{
\begin{itemchoice}
\item (Sai) Kết quả không phụ thuộc dữ kiện và có thể bỏ qua điều kiện. (Vì): Kết quả không phụ thuộc dữ kiện và có thể bỏ qua điều kiện. b) Đúng: Khi giải cần xác định đúng phương, chiều, điều kiện áp dụng và đơn vị. c) Sai: Lực từ luôn cùng chiều dòng điện trong dây. d) Đúng: Lực từ tác dụng lên đoạn dây có phương vuông góc với cả dòng điện và vectơ cảm ứng từ; chiều xác định bằng quy tắc bàn tay trái. Kiến thức cốt lõi: Lực từ tác dụng lên đoạn dây có phương vuông góc với cả dòng điện và vectơ cảm ứng từ; chiều xác định bằng quy tắc bàn tay trái
\item (Đúng) Khi giải cần xác định đúng phương, chiều, điều kiện áp dụng và đơn vị.
\item (Sai) Lực từ luôn cùng chiều dòng điện trong dây.
\item (Đúng) Lực từ tác dụng lên đoạn dây có phương vuông góc với cả dòng điện và vectơ cảm ứng từ; chiều xác định bằng quy tắc bàn tay trái.
\end{itemchoice}
}
\end{ex}

% ===== Câu 144 =====
% ID: L12C3B15-31-TL
% Mức: TH
\begin{bt}
% ID: L12C3B15-31-TL
% Mức: TH
Trình bày kiến thức cốt lõi của Xác định phương, chiều lực từ bằng quy tắc bàn tay trái và nêu điều kiện áp dụng.
\loigiai{
Cần nêu được: Lực từ tác dụng lên đoạn dây có phương vuông góc với cả dòng điện và vectơ cảm ứng từ; chiều xác định bằng quy tắc bàn tay trái. Khi vận dụng phải xác định đúng dữ kiện, phương chiều nếu có, điều kiện áp dụng, hệ đơn vị và kiểm tra tính hợp lí. Nhận định “Lực từ luôn cùng chiều dòng điện trong dây.” sai vì trái với kiến thức nêu trên.
}
\end{bt}

% ===== Câu 145 =====
% ID: L12C3B15-31-TLN
% Mức: VDC
\dangbt{Xác định cảm ứng từ từ phép đo lực từ}
\begin{ex}
% ID: L12C3B15-31-TLN
% Mức: VDC
Trong bốn phát biểu về Xác định cảm ứng từ từ phép đo lực từ, có bao nhiêu phát biểu đúng? (Chỉ ghi một số nguyên.) (1) Khi dây vuông góc từ trường, cảm ứng từ bằng FIl. (2) Khi dây vuông góc từ trường, cảm ứng từ được xác định bởi $B$ = F/(Il). (3) Phải dùng đúng định luật cho tình huống. (4) Khi dây vuông góc từ trường, cảm ứng từ được xác định bởi $B$ = F/(Il).
\shortans{3}
\loigiai{
Đối chiếu từng phát biểu với kiến thức: Khi dây vuông góc từ trường, cảm ứng từ được xác định bởi $B$ = F/(Il). Số phát biểu đúng là 3.
}
\end{ex}

% ===== Câu 146 =====
% ID: L12C3B15-31-TN
% Mức: NB
\dangbt{Tính lực từ tác dụng lên đoạn dây dẫn mang dòng điện}
\begin{ex}
% ID: L12C3B15-31-TN
% Mức: NB
Phát biểu nào sau đây đúng về Tính lực từ tác dụng lên đoạn dây dẫn mang dòng điện?
\choice
{Lực từ được tính bằng $F$ = BI/l và không phụ thuộc góc.}
{Có thể bỏ qua phương, chiều và đơn vị của các đại lượng trong mọi trường hợp.}
{Kết quả không phụ thuộc dữ kiện và điều kiện áp dụng.}
{\True Độ lớn lực từ lên đoạn dây thẳng dài l mang dòng $I$ trong từ trường đều là $F$ = BIl sin(alpha).}
\loigiai{
Độ lớn lực từ lên đoạn dây thẳng dài l mang dòng $I$ trong từ trường đều là $F$ = BIl sin(alpha). Vì vậy phương án $D$ đúng; các phương án còn lại trái với kiến thức hoặc điều kiện áp dụng.
}
\end{ex}

% ===== Câu 147 =====
% ID: L12C3B15-32-DS
% Mức: VD
\dangbt{Xác định phương, chiều lực từ bằng quy tắc bàn tay trái}
\begin{ex}
% ID: L12C3B15-32-DS
% Mức: VD
Xét Xác định phương, chiều lực từ bằng quy tắc bàn tay trái (Tình huống 5). Hãy xác định tính đúng, sai của bốn phát biểu sau.
\choiceTF
{\True Lực từ tác dụng lên đoạn dây có phương vuông góc với cả dòng điện và vectơ cảm ứng từ; chiều xác định bằng quy tắc bàn tay trái.}
{Kết quả không phụ thuộc dữ kiện và có thể bỏ qua điều kiện.}
{Lực từ luôn cùng chiều dòng điện trong dây.}
{\True Khi giải cần xác định đúng phương, chiều, điều kiện áp dụng và đơn vị.}
\loigiai{
\begin{itemchoice}
\item (Đúng) Lực từ tác dụng lên đoạn dây có phương vuông góc với cả dòng điện và vectơ cảm ứng từ; chiều xác định bằng quy tắc bàn tay trái. (Vì): Lực từ tác dụng lên đoạn dây có phương vuông góc với cả dòng điện và vectơ cảm ứng từ; chiều xác định bằng quy tắc bàn tay trái. b) Sai: Kết quả không phụ thuộc dữ kiện và có thể bỏ qua điều kiện. c) Sai: Lực từ luôn cùng chiều dòng điện trong dây. d) Đúng: Khi giải cần xác định đúng phương, chiều, điều kiện áp dụng và đơn vị. Kiến thức cốt lõi: Lực từ tác dụng lên đoạn dây có phương vuông góc với cả dòng điện và vectơ cảm ứng từ; chiều xác định bằng quy tắc bàn tay trái
\item (Sai) Kết quả không phụ thuộc dữ kiện và có thể bỏ qua điều kiện.
\item (Sai) Lực từ luôn cùng chiều dòng điện trong dây.
\item (Đúng) Khi giải cần xác định đúng phương, chiều, điều kiện áp dụng và đơn vị.
\end{itemchoice}
}
\end{ex}

% ===== Câu 148 =====
% ID: L12C3B15-32-TL
% Mức: TH
\begin{bt}
% ID: L12C3B15-32-TL
% Mức: TH
Giải thích vì sao nhận định sau đúng: “Lực từ tác dụng lên đoạn dây có phương vuông góc với cả dòng điện và vectơ cảm ứng từ; chiều xác định bằng quy tắc bàn tay trái.”
\loigiai{
Cần nêu được: Lực từ tác dụng lên đoạn dây có phương vuông góc với cả dòng điện và vectơ cảm ứng từ; chiều xác định bằng quy tắc bàn tay trái. Khi vận dụng phải xác định đúng dữ kiện, phương chiều nếu có, điều kiện áp dụng, hệ đơn vị và kiểm tra tính hợp lí. Nhận định “Lực từ luôn cùng chiều dòng điện trong dây.” sai vì trái với kiến thức nêu trên.
}
\end{bt}

% ===== Câu 149 =====
% ID: L12C3B15-32-TLN
% Mức: TH
\begin{ex}
% ID: L12C3B15-32-TLN
% Mức: TH
Trong bốn phát biểu về Xác định phương, chiều lực từ bằng quy tắc bàn tay trái, có bao nhiêu phát biểu đúng? (Chỉ ghi một số nguyên.) (1) Lực từ tác dụng lên đoạn dây có phương vuông góc với cả dòng điện và vectơ cảm ứng từ; chiều xác định bằng quy tắc bàn tay trái. (2) Cần kiểm tra điều kiện áp dụng. (3) Lực từ luôn cùng chiều dòng điện trong dây. (4) Có thể bỏ qua đơn vị của kết quả.
\shortans{2}
\loigiai{
Đối chiếu từng phát biểu với kiến thức: Lực từ tác dụng lên đoạn dây có phương vuông góc với cả dòng điện và vectơ cảm ứng từ; chiều xác định bằng quy tắc bàn tay trái. Số phát biểu đúng là 2.
}
\end{ex}

% ===== Câu 150 =====
% ID: L12C3B15-32-TN
% Mức: NB
\dangbt{Tính lực từ tác dụng lên đoạn dây dẫn mang dòng điện}
\begin{ex}
% ID: L12C3B15-32-TN
% Mức: NB
Trong nội dung Tính lực từ tác dụng lên đoạn dây dẫn mang dòng điện?
\choice
{\True Độ lớn lực từ lên đoạn dây thẳng dài l mang dòng $I$ trong từ trường đều là $F$ = BIl sin(alpha).}
{Lực từ được tính bằng $F$ = BI/l và không phụ thuộc góc.}
{Có thể bỏ qua phương, chiều và đơn vị của các đại lượng trong mọi trường hợp.}
{Kết quả không phụ thuộc dữ kiện và điều kiện áp dụng.}
\loigiai{
Độ lớn lực từ lên đoạn dây thẳng dài l mang dòng $I$ trong từ trường đều là $F$ = BIl sin(alpha). Vì vậy phương án $A$ đúng; các phương án còn lại trái với kiến thức hoặc điều kiện áp dụng.
}
\end{ex}

% ===== Câu 151 =====
% ID: L12C3B15-33-DS
% Mức: TH
\dangbt{Nhận biết lực từ và đặc trưng của vectơ cảm ứng từ}
\begin{ex}
% ID: L12C3B15-33-DS
% Mức: TH
Xét Nhận biết lực từ và đặc trưng của vectơ cảm ứng từ (Tình huống 1). Hãy xác định tính đúng, sai của bốn phát biểu sau.
\choiceTF
{\True Vectơ cảm ứng từ tại một điểm có hướng trùng với hướng của từ trường tại điểm đó; đơn vị là tesla (T).}
{Cảm ứng từ là đại lượng vô hướng và có đơn vị niutơn.}
{\True Khi giải cần xác định đúng phương, chiều, điều kiện áp dụng và đơn vị.}
{Kết quả không phụ thuộc dữ kiện và có thể bỏ qua điều kiện.}
\loigiai{
\begin{itemchoice}
\item (Đúng) Vectơ cảm ứng từ tại một điểm có hướng trùng với hướng của từ trường tại điểm đó; đơn vị là tesla (T). (Vì): Vectơ cảm ứng từ tại một điểm có hướng trùng với hướng của từ trường tại điểm đó; đơn vị là tesla (T). b) Sai: Cảm ứng từ là đại lượng vô hướng và có đơn vị niutơn. c) Đúng: Khi giải cần xác định đúng phương, chiều, điều kiện áp dụng và đơn vị. d) Sai: Kết quả không phụ thuộc dữ kiện và có thể bỏ qua điều kiện. Kiến thức cốt lõi: Vectơ cảm ứng từ tại một điểm có hướng trùng với hướng của từ trường tại điểm đó; đơn vị là tesla (T
\item (Sai) Cảm ứng từ là đại lượng vô hướng và có đơn vị niutơn.
\item (Đúng) Khi giải cần xác định đúng phương, chiều, điều kiện áp dụng và đơn vị.
\item (Sai) Kết quả không phụ thuộc dữ kiện và có thể bỏ qua điều kiện.
\end{itemchoice}
}
\end{ex}

% ===== Câu 152 =====
% ID: L12C3B15-33-TL
% Mức: VD
\dangbt{Xác định phương, chiều lực từ bằng quy tắc bàn tay trái}
\begin{bt}
% ID: L12C3B15-33-TL
% Mức: VD
Phân tích sai lầm trong nhận định: “Lực từ luôn cùng chiều dòng điện trong dây.”
\loigiai{
Cần nêu được: Lực từ tác dụng lên đoạn dây có phương vuông góc với cả dòng điện và vectơ cảm ứng từ; chiều xác định bằng quy tắc bàn tay trái. Khi vận dụng phải xác định đúng dữ kiện, phương chiều nếu có, điều kiện áp dụng, hệ đơn vị và kiểm tra tính hợp lí. Nhận định “Lực từ luôn cùng chiều dòng điện trong dây.” sai vì trái với kiến thức nêu trên.
}
\end{bt}

% ===== Câu 153 =====
% ID: L12C3B15-33-TLN
% Mức: TH
\begin{ex}
% ID: L12C3B15-33-TLN
% Mức: TH
Trong bốn phát biểu về Xác định phương, chiều lực từ bằng quy tắc bàn tay trái, có bao nhiêu phát biểu đúng? (Chỉ ghi một số nguyên.) (1) Lực từ luôn cùng chiều dòng điện trong dây. (2) Lực từ tác dụng lên đoạn dây có phương vuông góc với cả dòng điện và vectơ cảm ứng từ; chiều xác định bằng quy tắc bàn tay trái. (3) Lực từ luôn cùng chiều dòng điện trong dây. (4) Có thể bỏ qua đơn vị của kết quả.
\shortans{1}
\loigiai{
Đối chiếu từng phát biểu với kiến thức: Lực từ tác dụng lên đoạn dây có phương vuông góc với cả dòng điện và vectơ cảm ứng từ; chiều xác định bằng quy tắc bàn tay trái. Số phát biểu đúng là 1.
}
\end{ex}

% ===== Câu 154 =====
% ID: L12C3B15-33-TN
% Mức: NB
\dangbt{Tính lực từ tác dụng lên đoạn dây dẫn mang dòng điện}
\begin{ex}
% ID: L12C3B15-33-TN
% Mức: NB
Tình huống 3: chọn kết luận chính xác về Tính lực từ tác dụng lên đoạn dây dẫn mang dòng điện?
\choice
{Kết quả không phụ thuộc dữ kiện và điều kiện áp dụng.}
{\True Độ lớn lực từ lên đoạn dây thẳng dài l mang dòng $I$ trong từ trường đều là $F$ = BIl sin(alpha).}
{Lực từ được tính bằng $F$ = BI/l và không phụ thuộc góc.}
{Có thể bỏ qua phương, chiều và đơn vị của các đại lượng trong mọi trường hợp.}
\loigiai{
Độ lớn lực từ lên đoạn dây thẳng dài l mang dòng $I$ trong từ trường đều là $F$ = BIl sin(alpha). Vì vậy phương án $B$ đúng; các phương án còn lại trái với kiến thức hoặc điều kiện áp dụng.
}
\end{ex}

% ===== Câu 155 =====
% ID: L12C3B15-34-DS
% Mức: TH
\dangbt{Nhận biết lực từ và đặc trưng của vectơ cảm ứng từ}
\begin{ex}
% ID: L12C3B15-34-DS
% Mức: TH
Xét Nhận biết lực từ và đặc trưng của vectơ cảm ứng từ (Tình huống 2). Hãy xác định tính đúng, sai của bốn phát biểu sau.
\choiceTF
{Cảm ứng từ là đại lượng vô hướng và có đơn vị niutơn.}
{\True Vectơ cảm ứng từ tại một điểm có hướng trùng với hướng của từ trường tại điểm đó; đơn vị là tesla (T).}
{Kết quả không phụ thuộc dữ kiện và có thể bỏ qua điều kiện.}
{\True Khi giải cần xác định đúng phương, chiều, điều kiện áp dụng và đơn vị.}
\loigiai{
\begin{itemchoice}
\item (Sai) Cảm ứng từ là đại lượng vô hướng và có đơn vị niutơn. (Vì): Cảm ứng từ là đại lượng vô hướng và có đơn vị niutơn. b) Đúng: Vectơ cảm ứng từ tại một điểm có hướng trùng với hướng của từ trường tại điểm đó; đơn vị là tesla (T). c) Sai: Kết quả không phụ thuộc dữ kiện và có thể bỏ qua điều kiện. d) Đúng: Khi giải cần xác định đúng phương, chiều, điều kiện áp dụng và đơn vị. Kiến thức cốt lõi: Vectơ cảm ứng từ tại một điểm có hướng trùng với hướng của từ trường tại điểm đó; đơn vị là tesla (T
\item (Đúng) Vectơ cảm ứng từ tại một điểm có hướng trùng với hướng của từ trường tại điểm đó; đơn vị là tesla (T).
\item (Sai) Kết quả không phụ thuộc dữ kiện và có thể bỏ qua điều kiện.
\item (Đúng) Khi giải cần xác định đúng phương, chiều, điều kiện áp dụng và đơn vị.
\end{itemchoice}
}
\end{ex}

% ===== Câu 156 =====
% ID: L12C3B15-34-TL
% Mức: VD
\dangbt{Xác định phương, chiều lực từ bằng quy tắc bàn tay trái}
\begin{bt}
% ID: L12C3B15-34-TL
% Mức: VD
Đề xuất các bước giải một bài toán thuộc dạng Xác định phương, chiều lực từ bằng quy tắc bàn tay trái.
\loigiai{
Cần nêu được: Lực từ tác dụng lên đoạn dây có phương vuông góc với cả dòng điện và vectơ cảm ứng từ; chiều xác định bằng quy tắc bàn tay trái. Khi vận dụng phải xác định đúng dữ kiện, phương chiều nếu có, điều kiện áp dụng, hệ đơn vị và kiểm tra tính hợp lí. Nhận định “Lực từ luôn cùng chiều dòng điện trong dây.” sai vì trái với kiến thức nêu trên.
}
\end{bt}

% ===== Câu 157 =====
% ID: L12C3B15-34-TLN
% Mức: TH
\begin{ex}
% ID: L12C3B15-34-TLN
% Mức: TH
Trong bốn phát biểu về Xác định phương, chiều lực từ bằng quy tắc bàn tay trái, có bao nhiêu phát biểu đúng? (Chỉ ghi một số nguyên.) (1) Lực từ tác dụng lên đoạn dây có phương vuông góc với cả dòng điện và vectơ cảm ứng từ; chiều xác định bằng quy tắc bàn tay trái. (2) Cần kiểm tra điều kiện áp dụng. (3) Lực từ luôn cùng chiều dòng điện trong dây. (4) Có thể bỏ qua đơn vị của kết quả.
\shortans{2}
\loigiai{
Đối chiếu từng phát biểu với kiến thức: Lực từ tác dụng lên đoạn dây có phương vuông góc với cả dòng điện và vectơ cảm ứng từ; chiều xác định bằng quy tắc bàn tay trái. Số phát biểu đúng là 2.
}
\end{ex}

% ===== Câu 158 =====
% ID: L12C3B15-34-TN
% Mức: TH
\dangbt{Tính lực từ tác dụng lên đoạn dây dẫn mang dòng điện}
\begin{ex}
% ID: L12C3B15-34-TN
% Mức: TH
Nội dung nào mô tả đúng nhất Tính lực từ tác dụng lên đoạn dây dẫn mang dòng điện?
\choice
{Có thể bỏ qua phương, chiều và đơn vị của các đại lượng trong mọi trường hợp.}
{Kết quả không phụ thuộc dữ kiện và điều kiện áp dụng.}
{\True Độ lớn lực từ lên đoạn dây thẳng dài l mang dòng $I$ trong từ trường đều là $F$ = BIl sin(alpha).}
{Lực từ được tính bằng $F$ = BI/l và không phụ thuộc góc.}
\loigiai{
Độ lớn lực từ lên đoạn dây thẳng dài l mang dòng $I$ trong từ trường đều là $F$ = BIl sin(alpha). Vì vậy phương án $C$ đúng; các phương án còn lại trái với kiến thức hoặc điều kiện áp dụng.
}
\end{ex}

% ===== Câu 159 =====
% ID: L12C3B15-35-DS
% Mức: VD
\dangbt{Nhận biết lực từ và đặc trưng của vectơ cảm ứng từ}
\begin{ex}
% ID: L12C3B15-35-DS
% Mức: VD
Xét Nhận biết lực từ và đặc trưng của vectơ cảm ứng từ (Tình huống 3). Hãy xác định tính đúng, sai của bốn phát biểu sau.
\choiceTF
{\True Khi giải cần xác định đúng phương, chiều, điều kiện áp dụng và đơn vị.}
{Kết quả không phụ thuộc dữ kiện và có thể bỏ qua điều kiện.}
{\True Vectơ cảm ứng từ tại một điểm có hướng trùng với hướng của từ trường tại điểm đó; đơn vị là tesla (T).}
{Cảm ứng từ là đại lượng vô hướng và có đơn vị niutơn.}
\loigiai{
\begin{itemchoice}
\item (Đúng) Khi giải cần xác định đúng phương, chiều, điều kiện áp dụng và đơn vị. (Vì): Khi giải cần xác định đúng phương, chiều, điều kiện áp dụng và đơn vị. b) Sai: Kết quả không phụ thuộc dữ kiện và có thể bỏ qua điều kiện. c) Đúng: Vectơ cảm ứng từ tại một điểm có hướng trùng với hướng của từ trường tại điểm đó; đơn vị là tesla (T). d) Sai: Cảm ứng từ là đại lượng vô hướng và có đơn vị niutơn. Kiến thức cốt lõi: Vectơ cảm ứng từ tại một điểm có hướng trùng với hướng của từ trường tại điểm đó; đơn vị là tesla (T
\item (Sai) Kết quả không phụ thuộc dữ kiện và có thể bỏ qua điều kiện.
\item (Đúng) Vectơ cảm ứng từ tại một điểm có hướng trùng với hướng của từ trường tại điểm đó; đơn vị là tesla (T).
\item (Sai) Cảm ứng từ là đại lượng vô hướng và có đơn vị niutơn.
\end{itemchoice}
}
\end{ex}

% ===== Câu 160 =====
% ID: L12C3B15-35-TL
% Mức: VD
\dangbt{Xác định phương, chiều lực từ bằng quy tắc bàn tay trái}
\begin{bt}
% ID: L12C3B15-35-TL
% Mức: VD
Nêu một tình huống thực tiễn liên quan đến Xác định phương, chiều lực từ bằng quy tắc bàn tay trái và giải thích bằng kiến thức Vật lí.
\loigiai{
Cần nêu được: Lực từ tác dụng lên đoạn dây có phương vuông góc với cả dòng điện và vectơ cảm ứng từ; chiều xác định bằng quy tắc bàn tay trái. Khi vận dụng phải xác định đúng dữ kiện, phương chiều nếu có, điều kiện áp dụng, hệ đơn vị và kiểm tra tính hợp lí. Nhận định “Lực từ luôn cùng chiều dòng điện trong dây.” sai vì trái với kiến thức nêu trên.
}
\end{bt}

% ===== Câu 161 =====
% ID: L12C3B15-35-TLN
% Mức: VD
\begin{ex}
% ID: L12C3B15-35-TLN
% Mức: VD
Trong bốn phát biểu về Xác định phương, chiều lực từ bằng quy tắc bàn tay trái, có bao nhiêu phát biểu đúng? (Chỉ ghi một số nguyên.) (1) Lực từ luôn cùng chiều dòng điện trong dây. (2) Lực từ tác dụng lên đoạn dây có phương vuông góc với cả dòng điện và vectơ cảm ứng từ; chiều xác định bằng quy tắc bàn tay trái. (3) Lực từ luôn cùng chiều dòng điện trong dây. (4) Có thể bỏ qua đơn vị của kết quả.
\shortans{2}
\loigiai{
Đối chiếu từng phát biểu với kiến thức: Lực từ tác dụng lên đoạn dây có phương vuông góc với cả dòng điện và vectơ cảm ứng từ; chiều xác định bằng quy tắc bàn tay trái. Số phát biểu đúng là 2.
}
\end{ex}

% ===== Câu 162 =====
% ID: L12C3B15-35-TN
% Mức: TH
\dangbt{Tính lực từ tác dụng lên đoạn dây dẫn mang dòng điện}
\begin{ex}
% ID: L12C3B15-35-TN
% Mức: TH
Khi phân tích Tính lực từ tác dụng lên đoạn dây dẫn mang dòng điện?
\choice
{Lực từ được tính bằng $F$ = BI/l và không phụ thuộc góc.}
{Có thể bỏ qua phương, chiều và đơn vị của các đại lượng trong mọi trường hợp.}
{Kết quả không phụ thuộc dữ kiện và điều kiện áp dụng.}
{\True Độ lớn lực từ lên đoạn dây thẳng dài l mang dòng $I$ trong từ trường đều là $F$ = BIl sin(alpha).}
\loigiai{
Độ lớn lực từ lên đoạn dây thẳng dài l mang dòng $I$ trong từ trường đều là $F$ = BIl sin(alpha). Vì vậy phương án $D$ đúng; các phương án còn lại trái với kiến thức hoặc điều kiện áp dụng.
}
\end{ex}

% ===== Câu 163 =====
% ID: L12C3B15-36-DS
% Mức: VD
\dangbt{Nhận biết lực từ và đặc trưng của vectơ cảm ứng từ}
\begin{ex}
% ID: L12C3B15-36-DS
% Mức: VD
Xét Nhận biết lực từ và đặc trưng của vectơ cảm ứng từ (Tình huống 4). Hãy xác định tính đúng, sai của bốn phát biểu sau.
\choiceTF
{Kết quả không phụ thuộc dữ kiện và có thể bỏ qua điều kiện.}
{\True Khi giải cần xác định đúng phương, chiều, điều kiện áp dụng và đơn vị.}
{Cảm ứng từ là đại lượng vô hướng và có đơn vị niutơn.}
{\True Vectơ cảm ứng từ tại một điểm có hướng trùng với hướng của từ trường tại điểm đó; đơn vị là tesla (T).}
\loigiai{
\begin{itemchoice}
\item (Sai) Kết quả không phụ thuộc dữ kiện và có thể bỏ qua điều kiện. (Vì): Kết quả không phụ thuộc dữ kiện và có thể bỏ qua điều kiện. b) Đúng: Khi giải cần xác định đúng phương, chiều, điều kiện áp dụng và đơn vị. c) Sai: Cảm ứng từ là đại lượng vô hướng và có đơn vị niutơn. d) Đúng: Vectơ cảm ứng từ tại một điểm có hướng trùng với hướng của từ trường tại điểm đó; đơn vị là tesla (T). Kiến thức cốt lõi: Vectơ cảm ứng từ tại một điểm có hướng trùng với hướng của từ trường tại điểm đó; đơn vị là tesla (T
\item (Đúng) Khi giải cần xác định đúng phương, chiều, điều kiện áp dụng và đơn vị.
\item (Sai) Cảm ứng từ là đại lượng vô hướng và có đơn vị niutơn.
\item (Đúng) Vectơ cảm ứng từ tại một điểm có hướng trùng với hướng của từ trường tại điểm đó; đơn vị là tesla (T).
\end{itemchoice}
}
\end{ex}

% ===== Câu 164 =====
% ID: L12C3B15-36-TL
% Mức: VDC
\dangbt{Xác định phương, chiều lực từ bằng quy tắc bàn tay trái}
\begin{bt}
% ID: L12C3B15-36-TL
% Mức: VDC
So sánh nhận định đúng “Lực từ tác dụng lên đoạn dây có phương vuông góc với cả dòng điện và vectơ cảm ứng từ; chiều xác định bằng quy tắc bàn tay trái.” với nhận định sai “Lực từ luôn cùng chiều dòng điện trong dây.”; chỉ rõ căn cứ Vật lí.
\loigiai{
Cần nêu được: Lực từ tác dụng lên đoạn dây có phương vuông góc với cả dòng điện và vectơ cảm ứng từ; chiều xác định bằng quy tắc bàn tay trái. Khi vận dụng phải xác định đúng dữ kiện, phương chiều nếu có, điều kiện áp dụng, hệ đơn vị và kiểm tra tính hợp lí. Nhận định “Lực từ luôn cùng chiều dòng điện trong dây.” sai vì trái với kiến thức nêu trên.
}
\end{bt}

% ===== Câu 165 =====
% ID: L12C3B15-36-TLN
% Mức: VD
\begin{ex}
% ID: L12C3B15-36-TLN
% Mức: VD
Trong bốn phát biểu về Xác định phương, chiều lực từ bằng quy tắc bàn tay trái, có bao nhiêu phát biểu đúng? (Chỉ ghi một số nguyên.) (1) Lực từ tác dụng lên đoạn dây có phương vuông góc với cả dòng điện và vectơ cảm ứng từ; chiều xác định bằng quy tắc bàn tay trái. (2) Cần kiểm tra điều kiện áp dụng. (3) Lực từ luôn cùng chiều dòng điện trong dây. (4) Có thể bỏ qua đơn vị của kết quả.
\shortans{2}
\loigiai{
Đối chiếu từng phát biểu với kiến thức: Lực từ tác dụng lên đoạn dây có phương vuông góc với cả dòng điện và vectơ cảm ứng từ; chiều xác định bằng quy tắc bàn tay trái. Số phát biểu đúng là 2.
}
\end{ex}

% ===== Câu 166 =====
% ID: L12C3B15-36-TN
% Mức: TH
\dangbt{Tính lực từ tác dụng lên đoạn dây dẫn mang dòng điện}
\begin{ex}
% ID: L12C3B15-36-TN
% Mức: TH
Một học sinh ôn tập Tính lực từ tác dụng lên đoạn dây dẫn mang dòng điện?
\choice
{\True Độ lớn lực từ lên đoạn dây thẳng dài l mang dòng $I$ trong từ trường đều là $F$ = BIl sin(alpha).}
{Lực từ được tính bằng $F$ = BI/l và không phụ thuộc góc.}
{Có thể bỏ qua phương, chiều và đơn vị của các đại lượng trong mọi trường hợp.}
{Kết quả không phụ thuộc dữ kiện và điều kiện áp dụng.}
\loigiai{
Độ lớn lực từ lên đoạn dây thẳng dài l mang dòng $I$ trong từ trường đều là $F$ = BIl sin(alpha). Vì vậy phương án $A$ đúng; các phương án còn lại trái với kiến thức hoặc điều kiện áp dụng.
}
\end{ex}

% ===== Câu 167 =====
% ID: L12C3B15-37-DS
% Mức: VD
\dangbt{Nhận biết lực từ và đặc trưng của vectơ cảm ứng từ}
\begin{ex}
% ID: L12C3B15-37-DS
% Mức: VD
Xét Nhận biết lực từ và đặc trưng của vectơ cảm ứng từ (Tình huống 5). Hãy xác định tính đúng, sai của bốn phát biểu sau.
\choiceTF
{\True Vectơ cảm ứng từ tại một điểm có hướng trùng với hướng của từ trường tại điểm đó; đơn vị là tesla (T).}
{Kết quả không phụ thuộc dữ kiện và có thể bỏ qua điều kiện.}
{Cảm ứng từ là đại lượng vô hướng và có đơn vị niutơn.}
{\True Khi giải cần xác định đúng phương, chiều, điều kiện áp dụng và đơn vị.}
\loigiai{
\begin{itemchoice}
\item (Đúng) Vectơ cảm ứng từ tại một điểm có hướng trùng với hướng của từ trường tại điểm đó; đơn vị là tesla (T). (Vì): Vectơ cảm ứng từ tại một điểm có hướng trùng với hướng của từ trường tại điểm đó; đơn vị là tesla (T). b) Sai: Kết quả không phụ thuộc dữ kiện và có thể bỏ qua điều kiện. c) Sai: Cảm ứng từ là đại lượng vô hướng và có đơn vị niutơn. d) Đúng: Khi giải cần xác định đúng phương, chiều, điều kiện áp dụng và đơn vị. Kiến thức cốt lõi: Vectơ cảm ứng từ tại một điểm có hướng trùng với hướng của từ trường tại điểm đó; đơn vị là tesla (T
\item (Sai) Kết quả không phụ thuộc dữ kiện và có thể bỏ qua điều kiện.
\item (Sai) Cảm ứng từ là đại lượng vô hướng và có đơn vị niutơn.
\item (Đúng) Khi giải cần xác định đúng phương, chiều, điều kiện áp dụng và đơn vị.
\end{itemchoice}
}
\end{ex}

% ===== Câu 168 =====
% ID: L12C3B15-37-TL
% Mức: TH
\begin{bt}
% ID: L12C3B15-37-TL
% Mức: TH
Trình bày kiến thức cốt lõi của Nhận biết lực từ và đặc trưng của vectơ cảm ứng từ và nêu điều kiện áp dụng.
\loigiai{
Cần nêu được: Vectơ cảm ứng từ tại một điểm có hướng trùng với hướng của từ trường tại điểm đó; đơn vị là tesla (T). Khi vận dụng phải xác định đúng dữ kiện, phương chiều nếu có, điều kiện áp dụng, hệ đơn vị và kiểm tra tính hợp lí. Nhận định “Cảm ứng từ là đại lượng vô hướng và có đơn vị niutơn.” sai vì trái với kiến thức nêu trên.
}
\end{bt}

% ===== Câu 169 =====
% ID: L12C3B15-37-TLN
% Mức: VDC
\dangbt{Xác định phương, chiều lực từ bằng quy tắc bàn tay trái}
\begin{ex}
% ID: L12C3B15-37-TLN
% Mức: VDC
Trong bốn phát biểu về Xác định phương, chiều lực từ bằng quy tắc bàn tay trái, có bao nhiêu phát biểu đúng? (Chỉ ghi một số nguyên.) (1) Lực từ luôn cùng chiều dòng điện trong dây. (2) Lực từ tác dụng lên đoạn dây có phương vuông góc với cả dòng điện và vectơ cảm ứng từ; chiều xác định bằng quy tắc bàn tay trái. (3) Phải dùng đúng định luật cho tình huống. (4) Lực từ tác dụng lên đoạn dây có phương vuông góc với cả dòng điện và vectơ cảm ứng từ; chiều xác định bằng quy tắc bàn tay trái.
\shortans{3}
\loigiai{
Đối chiếu từng phát biểu với kiến thức: Lực từ tác dụng lên đoạn dây có phương vuông góc với cả dòng điện và vectơ cảm ứng từ; chiều xác định bằng quy tắc bàn tay trái. Số phát biểu đúng là 3.
}
\end{ex}

% ===== Câu 170 =====
% ID: L12C3B15-37-TN
% Mức: TH
\dangbt{Tính lực từ tác dụng lên đoạn dây dẫn mang dòng điện}
\begin{ex}
% ID: L12C3B15-37-TN
% Mức: TH
Chọn phát biểu không trái với kiến thức về Tính lực từ tác dụng lên đoạn dây dẫn mang dòng điện?
\choice
{Kết quả không phụ thuộc dữ kiện và điều kiện áp dụng.}
{\True Độ lớn lực từ lên đoạn dây thẳng dài l mang dòng $I$ trong từ trường đều là $F$ = BIl sin(alpha).}
{Lực từ được tính bằng $F$ = BI/l và không phụ thuộc góc.}
{Có thể bỏ qua phương, chiều và đơn vị của các đại lượng trong mọi trường hợp.}
\loigiai{
Độ lớn lực từ lên đoạn dây thẳng dài l mang dòng $I$ trong từ trường đều là $F$ = BIl sin(alpha). Vì vậy phương án $B$ đúng; các phương án còn lại trái với kiến thức hoặc điều kiện áp dụng.
}
\end{ex}

% ===== Câu 171 =====
% ID: L12C3B15-38-DS
% Mức: TH
\dangbt{Xác định phương, chiều lực từ bằng quy tắc bàn tay trái}
\begin{ex}
% ID: L12C3B15-38-DS
% Mức: TH
Xét Xác định phương, chiều lực từ bằng quy tắc bàn tay trái (Tình huống 1). Hãy xác định tính đúng, sai của bốn phát biểu sau.
\choiceTF
{\True Lực từ tác dụng lên đoạn dây có phương vuông góc với cả dòng điện và vectơ cảm ứng từ; chiều xác định bằng quy tắc bàn tay trái.}
{Lực từ luôn cùng chiều dòng điện trong dây.}
{\True Khi giải cần xác định đúng phương, chiều, điều kiện áp dụng và đơn vị.}
{Kết quả không phụ thuộc dữ kiện và có thể bỏ qua điều kiện.}
\loigiai{
\begin{itemchoice}
\item (Đúng) Lực từ tác dụng lên đoạn dây có phương vuông góc với cả dòng điện và vectơ cảm ứng từ; chiều xác định bằng quy tắc bàn tay trái. (Vì): Lực từ tác dụng lên đoạn dây có phương vuông góc với cả dòng điện và vectơ cảm ứng từ; chiều xác định bằng quy tắc bàn tay trái. b) Sai: Lực từ luôn cùng chiều dòng điện trong dây. c) Đúng: Khi giải cần xác định đúng phương, chiều, điều kiện áp dụng và đơn vị. d) Sai: Kết quả không phụ thuộc dữ kiện và có thể bỏ qua điều kiện. Kiến thức cốt lõi: Lực từ tác dụng lên đoạn dây có phương vuông góc với cả dòng điện và vectơ cảm ứng từ; chiều xác định bằng quy tắc bàn tay trái
\item (Sai) Lực từ luôn cùng chiều dòng điện trong dây.
\item (Đúng) Khi giải cần xác định đúng phương, chiều, điều kiện áp dụng và đơn vị.
\item (Sai) Kết quả không phụ thuộc dữ kiện và có thể bỏ qua điều kiện.
\end{itemchoice}
}
\end{ex}

% ===== Câu 172 =====
% ID: L12C3B15-38-TL
% Mức: TH
\dangbt{Nhận biết lực từ và đặc trưng của vectơ cảm ứng từ}
\begin{bt}
% ID: L12C3B15-38-TL
% Mức: TH
Giải thích vì sao nhận định sau đúng: “Vectơ cảm ứng từ tại một điểm có hướng trùng với hướng của từ trường tại điểm đó; đơn vị là tesla (T).”
\loigiai{
Cần nêu được: Vectơ cảm ứng từ tại một điểm có hướng trùng với hướng của từ trường tại điểm đó; đơn vị là tesla (T). Khi vận dụng phải xác định đúng dữ kiện, phương chiều nếu có, điều kiện áp dụng, hệ đơn vị và kiểm tra tính hợp lí. Nhận định “Cảm ứng từ là đại lượng vô hướng và có đơn vị niutơn.” sai vì trái với kiến thức nêu trên.
}
\end{bt}

% ===== Câu 173 =====
% ID: L12C3B15-38-TLN
% Mức: TH
\begin{ex}
% ID: L12C3B15-38-TLN
% Mức: TH
Trong bốn phát biểu về Nhận biết lực từ và đặc trưng của vectơ cảm ứng từ, có bao nhiêu phát biểu đúng? (Chỉ ghi một số nguyên.) (1) Vectơ cảm ứng từ tại một điểm có hướng trùng với hướng của từ trường tại điểm đó; đơn vị là tesla (T). (2) Cần kiểm tra điều kiện áp dụng. (3) Cảm ứng từ là đại lượng vô hướng và có đơn vị niutơn. (4) Có thể bỏ qua đơn vị của kết quả.
\shortans{2}
\loigiai{
Đối chiếu từng phát biểu với kiến thức: Vectơ cảm ứng từ tại một điểm có hướng trùng với hướng của từ trường tại điểm đó; đơn vị là tesla (T). Số phát biểu đúng là 2.
}
\end{ex}

% ===== Câu 174 =====
% ID: L12C3B15-38-TN
% Mức: VD
\dangbt{Tính lực từ tác dụng lên đoạn dây dẫn mang dòng điện}
\begin{ex}
% ID: L12C3B15-38-TN
% Mức: VD
Cơ sở Vật lí đúng dùng để giải Tính lực từ tác dụng lên đoạn dây dẫn mang dòng điện?
\choice
{Có thể bỏ qua phương, chiều và đơn vị của các đại lượng trong mọi trường hợp.}
{Kết quả không phụ thuộc dữ kiện và điều kiện áp dụng.}
{\True Độ lớn lực từ lên đoạn dây thẳng dài l mang dòng $I$ trong từ trường đều là $F$ = BIl sin(alpha).}
{Lực từ được tính bằng $F$ = BI/l và không phụ thuộc góc.}
\loigiai{
Độ lớn lực từ lên đoạn dây thẳng dài l mang dòng $I$ trong từ trường đều là $F$ = BIl sin(alpha). Vì vậy phương án $C$ đúng; các phương án còn lại trái với kiến thức hoặc điều kiện áp dụng.
}
\end{ex}

% ===== Câu 175 =====
% ID: L12C3B15-39-DS
% Mức: TH
\dangbt{Xác định phương, chiều lực từ bằng quy tắc bàn tay trái}
\begin{ex}
% ID: L12C3B15-39-DS
% Mức: TH
Xét Xác định phương, chiều lực từ bằng quy tắc bàn tay trái (Tình huống 2). Hãy xác định tính đúng, sai của bốn phát biểu sau.
\choiceTF
{Lực từ luôn cùng chiều dòng điện trong dây.}
{\True Lực từ tác dụng lên đoạn dây có phương vuông góc với cả dòng điện và vectơ cảm ứng từ; chiều xác định bằng quy tắc bàn tay trái.}
{Kết quả không phụ thuộc dữ kiện và có thể bỏ qua điều kiện.}
{\True Khi giải cần xác định đúng phương, chiều, điều kiện áp dụng và đơn vị.}
\loigiai{
\begin{itemchoice}
\item (Sai) Lực từ luôn cùng chiều dòng điện trong dây. (Vì): Lực từ luôn cùng chiều dòng điện trong dây. b) Đúng: Lực từ tác dụng lên đoạn dây có phương vuông góc với cả dòng điện và vectơ cảm ứng từ; chiều xác định bằng quy tắc bàn tay trái. c) Sai: Kết quả không phụ thuộc dữ kiện và có thể bỏ qua điều kiện. d) Đúng: Khi giải cần xác định đúng phương, chiều, điều kiện áp dụng và đơn vị. Kiến thức cốt lõi: Lực từ tác dụng lên đoạn dây có phương vuông góc với cả dòng điện và vectơ cảm ứng từ; chiều xác định bằng quy tắc bàn tay trái
\item (Đúng) Lực từ tác dụng lên đoạn dây có phương vuông góc với cả dòng điện và vectơ cảm ứng từ; chiều xác định bằng quy tắc bàn tay trái.
\item (Sai) Kết quả không phụ thuộc dữ kiện và có thể bỏ qua điều kiện.
\item (Đúng) Khi giải cần xác định đúng phương, chiều, điều kiện áp dụng và đơn vị.
\end{itemchoice}
}
\end{ex}

% ===== Câu 176 =====
% ID: L12C3B15-39-TL
% Mức: VD
\dangbt{Nhận biết lực từ và đặc trưng của vectơ cảm ứng từ}
\begin{bt}
% ID: L12C3B15-39-TL
% Mức: VD
Phân tích sai lầm trong nhận định: “Cảm ứng từ là đại lượng vô hướng và có đơn vị niutơn.”
\loigiai{
Cần nêu được: Vectơ cảm ứng từ tại một điểm có hướng trùng với hướng của từ trường tại điểm đó; đơn vị là tesla (T). Khi vận dụng phải xác định đúng dữ kiện, phương chiều nếu có, điều kiện áp dụng, hệ đơn vị và kiểm tra tính hợp lí. Nhận định “Cảm ứng từ là đại lượng vô hướng và có đơn vị niutơn.” sai vì trái với kiến thức nêu trên.
}
\end{bt}

% ===== Câu 177 =====
% ID: L12C3B15-39-TLN
% Mức: TH
\begin{ex}
% ID: L12C3B15-39-TLN
% Mức: TH
Trong bốn phát biểu về Nhận biết lực từ và đặc trưng của vectơ cảm ứng từ, có bao nhiêu phát biểu đúng? (Chỉ ghi một số nguyên.) (1) Cảm ứng từ là đại lượng vô hướng và có đơn vị niutơn. (2) Vectơ cảm ứng từ tại một điểm có hướng trùng với hướng của từ trường tại điểm đó; đơn vị là tesla (T). (3) Cảm ứng từ là đại lượng vô hướng và có đơn vị niutơn. (4) Có thể bỏ qua đơn vị của kết quả.
\shortans{1}
\loigiai{
Đối chiếu từng phát biểu với kiến thức: Vectơ cảm ứng từ tại một điểm có hướng trùng với hướng của từ trường tại điểm đó; đơn vị là tesla (T). Số phát biểu đúng là 1.
}
\end{ex}

% ===== Câu 178 =====
% ID: L12C3B15-39-TN
% Mức: VD
\dangbt{Tính lực từ tác dụng lên đoạn dây dẫn mang dòng điện}
\begin{ex}
% ID: L12C3B15-39-TN
% Mức: VD
Trong một bài thực hành về Tính lực từ tác dụng lên đoạn dây dẫn mang dòng điện?
\choice
{Lực từ được tính bằng $F$ = BI/l và không phụ thuộc góc.}
{Có thể bỏ qua phương, chiều và đơn vị của các đại lượng trong mọi trường hợp.}
{Kết quả không phụ thuộc dữ kiện và điều kiện áp dụng.}
{\True Độ lớn lực từ lên đoạn dây thẳng dài l mang dòng $I$ trong từ trường đều là $F$ = BIl sin(alpha).}
\loigiai{
Độ lớn lực từ lên đoạn dây thẳng dài l mang dòng $I$ trong từ trường đều là $F$ = BIl sin(alpha). Vì vậy phương án $D$ đúng; các phương án còn lại trái với kiến thức hoặc điều kiện áp dụng.
}
\end{ex}

% ===== Câu 179 =====
% ID: L12C3B15-40-DS
% Mức: VD
\dangbt{Xác định phương, chiều lực từ bằng quy tắc bàn tay trái}
\begin{ex}
% ID: L12C3B15-40-DS
% Mức: VD
Xét Xác định phương, chiều lực từ bằng quy tắc bàn tay trái (Tình huống 3). Hãy xác định tính đúng, sai của bốn phát biểu sau.
\choiceTF
{\True Khi giải cần xác định đúng phương, chiều, điều kiện áp dụng và đơn vị.}
{Kết quả không phụ thuộc dữ kiện và có thể bỏ qua điều kiện.}
{\True Lực từ tác dụng lên đoạn dây có phương vuông góc với cả dòng điện và vectơ cảm ứng từ; chiều xác định bằng quy tắc bàn tay trái.}
{Lực từ luôn cùng chiều dòng điện trong dây.}
\loigiai{
\begin{itemchoice}
\item (Đúng) Khi giải cần xác định đúng phương, chiều, điều kiện áp dụng và đơn vị. (Vì): Khi giải cần xác định đúng phương, chiều, điều kiện áp dụng và đơn vị. b) Sai: Kết quả không phụ thuộc dữ kiện và có thể bỏ qua điều kiện. c) Đúng: Lực từ tác dụng lên đoạn dây có phương vuông góc với cả dòng điện và vectơ cảm ứng từ; chiều xác định bằng quy tắc bàn tay trái. d) Sai: Lực từ luôn cùng chiều dòng điện trong dây. Kiến thức cốt lõi: Lực từ tác dụng lên đoạn dây có phương vuông góc với cả dòng điện và vectơ cảm ứng từ; chiều xác định bằng quy tắc bàn tay trái
\item (Sai) Kết quả không phụ thuộc dữ kiện và có thể bỏ qua điều kiện.
\item (Đúng) Lực từ tác dụng lên đoạn dây có phương vuông góc với cả dòng điện và vectơ cảm ứng từ; chiều xác định bằng quy tắc bàn tay trái.
\item (Sai) Lực từ luôn cùng chiều dòng điện trong dây.
\end{itemchoice}
}
\end{ex}

% ===== Câu 180 =====
% ID: L12C3B15-40-TL
% Mức: VD
\dangbt{Nhận biết lực từ và đặc trưng của vectơ cảm ứng từ}
\begin{bt}
% ID: L12C3B15-40-TL
% Mức: VD
Đề xuất các bước giải một bài toán thuộc dạng Nhận biết lực từ và đặc trưng của vectơ cảm ứng từ.
\loigiai{
Cần nêu được: Vectơ cảm ứng từ tại một điểm có hướng trùng với hướng của từ trường tại điểm đó; đơn vị là tesla (T). Khi vận dụng phải xác định đúng dữ kiện, phương chiều nếu có, điều kiện áp dụng, hệ đơn vị và kiểm tra tính hợp lí. Nhận định “Cảm ứng từ là đại lượng vô hướng và có đơn vị niutơn.” sai vì trái với kiến thức nêu trên.
}
\end{bt}

% ===== Câu 181 =====
% ID: L12C3B15-40-TLN
% Mức: TH
\begin{ex}
% ID: L12C3B15-40-TLN
% Mức: TH
Trong bốn phát biểu về Nhận biết lực từ và đặc trưng của vectơ cảm ứng từ, có bao nhiêu phát biểu đúng? (Chỉ ghi một số nguyên.) (1) Vectơ cảm ứng từ tại một điểm có hướng trùng với hướng của từ trường tại điểm đó; đơn vị là tesla (T). (2) Cần kiểm tra điều kiện áp dụng. (3) Cảm ứng từ là đại lượng vô hướng và có đơn vị niutơn. (4) Có thể bỏ qua đơn vị của kết quả.
\shortans{2}
\loigiai{
Đối chiếu từng phát biểu với kiến thức: Vectơ cảm ứng từ tại một điểm có hướng trùng với hướng của từ trường tại điểm đó; đơn vị là tesla (T). Số phát biểu đúng là 2.
}
\end{ex}

% ===== Câu 182 =====
% ID: L12C3B15-40-TN
% Mức: VD
\dangbt{Tính lực từ tác dụng lên đoạn dây dẫn mang dòng điện}
\begin{ex}
% ID: L12C3B15-40-TN
% Mức: VD
Kết luận đúng khi vận dụng Tính lực từ tác dụng lên đoạn dây dẫn mang dòng điện?
\choice
{\True Độ lớn lực từ lên đoạn dây thẳng dài l mang dòng $I$ trong từ trường đều là $F$ = BIl sin(alpha).}
{Lực từ được tính bằng $F$ = BI/l và không phụ thuộc góc.}
{Có thể bỏ qua phương, chiều và đơn vị của các đại lượng trong mọi trường hợp.}
{Kết quả không phụ thuộc dữ kiện và điều kiện áp dụng.}
\loigiai{
Độ lớn lực từ lên đoạn dây thẳng dài l mang dòng $I$ trong từ trường đều là $F$ = BIl sin(alpha). Vì vậy phương án $A$ đúng; các phương án còn lại trái với kiến thức hoặc điều kiện áp dụng.
}
\end{ex}

% ===== Câu 183 =====
% ID: L12C3B15-41-DS
% Mức: VD
\dangbt{Xác định phương, chiều lực từ bằng quy tắc bàn tay trái}
\begin{ex}
% ID: L12C3B15-41-DS
% Mức: VD
Xét Xác định phương, chiều lực từ bằng quy tắc bàn tay trái (Tình huống 4). Hãy xác định tính đúng, sai của bốn phát biểu sau.
\choiceTF
{Kết quả không phụ thuộc dữ kiện và có thể bỏ qua điều kiện.}
{\True Khi giải cần xác định đúng phương, chiều, điều kiện áp dụng và đơn vị.}
{Lực từ luôn cùng chiều dòng điện trong dây.}
{\True Lực từ tác dụng lên đoạn dây có phương vuông góc với cả dòng điện và vectơ cảm ứng từ; chiều xác định bằng quy tắc bàn tay trái.}
\loigiai{
\begin{itemchoice}
\item (Sai) Kết quả không phụ thuộc dữ kiện và có thể bỏ qua điều kiện. (Vì): Kết quả không phụ thuộc dữ kiện và có thể bỏ qua điều kiện. b) Đúng: Khi giải cần xác định đúng phương, chiều, điều kiện áp dụng và đơn vị. c) Sai: Lực từ luôn cùng chiều dòng điện trong dây. d) Đúng: Lực từ tác dụng lên đoạn dây có phương vuông góc với cả dòng điện và vectơ cảm ứng từ; chiều xác định bằng quy tắc bàn tay trái. Kiến thức cốt lõi: Lực từ tác dụng lên đoạn dây có phương vuông góc với cả dòng điện và vectơ cảm ứng từ; chiều xác định bằng quy tắc bàn tay trái
\item (Đúng) Khi giải cần xác định đúng phương, chiều, điều kiện áp dụng và đơn vị.
\item (Sai) Lực từ luôn cùng chiều dòng điện trong dây.
\item (Đúng) Lực từ tác dụng lên đoạn dây có phương vuông góc với cả dòng điện và vectơ cảm ứng từ; chiều xác định bằng quy tắc bàn tay trái.
\end{itemchoice}
}
\end{ex}

% ===== Câu 184 =====
% ID: L12C3B15-41-TL
% Mức: VD
\dangbt{Nhận biết lực từ và đặc trưng của vectơ cảm ứng từ}
\begin{bt}
% ID: L12C3B15-41-TL
% Mức: VD
Nêu một tình huống thực tiễn liên quan đến Nhận biết lực từ và đặc trưng của vectơ cảm ứng từ và giải thích bằng kiến thức Vật lí.
\loigiai{
Cần nêu được: Vectơ cảm ứng từ tại một điểm có hướng trùng với hướng của từ trường tại điểm đó; đơn vị là tesla (T). Khi vận dụng phải xác định đúng dữ kiện, phương chiều nếu có, điều kiện áp dụng, hệ đơn vị và kiểm tra tính hợp lí. Nhận định “Cảm ứng từ là đại lượng vô hướng và có đơn vị niutơn.” sai vì trái với kiến thức nêu trên.
}
\end{bt}

% ===== Câu 185 =====
% ID: L12C3B15-41-TLN
% Mức: VD
\begin{ex}
% ID: L12C3B15-41-TLN
% Mức: VD
Trong bốn phát biểu về Nhận biết lực từ và đặc trưng của vectơ cảm ứng từ, có bao nhiêu phát biểu đúng? (Chỉ ghi một số nguyên.) (1) Cảm ứng từ là đại lượng vô hướng và có đơn vị niutơn. (2) Vectơ cảm ứng từ tại một điểm có hướng trùng với hướng của từ trường tại điểm đó; đơn vị là tesla (T). (3) Cảm ứng từ là đại lượng vô hướng và có đơn vị niutơn. (4) Có thể bỏ qua đơn vị của kết quả.
\shortans{2}
\loigiai{
Đối chiếu từng phát biểu với kiến thức: Vectơ cảm ứng từ tại một điểm có hướng trùng với hướng của từ trường tại điểm đó; đơn vị là tesla (T). Số phát biểu đúng là 2.
}
\end{ex}

% ===== Câu 186 =====
% ID: L12C3B15-41-TN
% Mức: NB
\dangbt{Xác định cảm ứng từ từ phép đo lực từ}
\begin{ex}
% ID: L12C3B15-41-TN
% Mức: NB
Phát biểu nào sau đây đúng về Xác định cảm ứng từ từ phép đo lực từ?
\choice
{Khi dây vuông góc từ trường, cảm ứng từ bằng FIl.}
{Có thể bỏ qua phương, chiều và đơn vị của các đại lượng trong mọi trường hợp.}
{Kết quả không phụ thuộc dữ kiện và điều kiện áp dụng.}
{\True Khi dây vuông góc từ trường, cảm ứng từ được xác định bởi $B$ = F/(Il).}
\loigiai{
Khi dây vuông góc từ trường, cảm ứng từ được xác định bởi $B$ = F/(Il). Vì vậy phương án $D$ đúng; các phương án còn lại trái với kiến thức hoặc điều kiện áp dụng.
}
\end{ex}

% ===== Câu 187 =====
% ID: L12C3B15-42-DS
% Mức: VD
\dangbt{Xác định phương, chiều lực từ bằng quy tắc bàn tay trái}
\begin{ex}
% ID: L12C3B15-42-DS
% Mức: VD
Xét Xác định phương, chiều lực từ bằng quy tắc bàn tay trái (Tình huống 5). Hãy xác định tính đúng, sai của bốn phát biểu sau.
\choiceTF
{\True Lực từ tác dụng lên đoạn dây có phương vuông góc với cả dòng điện và vectơ cảm ứng từ; chiều xác định bằng quy tắc bàn tay trái.}
{Kết quả không phụ thuộc dữ kiện và có thể bỏ qua điều kiện.}
{Lực từ luôn cùng chiều dòng điện trong dây.}
{\True Khi giải cần xác định đúng phương, chiều, điều kiện áp dụng và đơn vị.}
\loigiai{
\begin{itemchoice}
\item (Đúng) Lực từ tác dụng lên đoạn dây có phương vuông góc với cả dòng điện và vectơ cảm ứng từ; chiều xác định bằng quy tắc bàn tay trái. (Vì): Lực từ tác dụng lên đoạn dây có phương vuông góc với cả dòng điện và vectơ cảm ứng từ; chiều xác định bằng quy tắc bàn tay trái. b) Sai: Kết quả không phụ thuộc dữ kiện và có thể bỏ qua điều kiện. c) Sai: Lực từ luôn cùng chiều dòng điện trong dây. d) Đúng: Khi giải cần xác định đúng phương, chiều, điều kiện áp dụng và đơn vị. Kiến thức cốt lõi: Lực từ tác dụng lên đoạn dây có phương vuông góc với cả dòng điện và vectơ cảm ứng từ; chiều xác định bằng quy tắc bàn tay trái
\item (Sai) Kết quả không phụ thuộc dữ kiện và có thể bỏ qua điều kiện.
\item (Sai) Lực từ luôn cùng chiều dòng điện trong dây.
\item (Đúng) Khi giải cần xác định đúng phương, chiều, điều kiện áp dụng và đơn vị.
\end{itemchoice}
}
\end{ex}

% ===== Câu 188 =====
% ID: L12C3B15-42-TL
% Mức: VDC
\dangbt{Nhận biết lực từ và đặc trưng của vectơ cảm ứng từ}
\begin{bt}
% ID: L12C3B15-42-TL
% Mức: VDC
So sánh nhận định đúng “Vectơ cảm ứng từ tại một điểm có hướng trùng với hướng của từ trường tại điểm đó; đơn vị là tesla (T).” với nhận định sai “Cảm ứng từ là đại lượng vô hướng và có đơn vị niutơn.”; chỉ rõ căn cứ Vật lí.
\loigiai{
Cần nêu được: Vectơ cảm ứng từ tại một điểm có hướng trùng với hướng của từ trường tại điểm đó; đơn vị là tesla (T). Khi vận dụng phải xác định đúng dữ kiện, phương chiều nếu có, điều kiện áp dụng, hệ đơn vị và kiểm tra tính hợp lí. Nhận định “Cảm ứng từ là đại lượng vô hướng và có đơn vị niutơn.” sai vì trái với kiến thức nêu trên.
}
\end{bt}

% ===== Câu 189 =====
% ID: L12C3B15-42-TLN
% Mức: VD
\begin{ex}
% ID: L12C3B15-42-TLN
% Mức: VD
Trong bốn phát biểu về Nhận biết lực từ và đặc trưng của vectơ cảm ứng từ, có bao nhiêu phát biểu đúng? (Chỉ ghi một số nguyên.) (1) Vectơ cảm ứng từ tại một điểm có hướng trùng với hướng của từ trường tại điểm đó; đơn vị là tesla (T). (2) Cần kiểm tra điều kiện áp dụng. (3) Cảm ứng từ là đại lượng vô hướng và có đơn vị niutơn. (4) Có thể bỏ qua đơn vị của kết quả.
\shortans{2}
\loigiai{
Đối chiếu từng phát biểu với kiến thức: Vectơ cảm ứng từ tại một điểm có hướng trùng với hướng của từ trường tại điểm đó; đơn vị là tesla (T). Số phát biểu đúng là 2.
}
\end{ex}

% ===== Câu 190 =====
% ID: L12C3B15-42-TN
% Mức: NB
\dangbt{Xác định cảm ứng từ từ phép đo lực từ}
\begin{ex}
% ID: L12C3B15-42-TN
% Mức: NB
Trong nội dung Xác định cảm ứng từ từ phép đo lực từ?
\choice
{\True Khi dây vuông góc từ trường, cảm ứng từ được xác định bởi $B$ = F/(Il).}
{Khi dây vuông góc từ trường, cảm ứng từ bằng FIl.}
{Có thể bỏ qua phương, chiều và đơn vị của các đại lượng trong mọi trường hợp.}
{Kết quả không phụ thuộc dữ kiện và điều kiện áp dụng.}
\loigiai{
Khi dây vuông góc từ trường, cảm ứng từ được xác định bởi $B$ = F/(Il). Vì vậy phương án $A$ đúng; các phương án còn lại trái với kiến thức hoặc điều kiện áp dụng.
}
\end{ex}

% ===== Câu 191 =====
% ID: L12C3B15-43-DS
% Mức: TH
\dangbt{Tính lực từ tác dụng lên đoạn dây dẫn mang dòng điện}
\begin{ex}
% ID: L12C3B15-43-DS
% Mức: TH
Xét Tính lực từ tác dụng lên đoạn dây dẫn mang dòng điện (Tình huống 1). Hãy xác định tính đúng, sai của bốn phát biểu sau.
\choiceTF
{\True Độ lớn lực từ lên đoạn dây thẳng dài l mang dòng $I$ trong từ trường đều là $F$ = BIl sin(alpha).}
{Lực từ được tính bằng $F$ = BI/l và không phụ thuộc góc.}
{\True Khi giải cần xác định đúng phương, chiều, điều kiện áp dụng và đơn vị.}
{Kết quả không phụ thuộc dữ kiện và có thể bỏ qua điều kiện.}
\loigiai{
\begin{itemchoice}
\item (Đúng) Độ lớn lực từ lên đoạn dây thẳng dài l mang dòng $I$ trong từ trường đều là $F$ = BIl sin(alpha). (Vì): lpha
\item (Sai) Lực từ được tính bằng $F$ = BI/l và không phụ thuộc góc.
\item (Đúng) Khi giải cần xác định đúng phương, chiều, điều kiện áp dụng và đơn vị.
\item (Sai) Kết quả không phụ thuộc dữ kiện và có thể bỏ qua điều kiện.
\end{itemchoice}
}
\end{ex}

% ===== Câu 192 =====
% ID: L12C3B15-43-TL
% Mức: TH
\dangbt{Xác định phương, chiều lực từ bằng quy tắc bàn tay trái}
\begin{bt}
% ID: L12C3B15-43-TL
% Mức: TH
Trình bày kiến thức cốt lõi của Xác định phương, chiều lực từ bằng quy tắc bàn tay trái và nêu điều kiện áp dụng.
\loigiai{
Cần nêu được: Lực từ tác dụng lên đoạn dây có phương vuông góc với cả dòng điện và vectơ cảm ứng từ; chiều xác định bằng quy tắc bàn tay trái. Khi vận dụng phải xác định đúng dữ kiện, phương chiều nếu có, điều kiện áp dụng, hệ đơn vị và kiểm tra tính hợp lí. Nhận định “Lực từ luôn cùng chiều dòng điện trong dây.” sai vì trái với kiến thức nêu trên.
}
\end{bt}

% ===== Câu 193 =====
% ID: L12C3B15-43-TLN
% Mức: VDC
\dangbt{Nhận biết lực từ và đặc trưng của vectơ cảm ứng từ}
\begin{ex}
% ID: L12C3B15-43-TLN
% Mức: VDC
Trong bốn phát biểu về Nhận biết lực từ và đặc trưng của vectơ cảm ứng từ, có bao nhiêu phát biểu đúng? (Chỉ ghi một số nguyên.) (1) Cảm ứng từ là đại lượng vô hướng và có đơn vị niutơn. (2) Vectơ cảm ứng từ tại một điểm có hướng trùng với hướng của từ trường tại điểm đó; đơn vị là tesla (T). (3) Phải dùng đúng định luật cho tình huống. (4) Vectơ cảm ứng từ tại một điểm có hướng trùng với hướng của từ trường tại điểm đó; đơn vị là tesla (T).
\shortans{3}
\loigiai{
Đối chiếu từng phát biểu với kiến thức: Vectơ cảm ứng từ tại một điểm có hướng trùng với hướng của từ trường tại điểm đó; đơn vị là tesla (T). Số phát biểu đúng là 3.
}
\end{ex}

% ===== Câu 194 =====
% ID: L12C3B15-43-TN
% Mức: NB
\dangbt{Xác định cảm ứng từ từ phép đo lực từ}
\begin{ex}
% ID: L12C3B15-43-TN
% Mức: NB
Tình huống 3: chọn kết luận chính xác về Xác định cảm ứng từ từ phép đo lực từ?
\choice
{Kết quả không phụ thuộc dữ kiện và điều kiện áp dụng.}
{\True Khi dây vuông góc từ trường, cảm ứng từ được xác định bởi $B$ = F/(Il).}
{Khi dây vuông góc từ trường, cảm ứng từ bằng FIl.}
{Có thể bỏ qua phương, chiều và đơn vị của các đại lượng trong mọi trường hợp.}
\loigiai{
Khi dây vuông góc từ trường, cảm ứng từ được xác định bởi $B$ = F/(Il). Vì vậy phương án $B$ đúng; các phương án còn lại trái với kiến thức hoặc điều kiện áp dụng.
}
\end{ex}

% ===== Câu 195 =====
% ID: L12C3B15-44-DS
% Mức: TH
\dangbt{Tính lực từ tác dụng lên đoạn dây dẫn mang dòng điện}
\begin{ex}
% ID: L12C3B15-44-DS
% Mức: TH
Xét Tính lực từ tác dụng lên đoạn dây dẫn mang dòng điện (Tình huống 2). Hãy xác định tính đúng, sai của bốn phát biểu sau.
\choiceTF
{Lực từ được tính bằng $F$ = BI/l và không phụ thuộc góc.}
{\True Độ lớn lực từ lên đoạn dây thẳng dài l mang dòng $I$ trong từ trường đều là $F$ = BIl sin(alpha).}
{Kết quả không phụ thuộc dữ kiện và có thể bỏ qua điều kiện.}
{\True Khi giải cần xác định đúng phương, chiều, điều kiện áp dụng và đơn vị.}
\loigiai{
\begin{itemchoice}
\item (Sai) Lực từ được tính bằng $F$ = BI/l và không phụ thuộc góc. (Vì): lpha
\item (Đúng) Độ lớn lực từ lên đoạn dây thẳng dài l mang dòng $I$ trong từ trường đều là $F$ = BIl sin(alpha).
\item (Sai) Kết quả không phụ thuộc dữ kiện và có thể bỏ qua điều kiện.
\item (Đúng) Khi giải cần xác định đúng phương, chiều, điều kiện áp dụng và đơn vị.
\end{itemchoice}
}
\end{ex}

% ===== Câu 196 =====
% ID: L12C3B15-44-TL
% Mức: TH
\dangbt{Xác định phương, chiều lực từ bằng quy tắc bàn tay trái}
\begin{bt}
% ID: L12C3B15-44-TL
% Mức: TH
Giải thích vì sao nhận định sau đúng: “Lực từ tác dụng lên đoạn dây có phương vuông góc với cả dòng điện và vectơ cảm ứng từ; chiều xác định bằng quy tắc bàn tay trái.”
\loigiai{
Cần nêu được: Lực từ tác dụng lên đoạn dây có phương vuông góc với cả dòng điện và vectơ cảm ứng từ; chiều xác định bằng quy tắc bàn tay trái. Khi vận dụng phải xác định đúng dữ kiện, phương chiều nếu có, điều kiện áp dụng, hệ đơn vị và kiểm tra tính hợp lí. Nhận định “Lực từ luôn cùng chiều dòng điện trong dây.” sai vì trái với kiến thức nêu trên.
}
\end{bt}

% ===== Câu 197 =====
% ID: L12C3B15-44-TLN
% Mức: TH
\begin{ex}
% ID: L12C3B15-44-TLN
% Mức: TH
Trong bốn phát biểu về Xác định phương, chiều lực từ bằng quy tắc bàn tay trái, có bao nhiêu phát biểu đúng? (Chỉ ghi một số nguyên.) (1) Lực từ tác dụng lên đoạn dây có phương vuông góc với cả dòng điện và vectơ cảm ứng từ; chiều xác định bằng quy tắc bàn tay trái. (2) Cần kiểm tra điều kiện áp dụng. (3) Lực từ luôn cùng chiều dòng điện trong dây. (4) Có thể bỏ qua đơn vị của kết quả.
\shortans{2}
\loigiai{
Đối chiếu từng phát biểu với kiến thức: Lực từ tác dụng lên đoạn dây có phương vuông góc với cả dòng điện và vectơ cảm ứng từ; chiều xác định bằng quy tắc bàn tay trái. Số phát biểu đúng là 2.
}
\end{ex}

% ===== Câu 198 =====
% ID: L12C3B15-44-TN
% Mức: TH
\dangbt{Xác định cảm ứng từ từ phép đo lực từ}
\begin{ex}
% ID: L12C3B15-44-TN
% Mức: TH
Nội dung nào mô tả đúng nhất Xác định cảm ứng từ từ phép đo lực từ?
\choice
{Có thể bỏ qua phương, chiều và đơn vị của các đại lượng trong mọi trường hợp.}
{Kết quả không phụ thuộc dữ kiện và điều kiện áp dụng.}
{\True Khi dây vuông góc từ trường, cảm ứng từ được xác định bởi $B$ = F/(Il).}
{Khi dây vuông góc từ trường, cảm ứng từ bằng FIl.}
\loigiai{
Khi dây vuông góc từ trường, cảm ứng từ được xác định bởi $B$ = F/(Il). Vì vậy phương án $C$ đúng; các phương án còn lại trái với kiến thức hoặc điều kiện áp dụng.
}
\end{ex}

% ===== Câu 199 =====
% ID: L12C3B15-45-DS
% Mức: VD
\dangbt{Tính lực từ tác dụng lên đoạn dây dẫn mang dòng điện}
\begin{ex}
% ID: L12C3B15-45-DS
% Mức: VD
Xét Tính lực từ tác dụng lên đoạn dây dẫn mang dòng điện (Tình huống 3). Hãy xác định tính đúng, sai của bốn phát biểu sau.
\choiceTF
{\True Khi giải cần xác định đúng phương, chiều, điều kiện áp dụng và đơn vị.}
{Kết quả không phụ thuộc dữ kiện và có thể bỏ qua điều kiện.}
{\True Độ lớn lực từ lên đoạn dây thẳng dài l mang dòng $I$ trong từ trường đều là $F$ = BIl sin(alpha).}
{Lực từ được tính bằng $F$ = BI/l và không phụ thuộc góc.}
\loigiai{
\begin{itemchoice}
\item (Đúng) Khi giải cần xác định đúng phương, chiều, điều kiện áp dụng và đơn vị. (Vì): lpha
\item (Sai) Kết quả không phụ thuộc dữ kiện và có thể bỏ qua điều kiện.
\item (Đúng) Độ lớn lực từ lên đoạn dây thẳng dài l mang dòng $I$ trong từ trường đều là $F$ = BIl sin(alpha).
\item (Sai) Lực từ được tính bằng $F$ = BI/l và không phụ thuộc góc.
\end{itemchoice}
}
\end{ex}

% ===== Câu 200 =====
% ID: L12C3B15-45-TL
% Mức: VD
\dangbt{Xác định phương, chiều lực từ bằng quy tắc bàn tay trái}
\begin{bt}
% ID: L12C3B15-45-TL
% Mức: VD
Phân tích sai lầm trong nhận định: “Lực từ luôn cùng chiều dòng điện trong dây.”
\loigiai{
Cần nêu được: Lực từ tác dụng lên đoạn dây có phương vuông góc với cả dòng điện và vectơ cảm ứng từ; chiều xác định bằng quy tắc bàn tay trái. Khi vận dụng phải xác định đúng dữ kiện, phương chiều nếu có, điều kiện áp dụng, hệ đơn vị và kiểm tra tính hợp lí. Nhận định “Lực từ luôn cùng chiều dòng điện trong dây.” sai vì trái với kiến thức nêu trên.
}
\end{bt}

% ===== Câu 201 =====
% ID: L12C3B15-45-TLN
% Mức: TH
\begin{ex}
% ID: L12C3B15-45-TLN
% Mức: TH
Trong bốn phát biểu về Xác định phương, chiều lực từ bằng quy tắc bàn tay trái, có bao nhiêu phát biểu đúng? (Chỉ ghi một số nguyên.) (1) Lực từ luôn cùng chiều dòng điện trong dây. (2) Lực từ tác dụng lên đoạn dây có phương vuông góc với cả dòng điện và vectơ cảm ứng từ; chiều xác định bằng quy tắc bàn tay trái. (3) Lực từ luôn cùng chiều dòng điện trong dây. (4) Có thể bỏ qua đơn vị của kết quả.
\shortans{1}
\loigiai{
Đối chiếu từng phát biểu với kiến thức: Lực từ tác dụng lên đoạn dây có phương vuông góc với cả dòng điện và vectơ cảm ứng từ; chiều xác định bằng quy tắc bàn tay trái. Số phát biểu đúng là 1.
}
\end{ex}

% ===== Câu 202 =====
% ID: L12C3B15-45-TN
% Mức: TH
\dangbt{Xác định cảm ứng từ từ phép đo lực từ}
\begin{ex}
% ID: L12C3B15-45-TN
% Mức: TH
Khi phân tích Xác định cảm ứng từ từ phép đo lực từ?
\choice
{Khi dây vuông góc từ trường, cảm ứng từ bằng FIl.}
{Có thể bỏ qua phương, chiều và đơn vị của các đại lượng trong mọi trường hợp.}
{Kết quả không phụ thuộc dữ kiện và điều kiện áp dụng.}
{\True Khi dây vuông góc từ trường, cảm ứng từ được xác định bởi $B$ = F/(Il).}
\loigiai{
Khi dây vuông góc từ trường, cảm ứng từ được xác định bởi $B$ = F/(Il). Vì vậy phương án $D$ đúng; các phương án còn lại trái với kiến thức hoặc điều kiện áp dụng.
}
\end{ex}

% ===== Câu 203 =====
% ID: L12C3B15-46-DS
% Mức: VD
\dangbt{Tính lực từ tác dụng lên đoạn dây dẫn mang dòng điện}
\begin{ex}
% ID: L12C3B15-46-DS
% Mức: VD
Xét Tính lực từ tác dụng lên đoạn dây dẫn mang dòng điện (Tình huống 4). Hãy xác định tính đúng, sai của bốn phát biểu sau.
\choiceTF
{Kết quả không phụ thuộc dữ kiện và có thể bỏ qua điều kiện.}
{\True Khi giải cần xác định đúng phương, chiều, điều kiện áp dụng và đơn vị.}
{Lực từ được tính bằng $F$ = BI/l và không phụ thuộc góc.}
{\True Độ lớn lực từ lên đoạn dây thẳng dài l mang dòng $I$ trong từ trường đều là $F$ = BIl sin(alpha).}
\loigiai{
\begin{itemchoice}
\item (Sai) Kết quả không phụ thuộc dữ kiện và có thể bỏ qua điều kiện. (Vì): lpha
\item (Đúng) Khi giải cần xác định đúng phương, chiều, điều kiện áp dụng và đơn vị.
\item (Sai) Lực từ được tính bằng $F$ = BI/l và không phụ thuộc góc.
\item (Đúng) Độ lớn lực từ lên đoạn dây thẳng dài l mang dòng $I$ trong từ trường đều là $F$ = BIl sin(alpha).
\end{itemchoice}
}
\end{ex}

% ===== Câu 204 =====
% ID: L12C3B15-46-TL
% Mức: VD
\dangbt{Xác định phương, chiều lực từ bằng quy tắc bàn tay trái}
\begin{bt}
% ID: L12C3B15-46-TL
% Mức: VD
Đề xuất các bước giải một bài toán thuộc dạng Xác định phương, chiều lực từ bằng quy tắc bàn tay trái.
\loigiai{
Cần nêu được: Lực từ tác dụng lên đoạn dây có phương vuông góc với cả dòng điện và vectơ cảm ứng từ; chiều xác định bằng quy tắc bàn tay trái. Khi vận dụng phải xác định đúng dữ kiện, phương chiều nếu có, điều kiện áp dụng, hệ đơn vị và kiểm tra tính hợp lí. Nhận định “Lực từ luôn cùng chiều dòng điện trong dây.” sai vì trái với kiến thức nêu trên.
}
\end{bt}

% ===== Câu 205 =====
% ID: L12C3B15-46-TLN
% Mức: TH
\begin{ex}
% ID: L12C3B15-46-TLN
% Mức: TH
Trong bốn phát biểu về Xác định phương, chiều lực từ bằng quy tắc bàn tay trái, có bao nhiêu phát biểu đúng? (Chỉ ghi một số nguyên.) (1) Lực từ tác dụng lên đoạn dây có phương vuông góc với cả dòng điện và vectơ cảm ứng từ; chiều xác định bằng quy tắc bàn tay trái. (2) Cần kiểm tra điều kiện áp dụng. (3) Lực từ luôn cùng chiều dòng điện trong dây. (4) Có thể bỏ qua đơn vị của kết quả.
\shortans{2}
\loigiai{
Đối chiếu từng phát biểu với kiến thức: Lực từ tác dụng lên đoạn dây có phương vuông góc với cả dòng điện và vectơ cảm ứng từ; chiều xác định bằng quy tắc bàn tay trái. Số phát biểu đúng là 2.
}
\end{ex}

% ===== Câu 206 =====
% ID: L12C3B15-46-TN
% Mức: TH
\dangbt{Xác định cảm ứng từ từ phép đo lực từ}
\begin{ex}
% ID: L12C3B15-46-TN
% Mức: TH
Một học sinh ôn tập Xác định cảm ứng từ từ phép đo lực từ?
\choice
{\True Khi dây vuông góc từ trường, cảm ứng từ được xác định bởi $B$ = F/(Il).}
{Khi dây vuông góc từ trường, cảm ứng từ bằng FIl.}
{Có thể bỏ qua phương, chiều và đơn vị của các đại lượng trong mọi trường hợp.}
{Kết quả không phụ thuộc dữ kiện và điều kiện áp dụng.}
\loigiai{
Khi dây vuông góc từ trường, cảm ứng từ được xác định bởi $B$ = F/(Il). Vì vậy phương án $A$ đúng; các phương án còn lại trái với kiến thức hoặc điều kiện áp dụng.
}
\end{ex}

% ===== Câu 207 =====
% ID: L12C3B15-47-DS
% Mức: VD
\dangbt{Tính lực từ tác dụng lên đoạn dây dẫn mang dòng điện}
\begin{ex}
% ID: L12C3B15-47-DS
% Mức: VD
Xét Tính lực từ tác dụng lên đoạn dây dẫn mang dòng điện (Tình huống 5). Hãy xác định tính đúng, sai của bốn phát biểu sau.
\choiceTF
{\True Độ lớn lực từ lên đoạn dây thẳng dài l mang dòng $I$ trong từ trường đều là $F$ = BIl sin(alpha).}
{Kết quả không phụ thuộc dữ kiện và có thể bỏ qua điều kiện.}
{Lực từ được tính bằng $F$ = BI/l và không phụ thuộc góc.}
{\True Khi giải cần xác định đúng phương, chiều, điều kiện áp dụng và đơn vị.}
\loigiai{
\begin{itemchoice}
\item (Đúng) Độ lớn lực từ lên đoạn dây thẳng dài l mang dòng $I$ trong từ trường đều là $F$ = BIl sin(alpha). (Vì): lpha
\item (Sai) Kết quả không phụ thuộc dữ kiện và có thể bỏ qua điều kiện.
\item (Sai) Lực từ được tính bằng $F$ = BI/l và không phụ thuộc góc.
\item (Đúng) Khi giải cần xác định đúng phương, chiều, điều kiện áp dụng và đơn vị.
\end{itemchoice}
}
\end{ex}

% ===== Câu 208 =====
% ID: L12C3B15-47-TL
% Mức: VD
\dangbt{Xác định phương, chiều lực từ bằng quy tắc bàn tay trái}
\begin{bt}
% ID: L12C3B15-47-TL
% Mức: VD
Nêu một tình huống thực tiễn liên quan đến Xác định phương, chiều lực từ bằng quy tắc bàn tay trái và giải thích bằng kiến thức Vật lí.
\loigiai{
Cần nêu được: Lực từ tác dụng lên đoạn dây có phương vuông góc với cả dòng điện và vectơ cảm ứng từ; chiều xác định bằng quy tắc bàn tay trái. Khi vận dụng phải xác định đúng dữ kiện, phương chiều nếu có, điều kiện áp dụng, hệ đơn vị và kiểm tra tính hợp lí. Nhận định “Lực từ luôn cùng chiều dòng điện trong dây.” sai vì trái với kiến thức nêu trên.
}
\end{bt}

% ===== Câu 209 =====
% ID: L12C3B15-47-TLN
% Mức: VD
\begin{ex}
% ID: L12C3B15-47-TLN
% Mức: VD
Trong bốn phát biểu về Xác định phương, chiều lực từ bằng quy tắc bàn tay trái, có bao nhiêu phát biểu đúng? (Chỉ ghi một số nguyên.) (1) Lực từ luôn cùng chiều dòng điện trong dây. (2) Lực từ tác dụng lên đoạn dây có phương vuông góc với cả dòng điện và vectơ cảm ứng từ; chiều xác định bằng quy tắc bàn tay trái. (3) Lực từ luôn cùng chiều dòng điện trong dây. (4) Có thể bỏ qua đơn vị của kết quả.
\shortans{2}
\loigiai{
Đối chiếu từng phát biểu với kiến thức: Lực từ tác dụng lên đoạn dây có phương vuông góc với cả dòng điện và vectơ cảm ứng từ; chiều xác định bằng quy tắc bàn tay trái. Số phát biểu đúng là 2.
}
\end{ex}

% ===== Câu 210 =====
% ID: L12C3B15-47-TN
% Mức: TH
\dangbt{Xác định cảm ứng từ từ phép đo lực từ}
\begin{ex}
% ID: L12C3B15-47-TN
% Mức: TH
Chọn phát biểu không trái với kiến thức về Xác định cảm ứng từ từ phép đo lực từ?
\choice
{Kết quả không phụ thuộc dữ kiện và điều kiện áp dụng.}
{\True Khi dây vuông góc từ trường, cảm ứng từ được xác định bởi $B$ = F/(Il).}
{Khi dây vuông góc từ trường, cảm ứng từ bằng FIl.}
{Có thể bỏ qua phương, chiều và đơn vị của các đại lượng trong mọi trường hợp.}
\loigiai{
Khi dây vuông góc từ trường, cảm ứng từ được xác định bởi $B$ = F/(Il). Vì vậy phương án $B$ đúng; các phương án còn lại trái với kiến thức hoặc điều kiện áp dụng.
}
\end{ex}

% ===== Câu 211 =====
% ID: L12C3B15-48-DS
% Mức: TH
\begin{ex}
% ID: L12C3B15-48-DS
% Mức: TH
Xét Xác định cảm ứng từ từ phép đo lực từ (Tình huống 1). Hãy xác định tính đúng, sai của bốn phát biểu sau.
\choiceTF
{\True Khi dây vuông góc từ trường, cảm ứng từ được xác định bởi $B$ = F/(Il).}
{Khi dây vuông góc từ trường, cảm ứng từ bằng FIl.}
{\True Khi giải cần xác định đúng phương, chiều, điều kiện áp dụng và đơn vị.}
{Kết quả không phụ thuộc dữ kiện và có thể bỏ qua điều kiện.}
\loigiai{
\begin{itemchoice}
\item (Đúng) Khi dây vuông góc từ trường, cảm ứng từ được xác định bởi $B$ = F/(Il). (Vì): Khi dây vuông góc từ trường, cảm ứng từ được xác định bởi $B$ = F/(Il). b) Sai: Khi dây vuông góc từ trường, cảm ứng từ bằng FIl. c) Đúng: Khi giải cần xác định đúng phương, chiều, điều kiện áp dụng và đơn vị. d) Sai: Kết quả không phụ thuộc dữ kiện và có thể bỏ qua điều kiện. Kiến thức cốt lõi: Khi dây vuông góc từ trường, cảm ứng từ được xác định bởi $B$ = F/(Il
\item (Sai) Khi dây vuông góc từ trường, cảm ứng từ bằng FIl.
\item (Đúng) Khi giải cần xác định đúng phương, chiều, điều kiện áp dụng và đơn vị.
\item (Sai) Kết quả không phụ thuộc dữ kiện và có thể bỏ qua điều kiện.
\end{itemchoice}
}
\end{ex}

% ===== Câu 212 =====
% ID: L12C3B15-48-TL
% Mức: VDC
\dangbt{Xác định phương, chiều lực từ bằng quy tắc bàn tay trái}
\begin{bt}
% ID: L12C3B15-48-TL
% Mức: VDC
So sánh nhận định đúng “Lực từ tác dụng lên đoạn dây có phương vuông góc với cả dòng điện và vectơ cảm ứng từ; chiều xác định bằng quy tắc bàn tay trái.” với nhận định sai “Lực từ luôn cùng chiều dòng điện trong dây.”; chỉ rõ căn cứ Vật lí.
\loigiai{
Cần nêu được: Lực từ tác dụng lên đoạn dây có phương vuông góc với cả dòng điện và vectơ cảm ứng từ; chiều xác định bằng quy tắc bàn tay trái. Khi vận dụng phải xác định đúng dữ kiện, phương chiều nếu có, điều kiện áp dụng, hệ đơn vị và kiểm tra tính hợp lí. Nhận định “Lực từ luôn cùng chiều dòng điện trong dây.” sai vì trái với kiến thức nêu trên.
}
\end{bt}

% ===== Câu 213 =====
% ID: L12C3B15-48-TLN
% Mức: VD
\begin{ex}
% ID: L12C3B15-48-TLN
% Mức: VD
Trong bốn phát biểu về Xác định phương, chiều lực từ bằng quy tắc bàn tay trái, có bao nhiêu phát biểu đúng? (Chỉ ghi một số nguyên.) (1) Lực từ tác dụng lên đoạn dây có phương vuông góc với cả dòng điện và vectơ cảm ứng từ; chiều xác định bằng quy tắc bàn tay trái. (2) Cần kiểm tra điều kiện áp dụng. (3) Lực từ luôn cùng chiều dòng điện trong dây. (4) Có thể bỏ qua đơn vị của kết quả.
\shortans{2}
\loigiai{
Đối chiếu từng phát biểu với kiến thức: Lực từ tác dụng lên đoạn dây có phương vuông góc với cả dòng điện và vectơ cảm ứng từ; chiều xác định bằng quy tắc bàn tay trái. Số phát biểu đúng là 2.
}
\end{ex}

% ===== Câu 214 =====
% ID: L12C3B15-48-TN
% Mức: VD
\dangbt{Xác định cảm ứng từ từ phép đo lực từ}
\begin{ex}
% ID: L12C3B15-48-TN
% Mức: VD
Cơ sở Vật lí đúng dùng để giải Xác định cảm ứng từ từ phép đo lực từ?
\choice
{Có thể bỏ qua phương, chiều và đơn vị của các đại lượng trong mọi trường hợp.}
{Kết quả không phụ thuộc dữ kiện và điều kiện áp dụng.}
{\True Khi dây vuông góc từ trường, cảm ứng từ được xác định bởi $B$ = F/(Il).}
{Khi dây vuông góc từ trường, cảm ứng từ bằng FIl.}
\loigiai{
Khi dây vuông góc từ trường, cảm ứng từ được xác định bởi $B$ = F/(Il). Vì vậy phương án $C$ đúng; các phương án còn lại trái với kiến thức hoặc điều kiện áp dụng.
}
\end{ex}

% ===== Câu 215 =====
% ID: L12C3B15-49-DS
% Mức: TH
\begin{ex}
% ID: L12C3B15-49-DS
% Mức: TH
Xét Xác định cảm ứng từ từ phép đo lực từ (Tình huống 2). Hãy xác định tính đúng, sai của bốn phát biểu sau.
\choiceTF
{Khi dây vuông góc từ trường, cảm ứng từ bằng FIl.}
{\True Khi dây vuông góc từ trường, cảm ứng từ được xác định bởi $B$ = F/(Il).}
{Kết quả không phụ thuộc dữ kiện và có thể bỏ qua điều kiện.}
{\True Khi giải cần xác định đúng phương, chiều, điều kiện áp dụng và đơn vị.}
\loigiai{
\begin{itemchoice}
\item (Sai) Khi dây vuông góc từ trường, cảm ứng từ bằng FIl. (Vì): Khi dây vuông góc từ trường, cảm ứng từ bằng FIl. b) Đúng: Khi dây vuông góc từ trường, cảm ứng từ được xác định bởi $B$ = F/(Il). c) Sai: Kết quả không phụ thuộc dữ kiện và có thể bỏ qua điều kiện. d) Đúng: Khi giải cần xác định đúng phương, chiều, điều kiện áp dụng và đơn vị. Kiến thức cốt lõi: Khi dây vuông góc từ trường, cảm ứng từ được xác định bởi $B$ = F/(Il
\item (Đúng) Khi dây vuông góc từ trường, cảm ứng từ được xác định bởi $B$ = F/(Il).
\item (Sai) Kết quả không phụ thuộc dữ kiện và có thể bỏ qua điều kiện.
\item (Đúng) Khi giải cần xác định đúng phương, chiều, điều kiện áp dụng và đơn vị.
\end{itemchoice}
}
\end{ex}

% ===== Câu 216 =====
% ID: L12C3B15-49-TL
% Mức: TH
\dangbt{Tính lực từ tác dụng lên đoạn dây dẫn mang dòng điện}
\begin{bt}
% ID: L12C3B15-49-TL
% Mức: TH
Trình bày kiến thức cốt lõi của Tính lực từ tác dụng lên đoạn dây dẫn mang dòng điện và nêu điều kiện áp dụng.
\loigiai{
Cần nêu được: Độ lớn lực từ lên đoạn dây thẳng dài l mang dòng $I$ trong từ trường đều là $F$ = BIl sin(alpha). Khi vận dụng phải xác định đúng dữ kiện, phương chiều nếu có, điều kiện áp dụng, hệ đơn vị và kiểm tra tính hợp lí. Nhận định “Lực từ được tính bằng $F$ = BI/l và không phụ thuộc góc.” sai vì trái với kiến thức nêu trên.
}
\end{bt}

% ===== Câu 217 =====
% ID: L12C3B15-49-TLN
% Mức: VDC
\dangbt{Xác định phương, chiều lực từ bằng quy tắc bàn tay trái}
\begin{ex}
% ID: L12C3B15-49-TLN
% Mức: VDC
Trong bốn phát biểu về Xác định phương, chiều lực từ bằng quy tắc bàn tay trái, có bao nhiêu phát biểu đúng? (Chỉ ghi một số nguyên.) (1) Lực từ luôn cùng chiều dòng điện trong dây. (2) Lực từ tác dụng lên đoạn dây có phương vuông góc với cả dòng điện và vectơ cảm ứng từ; chiều xác định bằng quy tắc bàn tay trái. (3) Phải dùng đúng định luật cho tình huống. (4) Lực từ tác dụng lên đoạn dây có phương vuông góc với cả dòng điện và vectơ cảm ứng từ; chiều xác định bằng quy tắc bàn tay trái.
\shortans{3}
\loigiai{
Đối chiếu từng phát biểu với kiến thức: Lực từ tác dụng lên đoạn dây có phương vuông góc với cả dòng điện và vectơ cảm ứng từ; chiều xác định bằng quy tắc bàn tay trái. Số phát biểu đúng là 3.
}
\end{ex}

% ===== Câu 218 =====
% ID: L12C3B15-49-TN
% Mức: VD
\dangbt{Xác định cảm ứng từ từ phép đo lực từ}
\begin{ex}
% ID: L12C3B15-49-TN
% Mức: VD
Trong một bài thực hành về Xác định cảm ứng từ từ phép đo lực từ?
\choice
{Khi dây vuông góc từ trường, cảm ứng từ bằng FIl.}
{Có thể bỏ qua phương, chiều và đơn vị của các đại lượng trong mọi trường hợp.}
{Kết quả không phụ thuộc dữ kiện và điều kiện áp dụng.}
{\True Khi dây vuông góc từ trường, cảm ứng từ được xác định bởi $B$ = F/(Il).}
\loigiai{
Khi dây vuông góc từ trường, cảm ứng từ được xác định bởi $B$ = F/(Il). Vì vậy phương án $D$ đúng; các phương án còn lại trái với kiến thức hoặc điều kiện áp dụng.
}
\end{ex}

% ===== Câu 219 =====
% ID: L12C3B15-50-DS
% Mức: VD
\begin{ex}
% ID: L12C3B15-50-DS
% Mức: VD
Xét Xác định cảm ứng từ từ phép đo lực từ (Tình huống 3). Hãy xác định tính đúng, sai của bốn phát biểu sau.
\choiceTF
{\True Khi giải cần xác định đúng phương, chiều, điều kiện áp dụng và đơn vị.}
{Kết quả không phụ thuộc dữ kiện và có thể bỏ qua điều kiện.}
{\True Khi dây vuông góc từ trường, cảm ứng từ được xác định bởi $B$ = F/(Il).}
{Khi dây vuông góc từ trường, cảm ứng từ bằng FIl.}
\loigiai{
\begin{itemchoice}
\item (Đúng) Khi giải cần xác định đúng phương, chiều, điều kiện áp dụng và đơn vị. (Vì): Khi giải cần xác định đúng phương, chiều, điều kiện áp dụng và đơn vị. b) Sai: Kết quả không phụ thuộc dữ kiện và có thể bỏ qua điều kiện. c) Đúng: Khi dây vuông góc từ trường, cảm ứng từ được xác định bởi $B$ = F/(Il). d) Sai: Khi dây vuông góc từ trường, cảm ứng từ bằng FIl. Kiến thức cốt lõi: Khi dây vuông góc từ trường, cảm ứng từ được xác định bởi $B$ = F/(Il
\item (Sai) Kết quả không phụ thuộc dữ kiện và có thể bỏ qua điều kiện.
\item (Đúng) Khi dây vuông góc từ trường, cảm ứng từ được xác định bởi $B$ = F/(Il).
\item (Sai) Khi dây vuông góc từ trường, cảm ứng từ bằng FIl.
\end{itemchoice}
}
\end{ex}

% ===== Câu 220 =====
% ID: L12C3B15-50-TL
% Mức: TH
\dangbt{Tính lực từ tác dụng lên đoạn dây dẫn mang dòng điện}
\begin{bt}
% ID: L12C3B15-50-TL
% Mức: TH
Giải thích vì sao nhận định sau đúng: “Độ lớn lực từ lên đoạn dây thẳng dài l mang dòng $I$ trong từ trường đều là $F$ = BIl sin(alpha).”
\loigiai{
Cần nêu được: Độ lớn lực từ lên đoạn dây thẳng dài l mang dòng $I$ trong từ trường đều là $F$ = BIl sin(alpha). Khi vận dụng phải xác định đúng dữ kiện, phương chiều nếu có, điều kiện áp dụng, hệ đơn vị và kiểm tra tính hợp lí. Nhận định “Lực từ được tính bằng $F$ = BI/l và không phụ thuộc góc.” sai vì trái với kiến thức nêu trên.
}
\end{bt}

% ===== Câu 221 =====
% ID: L12C3B15-50-TLN
% Mức: TH
\begin{ex}
% ID: L12C3B15-50-TLN
% Mức: TH
Trong bốn phát biểu về Tính lực từ tác dụng lên đoạn dây dẫn mang dòng điện, có bao nhiêu phát biểu đúng? (Chỉ ghi một số nguyên.) (1) Độ lớn lực từ lên đoạn dây thẳng dài l mang dòng $I$ trong từ trường đều là $F$ = BIl sin(alpha). (2) Cần kiểm tra điều kiện áp dụng. (3) Lực từ được tính bằng $F$ = BI/l và không phụ thuộc góc. (4) Có thể bỏ qua đơn vị của kết quả.
\shortans{2}
\loigiai{
Đối chiếu từng phát biểu với kiến thức: Độ lớn lực từ lên đoạn dây thẳng dài l mang dòng $I$ trong từ trường đều là $F$ = BIl sin(alpha). Số phát biểu đúng là 2.
}
\end{ex}

% ===== Câu 222 =====
% ID: L12C3B15-50-TN
% Mức: VD
\dangbt{Xác định cảm ứng từ từ phép đo lực từ}
\begin{ex}
% ID: L12C3B15-50-TN
% Mức: VD
Kết luận đúng khi vận dụng Xác định cảm ứng từ từ phép đo lực từ?
\choice
{\True Khi dây vuông góc từ trường, cảm ứng từ được xác định bởi $B$ = F/(Il).}
{Khi dây vuông góc từ trường, cảm ứng từ bằng FIl.}
{Có thể bỏ qua phương, chiều và đơn vị của các đại lượng trong mọi trường hợp.}
{Kết quả không phụ thuộc dữ kiện và điều kiện áp dụng.}
\loigiai{
Khi dây vuông góc từ trường, cảm ứng từ được xác định bởi $B$ = F/(Il). Vì vậy phương án $A$ đúng; các phương án còn lại trái với kiến thức hoặc điều kiện áp dụng.
}
\end{ex}

% ===== Câu 223 =====
% ID: L12C3B15-51-DS
% Mức: VD
\begin{ex}
% ID: L12C3B15-51-DS
% Mức: VD
Xét Xác định cảm ứng từ từ phép đo lực từ (Tình huống 4). Hãy xác định tính đúng, sai của bốn phát biểu sau.
\choiceTF
{Kết quả không phụ thuộc dữ kiện và có thể bỏ qua điều kiện.}
{\True Khi giải cần xác định đúng phương, chiều, điều kiện áp dụng và đơn vị.}
{Khi dây vuông góc từ trường, cảm ứng từ bằng FIl.}
{\True Khi dây vuông góc từ trường, cảm ứng từ được xác định bởi $B$ = F/(Il).}
\loigiai{
\begin{itemchoice}
\item (Sai) Kết quả không phụ thuộc dữ kiện và có thể bỏ qua điều kiện. (Vì): Kết quả không phụ thuộc dữ kiện và có thể bỏ qua điều kiện. b) Đúng: Khi giải cần xác định đúng phương, chiều, điều kiện áp dụng và đơn vị. c) Sai: Khi dây vuông góc từ trường, cảm ứng từ bằng FIl. d) Đúng: Khi dây vuông góc từ trường, cảm ứng từ được xác định bởi $B$ = F/(Il). Kiến thức cốt lõi: Khi dây vuông góc từ trường, cảm ứng từ được xác định bởi $B$ = F/(Il
\item (Đúng) Khi giải cần xác định đúng phương, chiều, điều kiện áp dụng và đơn vị.
\item (Sai) Khi dây vuông góc từ trường, cảm ứng từ bằng FIl.
\item (Đúng) Khi dây vuông góc từ trường, cảm ứng từ được xác định bởi $B$ = F/(Il).
\end{itemchoice}
}
\end{ex}

% ===== Câu 224 =====
% ID: L12C3B15-51-TL
% Mức: VD
\dangbt{Tính lực từ tác dụng lên đoạn dây dẫn mang dòng điện}
\begin{bt}
% ID: L12C3B15-51-TL
% Mức: VD
Phân tích sai lầm trong nhận định: “Lực từ được tính bằng $F$ = BI/l và không phụ thuộc góc.”
\loigiai{
Cần nêu được: Độ lớn lực từ lên đoạn dây thẳng dài l mang dòng $I$ trong từ trường đều là $F$ = BIl sin(alpha). Khi vận dụng phải xác định đúng dữ kiện, phương chiều nếu có, điều kiện áp dụng, hệ đơn vị và kiểm tra tính hợp lí. Nhận định “Lực từ được tính bằng $F$ = BI/l và không phụ thuộc góc.” sai vì trái với kiến thức nêu trên.
}
\end{bt}

% ===== Câu 225 =====
% ID: L12C3B15-51-TLN
% Mức: TH
\begin{ex}
% ID: L12C3B15-51-TLN
% Mức: TH
Trong bốn phát biểu về Tính lực từ tác dụng lên đoạn dây dẫn mang dòng điện, có bao nhiêu phát biểu đúng? (Chỉ ghi một số nguyên.) (1) Lực từ được tính bằng $F$ = BI/l và không phụ thuộc góc. (2) Độ lớn lực từ lên đoạn dây thẳng dài l mang dòng $I$ trong từ trường đều là $F$ = BIl sin(alpha). (3) Lực từ được tính bằng $F$ = BI/l và không phụ thuộc góc. (4) Có thể bỏ qua đơn vị của kết quả.
\shortans{1}
\loigiai{
Đối chiếu từng phát biểu với kiến thức: Độ lớn lực từ lên đoạn dây thẳng dài l mang dòng $I$ trong từ trường đều là $F$ = BIl sin(alpha). Số phát biểu đúng là 1.
}
\end{ex}

% ===== Câu 226 =====
% ID: L12C3B15-51-TN
% Mức: VD
\dangbt{Xác định cảm ứng từ từ phép đo lực từ}
\begin{ex}
% ID: L12C3B15-51-TN
% Mức: VD
Một dây dẫn thẳng nằm ngang có dòng điện I=16Achạy qua theo hướng từ Tây sang Đông. Tại vị trí này, cảm ứng từ $\overrightarrow{B}$ của từ trường Trái Đất có phương nằm ngang và hướng về phía Bắc. Lực từ tác dụng lên $1\,\text{m}$ chiều dài của đoạn dây là F=0,6mN. Độ lớn của $\overrightarrow{B}$ là
\choice
{27 $\muT$}
{30 $\muT$}
{\True 37,5 $\muT$}
{70 $\muT$}
\loigiai{
:Độ lớn cảm ứng từ là Ta có: $F=IB\ell \sin \alpha \Rightarrow B=\dfrac{F}{IB\ell \sin \alpha }=\dfrac{0,{{6.10}^{-3}}}{16.1.sin{{90}^{0}}}=3,{{75.10}^{-5}T=37,5\left( \mu T \right)$
}
\end{ex}

% ===== Câu 227 =====
% ID: L12C3B15-52-DS
% Mức: VD
\begin{ex}
% ID: L12C3B15-52-DS
% Mức: VD
Xét Xác định cảm ứng từ từ phép đo lực từ (Tình huống 5). Hãy xác định tính đúng, sai của bốn phát biểu sau.
\choiceTF
{\True Khi dây vuông góc từ trường, cảm ứng từ được xác định bởi $B$ = F/(Il).}
{Kết quả không phụ thuộc dữ kiện và có thể bỏ qua điều kiện.}
{Khi dây vuông góc từ trường, cảm ứng từ bằng FIl.}
{\True Khi giải cần xác định đúng phương, chiều, điều kiện áp dụng và đơn vị.}
\loigiai{
\begin{itemchoice}
\item (Đúng) Khi dây vuông góc từ trường, cảm ứng từ được xác định bởi $B$ = F/(Il). (Vì): Khi dây vuông góc từ trường, cảm ứng từ được xác định bởi $B$ = F/(Il). b) Sai: Kết quả không phụ thuộc dữ kiện và có thể bỏ qua điều kiện. c) Sai: Khi dây vuông góc từ trường, cảm ứng từ bằng FIl. d) Đúng: Khi giải cần xác định đúng phương, chiều, điều kiện áp dụng và đơn vị. Kiến thức cốt lõi: Khi dây vuông góc từ trường, cảm ứng từ được xác định bởi $B$ = F/(Il
\item (Sai) Kết quả không phụ thuộc dữ kiện và có thể bỏ qua điều kiện.
\item (Sai) Khi dây vuông góc từ trường, cảm ứng từ bằng FIl.
\item (Đúng) Khi giải cần xác định đúng phương, chiều, điều kiện áp dụng và đơn vị.
\end{itemchoice}
}
\end{ex}

% ===== Câu 228 =====
% ID: L12C3B15-52-TL
% Mức: VD
\dangbt{Tính lực từ tác dụng lên đoạn dây dẫn mang dòng điện}
\begin{bt}
% ID: L12C3B15-52-TL
% Mức: VD
Đề xuất các bước giải một bài toán thuộc dạng Tính lực từ tác dụng lên đoạn dây dẫn mang dòng điện.
\loigiai{
Cần nêu được: Độ lớn lực từ lên đoạn dây thẳng dài l mang dòng $I$ trong từ trường đều là $F$ = BIl sin(alpha). Khi vận dụng phải xác định đúng dữ kiện, phương chiều nếu có, điều kiện áp dụng, hệ đơn vị và kiểm tra tính hợp lí. Nhận định “Lực từ được tính bằng $F$ = BI/l và không phụ thuộc góc.” sai vì trái với kiến thức nêu trên.
}
\end{bt}

% ===== Câu 229 =====
% ID: L12C3B15-52-TLN
% Mức: TH
\begin{ex}
% ID: L12C3B15-52-TLN
% Mức: TH
Trong bốn phát biểu về Tính lực từ tác dụng lên đoạn dây dẫn mang dòng điện, có bao nhiêu phát biểu đúng? (Chỉ ghi một số nguyên.) (1) Độ lớn lực từ lên đoạn dây thẳng dài l mang dòng $I$ trong từ trường đều là $F$ = BIl sin(alpha). (2) Cần kiểm tra điều kiện áp dụng. (3) Lực từ được tính bằng $F$ = BI/l và không phụ thuộc góc. (4) Có thể bỏ qua đơn vị của kết quả.
\shortans{2}
\loigiai{
Đối chiếu từng phát biểu với kiến thức: Độ lớn lực từ lên đoạn dây thẳng dài l mang dòng $I$ trong từ trường đều là $F$ = BIl sin(alpha). Số phát biểu đúng là 2.
}
\end{ex}

% ===== Câu 230 =====
% ID: L12C3B15-52-TN
% Mức: NB
\dangbt{Xác định phương, chiều lực từ bằng quy tắc bàn tay trái}
\begin{ex}
% ID: L12C3B15-52-TN
% Mức: NB
Phát biểu nào sau đây đúng về Xác định phương, chiều lực từ bằng quy tắc bàn tay trái?
\choice
{Lực từ luôn cùng chiều dòng điện trong dây.}
{Có thể bỏ qua phương, chiều và đơn vị của các đại lượng trong mọi trường hợp.}
{Kết quả không phụ thuộc dữ kiện và điều kiện áp dụng.}
{\True Lực từ tác dụng lên đoạn dây có phương vuông góc với cả dòng điện và vectơ cảm ứng từ; chiều xác định bằng quy tắc bàn tay trái.}
\loigiai{
Lực từ tác dụng lên đoạn dây có phương vuông góc với cả dòng điện và vectơ cảm ứng từ; chiều xác định bằng quy tắc bàn tay trái. Vì vậy phương án $D$ đúng; các phương án còn lại trái với kiến thức hoặc điều kiện áp dụng.
}
\end{ex}

% ===== Câu 231 =====
% ID: L12C3B15-53-DS
% Mức: TH
\dangbt{Cân bằng và chuyển động của dây dẫn mang dòng điện trong từ trường}
\begin{ex}
% ID: L12C3B15-53-DS
% Mức: TH
Xét Cân bằng và chuyển động của dây dẫn mang dòng điện trong từ trường (Tình huống 1). Hãy xác định tính đúng, sai của bốn phát biểu sau.
\choiceTF
{\True Dây cân bằng khi tổng vectơ các lực, gồm lực từ và các lực cơ học, bằng vectơ không.}
{Dây cân bằng chỉ cần lực từ khác không, không cần xét trọng lực hay lực căng.}
{\True Khi giải cần xác định đúng phương, chiều, điều kiện áp dụng và đơn vị.}
{Kết quả không phụ thuộc dữ kiện và có thể bỏ qua điều kiện.}
\loigiai{
\begin{itemchoice}
\item (Đúng) Dây cân bằng khi tổng vectơ các lực, gồm lực từ và các lực cơ học, bằng vectơ không. (Vì): ây cân bằng khi tổng vectơ các lực, gồm lực từ và các lực cơ học, bằng vectơ không. b) Sai: Dây cân bằng chỉ cần lực từ khác không, không cần xét trọng lực hay lực căng. c) Đúng: Khi giải cần xác định đúng phương, chiều, điều kiện áp dụng và đơn vị. d) Sai: Kết quả không phụ thuộc dữ kiện và có thể bỏ qua điều kiện. Kiến thức cốt lõi: Dây cân bằng khi tổng vectơ các lực, gồm lực từ và các lực cơ học, bằng vectơ không
\item (Sai) Dây cân bằng chỉ cần lực từ khác không, không cần xét trọng lực hay lực căng.
\item (Đúng) Khi giải cần xác định đúng phương, chiều, điều kiện áp dụng và đơn vị.
\item (Sai) Kết quả không phụ thuộc dữ kiện và có thể bỏ qua điều kiện.
\end{itemchoice}
}
\end{ex}

% ===== Câu 232 =====
% ID: L12C3B15-53-TL
% Mức: VD
\dangbt{Tính lực từ tác dụng lên đoạn dây dẫn mang dòng điện}
\begin{bt}
% ID: L12C3B15-53-TL
% Mức: VD
Nêu một tình huống thực tiễn liên quan đến Tính lực từ tác dụng lên đoạn dây dẫn mang dòng điện và giải thích bằng kiến thức Vật lí.
\loigiai{
Cần nêu được: Độ lớn lực từ lên đoạn dây thẳng dài l mang dòng $I$ trong từ trường đều là $F$ = BIl sin(alpha). Khi vận dụng phải xác định đúng dữ kiện, phương chiều nếu có, điều kiện áp dụng, hệ đơn vị và kiểm tra tính hợp lí. Nhận định “Lực từ được tính bằng $F$ = BI/l và không phụ thuộc góc.” sai vì trái với kiến thức nêu trên.
}
\end{bt}

% ===== Câu 233 =====
% ID: L12C3B15-53-TLN
% Mức: VD
\begin{ex}
% ID: L12C3B15-53-TLN
% Mức: VD
Trong bốn phát biểu về Tính lực từ tác dụng lên đoạn dây dẫn mang dòng điện, có bao nhiêu phát biểu đúng? (Chỉ ghi một số nguyên.) (1) Lực từ được tính bằng $F$ = BI/l và không phụ thuộc góc. (2) Độ lớn lực từ lên đoạn dây thẳng dài l mang dòng $I$ trong từ trường đều là $F$ = BIl sin(alpha). (3) Lực từ được tính bằng $F$ = BI/l và không phụ thuộc góc. (4) Có thể bỏ qua đơn vị của kết quả.
\shortans{2}
\loigiai{
Đối chiếu từng phát biểu với kiến thức: Độ lớn lực từ lên đoạn dây thẳng dài l mang dòng $I$ trong từ trường đều là $F$ = BIl sin(alpha). Số phát biểu đúng là 2.
}
\end{ex}

% ===== Câu 234 =====
% ID: L12C3B15-53-TN
% Mức: NB
\dangbt{Xác định phương, chiều lực từ bằng quy tắc bàn tay trái}
\begin{ex}
% ID: L12C3B15-53-TN
% Mức: NB
Trong nội dung Xác định phương, chiều lực từ bằng quy tắc bàn tay trái?
\choice
{\True Lực từ tác dụng lên đoạn dây có phương vuông góc với cả dòng điện và vectơ cảm ứng từ; chiều xác định bằng quy tắc bàn tay trái.}
{Lực từ luôn cùng chiều dòng điện trong dây.}
{Có thể bỏ qua phương, chiều và đơn vị của các đại lượng trong mọi trường hợp.}
{Kết quả không phụ thuộc dữ kiện và điều kiện áp dụng.}
\loigiai{
Lực từ tác dụng lên đoạn dây có phương vuông góc với cả dòng điện và vectơ cảm ứng từ; chiều xác định bằng quy tắc bàn tay trái. Vì vậy phương án $A$ đúng; các phương án còn lại trái với kiến thức hoặc điều kiện áp dụng.
}
\end{ex}

% ===== Câu 235 =====
% ID: L12C3B15-54-DS
% Mức: TH
\dangbt{Cân bằng và chuyển động của dây dẫn mang dòng điện trong từ trường}
\begin{ex}
% ID: L12C3B15-54-DS
% Mức: TH
Xét Cân bằng và chuyển động của dây dẫn mang dòng điện trong từ trường (Tình huống 2). Hãy xác định tính đúng, sai của bốn phát biểu sau.
\choiceTF
{Dây cân bằng chỉ cần lực từ khác không, không cần xét trọng lực hay lực căng.}
{\True Dây cân bằng khi tổng vectơ các lực, gồm lực từ và các lực cơ học, bằng vectơ không.}
{Kết quả không phụ thuộc dữ kiện và có thể bỏ qua điều kiện.}
{\True Khi giải cần xác định đúng phương, chiều, điều kiện áp dụng và đơn vị.}
\loigiai{
\begin{itemchoice}
\item (Sai) Dây cân bằng chỉ cần lực từ khác không, không cần xét trọng lực hay lực căng. (Vì): Dây cân bằng chỉ cần lực từ khác không, không cần xét trọng lực hay lực căng. b) Đúng: Dây cân bằng khi tổng vectơ các lực, gồm lực từ và các lực cơ học, bằng vectơ không. c) Sai: Kết quả không phụ thuộc dữ kiện và có thể bỏ qua điều kiện. d) Đúng: Khi giải cần xác định đúng phương, chiều, điều kiện áp dụng và đơn vị. Kiến thức cốt lõi: Dây cân bằng khi tổng vectơ các lực, gồm lực từ và các lực cơ học, bằng vectơ không
\item (Đúng) Dây cân bằng khi tổng vectơ các lực, gồm lực từ và các lực cơ học, bằng vectơ không.
\item (Sai) Kết quả không phụ thuộc dữ kiện và có thể bỏ qua điều kiện.
\item (Đúng) Khi giải cần xác định đúng phương, chiều, điều kiện áp dụng và đơn vị.
\end{itemchoice}
}
\end{ex}

% ===== Câu 236 =====
% ID: L12C3B15-54-TL
% Mức: VDC
\dangbt{Tính lực từ tác dụng lên đoạn dây dẫn mang dòng điện}
\begin{bt}
% ID: L12C3B15-54-TL
% Mức: VDC
So sánh nhận định đúng “Độ lớn lực từ lên đoạn dây thẳng dài l mang dòng $I$ trong từ trường đều là $F$ = BIl sin(alpha).” với nhận định sai “Lực từ được tính bằng $F$ = BI/l và không phụ thuộc góc.”; chỉ rõ căn cứ Vật lí.
\loigiai{
Cần nêu được: Độ lớn lực từ lên đoạn dây thẳng dài l mang dòng $I$ trong từ trường đều là $F$ = BIl sin(alpha). Khi vận dụng phải xác định đúng dữ kiện, phương chiều nếu có, điều kiện áp dụng, hệ đơn vị và kiểm tra tính hợp lí. Nhận định “Lực từ được tính bằng $F$ = BI/l và không phụ thuộc góc.” sai vì trái với kiến thức nêu trên.
}
\end{bt}

% ===== Câu 237 =====
% ID: L12C3B15-54-TLN
% Mức: VD
\begin{ex}
% ID: L12C3B15-54-TLN
% Mức: VD
Trong bốn phát biểu về Tính lực từ tác dụng lên đoạn dây dẫn mang dòng điện, có bao nhiêu phát biểu đúng? (Chỉ ghi một số nguyên.) (1) Độ lớn lực từ lên đoạn dây thẳng dài l mang dòng $I$ trong từ trường đều là $F$ = BIl sin(alpha). (2) Cần kiểm tra điều kiện áp dụng. (3) Lực từ được tính bằng $F$ = BI/l và không phụ thuộc góc. (4) Có thể bỏ qua đơn vị của kết quả.
\shortans{2}
\loigiai{
Đối chiếu từng phát biểu với kiến thức: Độ lớn lực từ lên đoạn dây thẳng dài l mang dòng $I$ trong từ trường đều là $F$ = BIl sin(alpha). Số phát biểu đúng là 2.
}
\end{ex}

% ===== Câu 238 =====
% ID: L12C3B15-54-TN
% Mức: NB
\dangbt{Xác định phương, chiều lực từ bằng quy tắc bàn tay trái}
\begin{ex}
% ID: L12C3B15-54-TN
% Mức: NB
Tình huống 3: chọn kết luận chính xác về Xác định phương, chiều lực từ bằng quy tắc bàn tay trái?
\choice
{Kết quả không phụ thuộc dữ kiện và điều kiện áp dụng.}
{\True Lực từ tác dụng lên đoạn dây có phương vuông góc với cả dòng điện và vectơ cảm ứng từ; chiều xác định bằng quy tắc bàn tay trái.}
{Lực từ luôn cùng chiều dòng điện trong dây.}
{Có thể bỏ qua phương, chiều và đơn vị của các đại lượng trong mọi trường hợp.}
\loigiai{
Lực từ tác dụng lên đoạn dây có phương vuông góc với cả dòng điện và vectơ cảm ứng từ; chiều xác định bằng quy tắc bàn tay trái. Vì vậy phương án $B$ đúng; các phương án còn lại trái với kiến thức hoặc điều kiện áp dụng.
}
\end{ex}

% ===== Câu 239 =====
% ID: L12C3B15-55-DS
% Mức: VD
\dangbt{Cân bằng và chuyển động của dây dẫn mang dòng điện trong từ trường}
\begin{ex}
% ID: L12C3B15-55-DS
% Mức: VD
Xét Cân bằng và chuyển động của dây dẫn mang dòng điện trong từ trường (Tình huống 3). Hãy xác định tính đúng, sai của bốn phát biểu sau.
\choiceTF
{\True Khi giải cần xác định đúng phương, chiều, điều kiện áp dụng và đơn vị.}
{Kết quả không phụ thuộc dữ kiện và có thể bỏ qua điều kiện.}
{\True Dây cân bằng khi tổng vectơ các lực, gồm lực từ và các lực cơ học, bằng vectơ không.}
{Dây cân bằng chỉ cần lực từ khác không, không cần xét trọng lực hay lực căng.}
\loigiai{
\begin{itemchoice}
\item (Đúng) Khi giải cần xác định đúng phương, chiều, điều kiện áp dụng và đơn vị. (Vì): Khi giải cần xác định đúng phương, chiều, điều kiện áp dụng và đơn vị. b) Sai: Kết quả không phụ thuộc dữ kiện và có thể bỏ qua điều kiện. c) Đúng: Dây cân bằng khi tổng vectơ các lực, gồm lực từ và các lực cơ học, bằng vectơ không. d) Sai: Dây cân bằng chỉ cần lực từ khác không, không cần xét trọng lực hay lực căng. Kiến thức cốt lõi: Dây cân bằng khi tổng vectơ các lực, gồm lực từ và các lực cơ học, bằng vectơ không
\item (Sai) Kết quả không phụ thuộc dữ kiện và có thể bỏ qua điều kiện.
\item (Đúng) Dây cân bằng khi tổng vectơ các lực, gồm lực từ và các lực cơ học, bằng vectơ không.
\item (Sai) Dây cân bằng chỉ cần lực từ khác không, không cần xét trọng lực hay lực căng.
\end{itemchoice}
}
\end{ex}

% ===== Câu 240 =====
% ID: L12C3B15-55-TL
% Mức: TH
\dangbt{Xác định cảm ứng từ từ phép đo lực từ}
\begin{bt}
% ID: L12C3B15-55-TL
% Mức: TH
Trình bày kiến thức cốt lõi của Xác định cảm ứng từ từ phép đo lực từ và nêu điều kiện áp dụng.
\loigiai{
Cần nêu được: Khi dây vuông góc từ trường, cảm ứng từ được xác định bởi $B$ = F/(Il). Khi vận dụng phải xác định đúng dữ kiện, phương chiều nếu có, điều kiện áp dụng, hệ đơn vị và kiểm tra tính hợp lí. Nhận định “Khi dây vuông góc từ trường, cảm ứng từ bằng FIl.” sai vì trái với kiến thức nêu trên.
}
\end{bt}

% ===== Câu 241 =====
% ID: L12C3B15-55-TLN
% Mức: VDC
\dangbt{Tính lực từ tác dụng lên đoạn dây dẫn mang dòng điện}
\begin{ex}
% ID: L12C3B15-55-TLN
% Mức: VDC
Trong bốn phát biểu về Tính lực từ tác dụng lên đoạn dây dẫn mang dòng điện, có bao nhiêu phát biểu đúng? (Chỉ ghi một số nguyên.) (1) Lực từ được tính bằng $F$ = BI/l và không phụ thuộc góc. (2) Độ lớn lực từ lên đoạn dây thẳng dài l mang dòng $I$ trong từ trường đều là $F$ = BIl sin(alpha). (3) Phải dùng đúng định luật cho tình huống. (4) Độ lớn lực từ lên đoạn dây thẳng dài l mang dòng $I$ trong từ trường đều là $F$ = BIl sin(alpha).
\shortans{3}
\loigiai{
Đối chiếu từng phát biểu với kiến thức: Độ lớn lực từ lên đoạn dây thẳng dài l mang dòng $I$ trong từ trường đều là $F$ = BIl sin(alpha). Số phát biểu đúng là 3.
}
\end{ex}

% ===== Câu 242 =====
% ID: L12C3B15-55-TN
% Mức: TH
\dangbt{Xác định phương, chiều lực từ bằng quy tắc bàn tay trái}
\begin{ex}
% ID: L12C3B15-55-TN
% Mức: TH
Nội dung nào mô tả đúng nhất Xác định phương, chiều lực từ bằng quy tắc bàn tay trái?
\choice
{Có thể bỏ qua phương, chiều và đơn vị của các đại lượng trong mọi trường hợp.}
{Kết quả không phụ thuộc dữ kiện và điều kiện áp dụng.}
{\True Lực từ tác dụng lên đoạn dây có phương vuông góc với cả dòng điện và vectơ cảm ứng từ; chiều xác định bằng quy tắc bàn tay trái.}
{Lực từ luôn cùng chiều dòng điện trong dây.}
\loigiai{
Lực từ tác dụng lên đoạn dây có phương vuông góc với cả dòng điện và vectơ cảm ứng từ; chiều xác định bằng quy tắc bàn tay trái. Vì vậy phương án $C$ đúng; các phương án còn lại trái với kiến thức hoặc điều kiện áp dụng.
}
\end{ex}

% ===== Câu 243 =====
% ID: L12C3B15-56-DS
% Mức: VD
\dangbt{Cân bằng và chuyển động của dây dẫn mang dòng điện trong từ trường}
\begin{ex}
% ID: L12C3B15-56-DS
% Mức: VD
Xét Cân bằng và chuyển động của dây dẫn mang dòng điện trong từ trường (Tình huống 4). Hãy xác định tính đúng, sai của bốn phát biểu sau.
\choiceTF
{Kết quả không phụ thuộc dữ kiện và có thể bỏ qua điều kiện.}
{\True Khi giải cần xác định đúng phương, chiều, điều kiện áp dụng và đơn vị.}
{Dây cân bằng chỉ cần lực từ khác không, không cần xét trọng lực hay lực căng.}
{\True Dây cân bằng khi tổng vectơ các lực, gồm lực từ và các lực cơ học, bằng vectơ không.}
\loigiai{
\begin{itemchoice}
\item (Sai) Kết quả không phụ thuộc dữ kiện và có thể bỏ qua điều kiện. (Vì): Kết quả không phụ thuộc dữ kiện và có thể bỏ qua điều kiện. b) Đúng: Khi giải cần xác định đúng phương, chiều, điều kiện áp dụng và đơn vị. c) Sai: Dây cân bằng chỉ cần lực từ khác không, không cần xét trọng lực hay lực căng. d) Đúng: Dây cân bằng khi tổng vectơ các lực, gồm lực từ và các lực cơ học, bằng vectơ không. Kiến thức cốt lõi: Dây cân bằng khi tổng vectơ các lực, gồm lực từ và các lực cơ học, bằng vectơ không
\item (Đúng) Khi giải cần xác định đúng phương, chiều, điều kiện áp dụng và đơn vị.
\item (Sai) Dây cân bằng chỉ cần lực từ khác không, không cần xét trọng lực hay lực căng.
\item (Đúng) Dây cân bằng khi tổng vectơ các lực, gồm lực từ và các lực cơ học, bằng vectơ không.
\end{itemchoice}
}
\end{ex}

% ===== Câu 244 =====
% ID: L12C3B15-56-TL
% Mức: TH
\dangbt{Xác định cảm ứng từ từ phép đo lực từ}
\begin{bt}
% ID: L12C3B15-56-TL
% Mức: TH
Giải thích vì sao nhận định sau đúng: “Khi dây vuông góc từ trường, cảm ứng từ được xác định bởi $B$ = F/(Il).”
\loigiai{
Cần nêu được: Khi dây vuông góc từ trường, cảm ứng từ được xác định bởi $B$ = F/(Il). Khi vận dụng phải xác định đúng dữ kiện, phương chiều nếu có, điều kiện áp dụng, hệ đơn vị và kiểm tra tính hợp lí. Nhận định “Khi dây vuông góc từ trường, cảm ứng từ bằng FIl.” sai vì trái với kiến thức nêu trên.
}
\end{bt}

% ===== Câu 245 =====
% ID: L12C3B15-56-TLN
% Mức: TH
\begin{ex}
% ID: L12C3B15-56-TLN
% Mức: TH
Trong bốn phát biểu về Xác định cảm ứng từ từ phép đo lực từ, có bao nhiêu phát biểu đúng? (Chỉ ghi một số nguyên.) (1) Khi dây vuông góc từ trường, cảm ứng từ được xác định bởi $B$ = F/(Il). (2) Cần kiểm tra điều kiện áp dụng. (3) Khi dây vuông góc từ trường, cảm ứng từ bằng FIl. (4) Có thể bỏ qua đơn vị của kết quả.
\shortans{2}
\loigiai{
Đối chiếu từng phát biểu với kiến thức: Khi dây vuông góc từ trường, cảm ứng từ được xác định bởi $B$ = F/(Il). Số phát biểu đúng là 2.
}
\end{ex}

% ===== Câu 246 =====
% ID: L12C3B15-56-TN
% Mức: TH
\dangbt{Xác định phương, chiều lực từ bằng quy tắc bàn tay trái}
\begin{ex}
% ID: L12C3B15-56-TN
% Mức: TH
Khi phân tích Xác định phương, chiều lực từ bằng quy tắc bàn tay trái?
\choice
{Lực từ luôn cùng chiều dòng điện trong dây.}
{Có thể bỏ qua phương, chiều và đơn vị của các đại lượng trong mọi trường hợp.}
{Kết quả không phụ thuộc dữ kiện và điều kiện áp dụng.}
{\True Lực từ tác dụng lên đoạn dây có phương vuông góc với cả dòng điện và vectơ cảm ứng từ; chiều xác định bằng quy tắc bàn tay trái.}
\loigiai{
Lực từ tác dụng lên đoạn dây có phương vuông góc với cả dòng điện và vectơ cảm ứng từ; chiều xác định bằng quy tắc bàn tay trái. Vì vậy phương án $D$ đúng; các phương án còn lại trái với kiến thức hoặc điều kiện áp dụng.
}
\end{ex}

% ===== Câu 247 =====
% ID: L12C3B15-57-DS
% Mức: VD
\dangbt{Cân bằng và chuyển động của dây dẫn mang dòng điện trong từ trường}
\begin{ex}
% ID: L12C3B15-57-DS
% Mức: VD
Xét Cân bằng và chuyển động của dây dẫn mang dòng điện trong từ trường (Tình huống 5). Hãy xác định tính đúng, sai của bốn phát biểu sau.
\choiceTF
{\True Dây cân bằng khi tổng vectơ các lực, gồm lực từ và các lực cơ học, bằng vectơ không.}
{Kết quả không phụ thuộc dữ kiện và có thể bỏ qua điều kiện.}
{Dây cân bằng chỉ cần lực từ khác không, không cần xét trọng lực hay lực căng.}
{\True Khi giải cần xác định đúng phương, chiều, điều kiện áp dụng và đơn vị.}
\loigiai{
\begin{itemchoice}
\item (Đúng) Dây cân bằng khi tổng vectơ các lực, gồm lực từ và các lực cơ học, bằng vectơ không. (Vì): ây cân bằng khi tổng vectơ các lực, gồm lực từ và các lực cơ học, bằng vectơ không. b) Sai: Kết quả không phụ thuộc dữ kiện và có thể bỏ qua điều kiện. c) Sai: Dây cân bằng chỉ cần lực từ khác không, không cần xét trọng lực hay lực căng. d) Đúng: Khi giải cần xác định đúng phương, chiều, điều kiện áp dụng và đơn vị. Kiến thức cốt lõi: Dây cân bằng khi tổng vectơ các lực, gồm lực từ và các lực cơ học, bằng vectơ không
\item (Sai) Kết quả không phụ thuộc dữ kiện và có thể bỏ qua điều kiện.
\item (Sai) Dây cân bằng chỉ cần lực từ khác không, không cần xét trọng lực hay lực căng.
\item (Đúng) Khi giải cần xác định đúng phương, chiều, điều kiện áp dụng và đơn vị.
\end{itemchoice}
}
\end{ex}

% ===== Câu 248 =====
% ID: L12C3B15-57-TL
% Mức: VD
\dangbt{Xác định cảm ứng từ từ phép đo lực từ}
\begin{bt}
% ID: L12C3B15-57-TL
% Mức: VD
Phân tích sai lầm trong nhận định: “Khi dây vuông góc từ trường, cảm ứng từ bằng FIl.”
\loigiai{
Cần nêu được: Khi dây vuông góc từ trường, cảm ứng từ được xác định bởi $B$ = F/(Il). Khi vận dụng phải xác định đúng dữ kiện, phương chiều nếu có, điều kiện áp dụng, hệ đơn vị và kiểm tra tính hợp lí. Nhận định “Khi dây vuông góc từ trường, cảm ứng từ bằng FIl.” sai vì trái với kiến thức nêu trên.
}
\end{bt}

% ===== Câu 249 =====
% ID: L12C3B15-57-TLN
% Mức: TH
\begin{ex}
% ID: L12C3B15-57-TLN
% Mức: TH
Trong bốn phát biểu về Xác định cảm ứng từ từ phép đo lực từ, có bao nhiêu phát biểu đúng? (Chỉ ghi một số nguyên.) (1) Khi dây vuông góc từ trường, cảm ứng từ bằng FIl. (2) Khi dây vuông góc từ trường, cảm ứng từ được xác định bởi $B$ = F/(Il). (3) Khi dây vuông góc từ trường, cảm ứng từ bằng FIl. (4) Có thể bỏ qua đơn vị của kết quả.
\shortans{1}
\loigiai{
Đối chiếu từng phát biểu với kiến thức: Khi dây vuông góc từ trường, cảm ứng từ được xác định bởi $B$ = F/(Il). Số phát biểu đúng là 1.
}
\end{ex}

% ===== Câu 250 =====
% ID: L12C3B15-57-TN
% Mức: TH
\dangbt{Xác định phương, chiều lực từ bằng quy tắc bàn tay trái}
\begin{ex}
% ID: L12C3B15-57-TN
% Mức: TH
Một học sinh ôn tập Xác định phương, chiều lực từ bằng quy tắc bàn tay trái?
\choice
{\True Lực từ tác dụng lên đoạn dây có phương vuông góc với cả dòng điện và vectơ cảm ứng từ; chiều xác định bằng quy tắc bàn tay trái.}
{Lực từ luôn cùng chiều dòng điện trong dây.}
{Có thể bỏ qua phương, chiều và đơn vị của các đại lượng trong mọi trường hợp.}
{Kết quả không phụ thuộc dữ kiện và điều kiện áp dụng.}
\loigiai{
Lực từ tác dụng lên đoạn dây có phương vuông góc với cả dòng điện và vectơ cảm ứng từ; chiều xác định bằng quy tắc bàn tay trái. Vì vậy phương án $A$ đúng; các phương án còn lại trái với kiến thức hoặc điều kiện áp dụng.
}
\end{ex}

% ===== Câu 251 =====
% ID: L12C3B15-58-DS
% Mức: TH
\dangbt{Bài toán tổng hợp lực từ gắn với đồ thị, thí nghiệm và thực tiễn}
\begin{ex}
% ID: L12C3B15-58-DS
% Mức: TH
Xét Bài toán tổng hợp lực từ gắn với đồ thị, thí nghiệm và thực tiễn (Tình huống 1). Hãy xác định tính đúng, sai của bốn phát biểu sau.
\choiceTF
{\True Với $B$ và l không đổi, đồ thị $F$ theo $I$ là đường thẳng qua gốc có hệ số góc Bl sin(alpha).}
{Đồ thị $F$ theo $I$ luôn là đường cong và không phụ thuộc B.}
{\True Khi giải cần xác định đúng phương, chiều, điều kiện áp dụng và đơn vị.}
{Kết quả không phụ thuộc dữ kiện và có thể bỏ qua điều kiện.}
\loigiai{
\begin{itemchoice}
\item (Đúng) Với $B$ và l không đổi, đồ thị $F$ theo $I$ là đường thẳng qua gốc có hệ số góc Bl sin(alpha). (Vì): lpha
\item (Sai) Đồ thị $F$ theo $I$ luôn là đường cong và không phụ thuộc B.
\item (Đúng) Khi giải cần xác định đúng phương, chiều, điều kiện áp dụng và đơn vị.
\item (Sai) Kết quả không phụ thuộc dữ kiện và có thể bỏ qua điều kiện.
\end{itemchoice}
}
\end{ex}

% ===== Câu 252 =====
% ID: L12C3B15-58-TL
% Mức: VD
\dangbt{Xác định cảm ứng từ từ phép đo lực từ}
\begin{bt}
% ID: L12C3B15-58-TL
% Mức: VD
Đề xuất các bước giải một bài toán thuộc dạng Xác định cảm ứng từ từ phép đo lực từ.
\loigiai{
Cần nêu được: Khi dây vuông góc từ trường, cảm ứng từ được xác định bởi $B$ = F/(Il). Khi vận dụng phải xác định đúng dữ kiện, phương chiều nếu có, điều kiện áp dụng, hệ đơn vị và kiểm tra tính hợp lí. Nhận định “Khi dây vuông góc từ trường, cảm ứng từ bằng FIl.” sai vì trái với kiến thức nêu trên.
}
\end{bt}

% ===== Câu 253 =====
% ID: L12C3B15-58-TLN
% Mức: TH
\begin{ex}
% ID: L12C3B15-58-TLN
% Mức: TH
Trong bốn phát biểu về Xác định cảm ứng từ từ phép đo lực từ, có bao nhiêu phát biểu đúng? (Chỉ ghi một số nguyên.) (1) Khi dây vuông góc từ trường, cảm ứng từ được xác định bởi $B$ = F/(Il). (2) Cần kiểm tra điều kiện áp dụng. (3) Khi dây vuông góc từ trường, cảm ứng từ bằng FIl. (4) Có thể bỏ qua đơn vị của kết quả.
\shortans{2}
\loigiai{
Đối chiếu từng phát biểu với kiến thức: Khi dây vuông góc từ trường, cảm ứng từ được xác định bởi $B$ = F/(Il). Số phát biểu đúng là 2.
}
\end{ex}

% ===== Câu 254 =====
% ID: L12C3B15-58-TN
% Mức: TH
\dangbt{Xác định phương, chiều lực từ bằng quy tắc bàn tay trái}
\begin{ex}
% ID: L12C3B15-58-TN
% Mức: TH
Chọn phát biểu không trái với kiến thức về Xác định phương, chiều lực từ bằng quy tắc bàn tay trái?
\choice
{Kết quả không phụ thuộc dữ kiện và điều kiện áp dụng.}
{\True Lực từ tác dụng lên đoạn dây có phương vuông góc với cả dòng điện và vectơ cảm ứng từ; chiều xác định bằng quy tắc bàn tay trái.}
{Lực từ luôn cùng chiều dòng điện trong dây.}
{Có thể bỏ qua phương, chiều và đơn vị của các đại lượng trong mọi trường hợp.}
\loigiai{
Lực từ tác dụng lên đoạn dây có phương vuông góc với cả dòng điện và vectơ cảm ứng từ; chiều xác định bằng quy tắc bàn tay trái. Vì vậy phương án $B$ đúng; các phương án còn lại trái với kiến thức hoặc điều kiện áp dụng.
}
\end{ex}

% ===== Câu 255 =====
% ID: L12C3B15-59-DS
% Mức: TH
\dangbt{Bài toán tổng hợp lực từ gắn với đồ thị, thí nghiệm và thực tiễn}
\begin{ex}
% ID: L12C3B15-59-DS
% Mức: TH
Xét Bài toán tổng hợp lực từ gắn với đồ thị, thí nghiệm và thực tiễn (Tình huống 2). Hãy xác định tính đúng, sai của bốn phát biểu sau.
\choiceTF
{Đồ thị $F$ theo $I$ luôn là đường cong và không phụ thuộc B.}
{\True Với $B$ và l không đổi, đồ thị $F$ theo $I$ là đường thẳng qua gốc có hệ số góc Bl sin(alpha).}
{Kết quả không phụ thuộc dữ kiện và có thể bỏ qua điều kiện.}
{\True Khi giải cần xác định đúng phương, chiều, điều kiện áp dụng và đơn vị.}
\loigiai{
\begin{itemchoice}
\item (Sai) Đồ thị $F$ theo $I$ luôn là đường cong và không phụ thuộc B. (Vì): lpha
\item (Đúng) Với $B$ và l không đổi, đồ thị $F$ theo $I$ là đường thẳng qua gốc có hệ số góc Bl sin(alpha).
\item (Sai) Kết quả không phụ thuộc dữ kiện và có thể bỏ qua điều kiện.
\item (Đúng) Khi giải cần xác định đúng phương, chiều, điều kiện áp dụng và đơn vị.
\end{itemchoice}
}
\end{ex}

% ===== Câu 256 =====
% ID: L12C3B15-59-TL
% Mức: VD
\dangbt{Xác định cảm ứng từ từ phép đo lực từ}
\begin{bt}
% ID: L12C3B15-59-TL
% Mức: VD
Nêu một tình huống thực tiễn liên quan đến Xác định cảm ứng từ từ phép đo lực từ và giải thích bằng kiến thức Vật lí.
\loigiai{
Cần nêu được: Khi dây vuông góc từ trường, cảm ứng từ được xác định bởi $B$ = F/(Il). Khi vận dụng phải xác định đúng dữ kiện, phương chiều nếu có, điều kiện áp dụng, hệ đơn vị và kiểm tra tính hợp lí. Nhận định “Khi dây vuông góc từ trường, cảm ứng từ bằng FIl.” sai vì trái với kiến thức nêu trên.
}
\end{bt}

% ===== Câu 257 =====
% ID: L12C3B15-59-TLN
% Mức: VD
\begin{ex}
% ID: L12C3B15-59-TLN
% Mức: VD
Trong bốn phát biểu về Xác định cảm ứng từ từ phép đo lực từ, có bao nhiêu phát biểu đúng? (Chỉ ghi một số nguyên.) (1) Khi dây vuông góc từ trường, cảm ứng từ bằng FIl. (2) Khi dây vuông góc từ trường, cảm ứng từ được xác định bởi $B$ = F/(Il). (3) Khi dây vuông góc từ trường, cảm ứng từ bằng FIl. (4) Có thể bỏ qua đơn vị của kết quả.
\shortans{2}
\loigiai{
Đối chiếu từng phát biểu với kiến thức: Khi dây vuông góc từ trường, cảm ứng từ được xác định bởi $B$ = F/(Il). Số phát biểu đúng là 2.
}
\end{ex}

% ===== Câu 258 =====
% ID: L12C3B15-59-TN
% Mức: VD
\dangbt{Xác định phương, chiều lực từ bằng quy tắc bàn tay trái}
\begin{ex}
% ID: L12C3B15-59-TN
% Mức: VD
Cơ sở Vật lí đúng dùng để giải Xác định phương, chiều lực từ bằng quy tắc bàn tay trái?
\choice
{Có thể bỏ qua phương, chiều và đơn vị của các đại lượng trong mọi trường hợp.}
{Kết quả không phụ thuộc dữ kiện và điều kiện áp dụng.}
{\True Lực từ tác dụng lên đoạn dây có phương vuông góc với cả dòng điện và vectơ cảm ứng từ; chiều xác định bằng quy tắc bàn tay trái.}
{Lực từ luôn cùng chiều dòng điện trong dây.}
\loigiai{
Lực từ tác dụng lên đoạn dây có phương vuông góc với cả dòng điện và vectơ cảm ứng từ; chiều xác định bằng quy tắc bàn tay trái. Vì vậy phương án $C$ đúng; các phương án còn lại trái với kiến thức hoặc điều kiện áp dụng.
}
\end{ex}

% ===== Câu 259 =====
% ID: L12C3B15-60-DS
% Mức: VD
\dangbt{Bài toán tổng hợp lực từ gắn với đồ thị, thí nghiệm và thực tiễn}
\begin{ex}
% ID: L12C3B15-60-DS
% Mức: VD
Xét Bài toán tổng hợp lực từ gắn với đồ thị, thí nghiệm và thực tiễn (Tình huống 3). Hãy xác định tính đúng, sai của bốn phát biểu sau.
\choiceTF
{\True Khi giải cần xác định đúng phương, chiều, điều kiện áp dụng và đơn vị.}
{Kết quả không phụ thuộc dữ kiện và có thể bỏ qua điều kiện.}
{\True Với $B$ và l không đổi, đồ thị $F$ theo $I$ là đường thẳng qua gốc có hệ số góc Bl sin(alpha).}
{Đồ thị $F$ theo $I$ luôn là đường cong và không phụ thuộc B.}
\loigiai{
\begin{itemchoice}
\item (Đúng) Khi giải cần xác định đúng phương, chiều, điều kiện áp dụng và đơn vị. (Vì): lpha
\item (Sai) Kết quả không phụ thuộc dữ kiện và có thể bỏ qua điều kiện.
\item (Đúng) Với $B$ và l không đổi, đồ thị $F$ theo $I$ là đường thẳng qua gốc có hệ số góc Bl sin(alpha).
\item (Sai) Đồ thị $F$ theo $I$ luôn là đường cong và không phụ thuộc B.
\end{itemchoice}
}
\end{ex}

% ===== Câu 260 =====
% ID: L12C3B15-60-TL
% Mức: VDC
\dangbt{Xác định cảm ứng từ từ phép đo lực từ}
\begin{bt}
% ID: L12C3B15-60-TL
% Mức: VDC
So sánh nhận định đúng “Khi dây vuông góc từ trường, cảm ứng từ được xác định bởi $B$ = F/(Il).” với nhận định sai “Khi dây vuông góc từ trường, cảm ứng từ bằng FIl.”; chỉ rõ căn cứ Vật lí.
\loigiai{
Cần nêu được: Khi dây vuông góc từ trường, cảm ứng từ được xác định bởi $B$ = F/(Il). Khi vận dụng phải xác định đúng dữ kiện, phương chiều nếu có, điều kiện áp dụng, hệ đơn vị và kiểm tra tính hợp lí. Nhận định “Khi dây vuông góc từ trường, cảm ứng từ bằng FIl.” sai vì trái với kiến thức nêu trên.
}
\end{bt}

% ===== Câu 261 =====
% ID: L12C3B15-60-TLN
% Mức: VD
\begin{ex}
% ID: L12C3B15-60-TLN
% Mức: VD
Trong bốn phát biểu về Xác định cảm ứng từ từ phép đo lực từ, có bao nhiêu phát biểu đúng? (Chỉ ghi một số nguyên.) (1) Khi dây vuông góc từ trường, cảm ứng từ được xác định bởi $B$ = F/(Il). (2) Cần kiểm tra điều kiện áp dụng. (3) Khi dây vuông góc từ trường, cảm ứng từ bằng FIl. (4) Có thể bỏ qua đơn vị của kết quả.
\shortans{2}
\loigiai{
Đối chiếu từng phát biểu với kiến thức: Khi dây vuông góc từ trường, cảm ứng từ được xác định bởi $B$ = F/(Il). Số phát biểu đúng là 2.
}
\end{ex}

% ===== Câu 262 =====
% ID: L12C3B15-60-TN
% Mức: VD
\dangbt{Xác định phương, chiều lực từ bằng quy tắc bàn tay trái}
\begin{ex}
% ID: L12C3B15-60-TN
% Mức: VD
Trong một bài thực hành về Xác định phương, chiều lực từ bằng quy tắc bàn tay trái?
\choice
{Lực từ luôn cùng chiều dòng điện trong dây.}
{Có thể bỏ qua phương, chiều và đơn vị của các đại lượng trong mọi trường hợp.}
{Kết quả không phụ thuộc dữ kiện và điều kiện áp dụng.}
{\True Lực từ tác dụng lên đoạn dây có phương vuông góc với cả dòng điện và vectơ cảm ứng từ; chiều xác định bằng quy tắc bàn tay trái.}
\loigiai{
Lực từ tác dụng lên đoạn dây có phương vuông góc với cả dòng điện và vectơ cảm ứng từ; chiều xác định bằng quy tắc bàn tay trái. Vì vậy phương án $D$ đúng; các phương án còn lại trái với kiến thức hoặc điều kiện áp dụng.
}
\end{ex}

% ===== Câu 263 =====
% ID: L12C3B15-61-DS
% Mức: VD
\dangbt{Bài toán tổng hợp lực từ gắn với đồ thị, thí nghiệm và thực tiễn}
\begin{ex}
% ID: L12C3B15-61-DS
% Mức: VD
Xét Bài toán tổng hợp lực từ gắn với đồ thị, thí nghiệm và thực tiễn (Tình huống 4). Hãy xác định tính đúng, sai của bốn phát biểu sau.
\choiceTF
{Kết quả không phụ thuộc dữ kiện và có thể bỏ qua điều kiện.}
{\True Khi giải cần xác định đúng phương, chiều, điều kiện áp dụng và đơn vị.}
{Đồ thị $F$ theo $I$ luôn là đường cong và không phụ thuộc B.}
{\True Với $B$ và l không đổi, đồ thị $F$ theo $I$ là đường thẳng qua gốc có hệ số góc Bl sin(alpha).}
\loigiai{
\begin{itemchoice}
\item (Sai) Kết quả không phụ thuộc dữ kiện và có thể bỏ qua điều kiện. (Vì): lpha
\item (Đúng) Khi giải cần xác định đúng phương, chiều, điều kiện áp dụng và đơn vị.
\item (Sai) Đồ thị $F$ theo $I$ luôn là đường cong và không phụ thuộc B.
\item (Đúng) Với $B$ và l không đổi, đồ thị $F$ theo $I$ là đường thẳng qua gốc có hệ số góc Bl sin(alpha).
\end{itemchoice}
}
\end{ex}

% ===== Câu 264 =====
% ID: L12C3B15-61-TL
% Mức: TH
\dangbt{Cân bằng và chuyển động của dây dẫn mang dòng điện trong từ trường}
\begin{bt}
% ID: L12C3B15-61-TL
% Mức: TH
Trình bày kiến thức cốt lõi của Cân bằng và chuyển động của dây dẫn mang dòng điện trong từ trường và nêu điều kiện áp dụng.
\loigiai{
Cần nêu được: Dây cân bằng khi tổng vectơ các lực, gồm lực từ và các lực cơ học, bằng vectơ không. Khi vận dụng phải xác định đúng dữ kiện, phương chiều nếu có, điều kiện áp dụng, hệ đơn vị và kiểm tra tính hợp lí. Nhận định “Dây cân bằng chỉ cần lực từ khác không, không cần xét trọng lực hay lực căng.” sai vì trái với kiến thức nêu trên.
}
\end{bt}

% ===== Câu 265 =====
% ID: L12C3B15-61-TLN
% Mức: VDC
\dangbt{Xác định cảm ứng từ từ phép đo lực từ}
\begin{ex}
% ID: L12C3B15-61-TLN
% Mức: VDC
Trong bốn phát biểu về Xác định cảm ứng từ từ phép đo lực từ, có bao nhiêu phát biểu đúng? (Chỉ ghi một số nguyên.) (1) Khi dây vuông góc từ trường, cảm ứng từ bằng FIl. (2) Khi dây vuông góc từ trường, cảm ứng từ được xác định bởi $B$ = F/(Il). (3) Phải dùng đúng định luật cho tình huống. (4) Khi dây vuông góc từ trường, cảm ứng từ được xác định bởi $B$ = F/(Il).
\shortans{3}
\loigiai{
Đối chiếu từng phát biểu với kiến thức: Khi dây vuông góc từ trường, cảm ứng từ được xác định bởi $B$ = F/(Il). Số phát biểu đúng là 3.
}
\end{ex}

% ===== Câu 266 =====
% ID: L12C3B15-61-TN
% Mức: VD
\dangbt{Xác định phương, chiều lực từ bằng quy tắc bàn tay trái}
\begin{ex}
% ID: L12C3B15-61-TN
% Mức: VD
Kết luận đúng khi vận dụng Xác định phương, chiều lực từ bằng quy tắc bàn tay trái?
\choice
{\True Lực từ tác dụng lên đoạn dây có phương vuông góc với cả dòng điện và vectơ cảm ứng từ; chiều xác định bằng quy tắc bàn tay trái.}
{Lực từ luôn cùng chiều dòng điện trong dây.}
{Có thể bỏ qua phương, chiều và đơn vị của các đại lượng trong mọi trường hợp.}
{Kết quả không phụ thuộc dữ kiện và điều kiện áp dụng.}
\loigiai{
Lực từ tác dụng lên đoạn dây có phương vuông góc với cả dòng điện và vectơ cảm ứng từ; chiều xác định bằng quy tắc bàn tay trái. Vì vậy phương án $A$ đúng; các phương án còn lại trái với kiến thức hoặc điều kiện áp dụng.
}
\end{ex}

% ===== Câu 267 =====
% ID: L12C3B15-62-DS
% Mức: VD
\dangbt{Bài toán tổng hợp lực từ gắn với đồ thị, thí nghiệm và thực tiễn}
\begin{ex}
% ID: L12C3B15-62-DS
% Mức: VD
Xét Bài toán tổng hợp lực từ gắn với đồ thị, thí nghiệm và thực tiễn (Tình huống 5). Hãy xác định tính đúng, sai của bốn phát biểu sau.
\choiceTF
{\True Với $B$ và l không đổi, đồ thị $F$ theo $I$ là đường thẳng qua gốc có hệ số góc Bl sin(alpha).}
{Kết quả không phụ thuộc dữ kiện và có thể bỏ qua điều kiện.}
{Đồ thị $F$ theo $I$ luôn là đường cong và không phụ thuộc B.}
{\True Khi giải cần xác định đúng phương, chiều, điều kiện áp dụng và đơn vị.}
\loigiai{
\begin{itemchoice}
\item (Đúng) Với $B$ và l không đổi, đồ thị $F$ theo $I$ là đường thẳng qua gốc có hệ số góc Bl sin(alpha). (Vì): lpha
\item (Sai) Kết quả không phụ thuộc dữ kiện và có thể bỏ qua điều kiện.
\item (Sai) Đồ thị $F$ theo $I$ luôn là đường cong và không phụ thuộc B.
\item (Đúng) Khi giải cần xác định đúng phương, chiều, điều kiện áp dụng và đơn vị.
\end{itemchoice}
}
\end{ex}

% ===== Câu 268 =====
% ID: L12C3B15-62-TL
% Mức: TH
\dangbt{Cân bằng và chuyển động của dây dẫn mang dòng điện trong từ trường}
\begin{bt}
% ID: L12C3B15-62-TL
% Mức: TH
Giải thích vì sao nhận định sau đúng: “Dây cân bằng khi tổng vectơ các lực, gồm lực từ và các lực cơ học, bằng vectơ không.”
\loigiai{
Cần nêu được: Dây cân bằng khi tổng vectơ các lực, gồm lực từ và các lực cơ học, bằng vectơ không. Khi vận dụng phải xác định đúng dữ kiện, phương chiều nếu có, điều kiện áp dụng, hệ đơn vị và kiểm tra tính hợp lí. Nhận định “Dây cân bằng chỉ cần lực từ khác không, không cần xét trọng lực hay lực căng.” sai vì trái với kiến thức nêu trên.
}
\end{bt}

% ===== Câu 269 =====
% ID: L12C3B15-62-TLN
% Mức: TH
\begin{ex}
% ID: L12C3B15-62-TLN
% Mức: TH
Trong bốn phát biểu về Cân bằng và chuyển động của dây dẫn mang dòng điện trong từ trường, có bao nhiêu phát biểu đúng? (Chỉ ghi một số nguyên.) (1) Dây cân bằng khi tổng vectơ các lực, gồm lực từ và các lực cơ học, bằng vectơ không. (2) Cần kiểm tra điều kiện áp dụng. (3) Dây cân bằng chỉ cần lực từ khác không, không cần xét trọng lực hay lực căng. (4) Có thể bỏ qua đơn vị của kết quả.
\shortans{2}
\loigiai{
Đối chiếu từng phát biểu với kiến thức: Dây cân bằng khi tổng vectơ các lực, gồm lực từ và các lực cơ học, bằng vectơ không. Số phát biểu đúng là 2.
}
\end{ex}

% ===== Câu 270 =====
% ID: L12C3B15-62-TN
% Mức: NB
\dangbt{Nhận biết lực từ và đặc trưng của vectơ cảm ứng từ}
\begin{ex}
% ID: L12C3B15-62-TN
% Mức: NB
Phát biểu nào sau đây đúng về Nhận biết lực từ và đặc trưng của vectơ cảm ứng từ?
\choice
{Cảm ứng từ là đại lượng vô hướng và có đơn vị niutơn.}
{Có thể bỏ qua phương, chiều và đơn vị của các đại lượng trong mọi trường hợp.}
{Kết quả không phụ thuộc dữ kiện và điều kiện áp dụng.}
{\True Vectơ cảm ứng từ tại một điểm có hướng trùng với hướng của từ trường tại điểm đó; đơn vị là tesla (T).}
\loigiai{
Vectơ cảm ứng từ tại một điểm có hướng trùng với hướng của từ trường tại điểm đó; đơn vị là tesla (T). Vì vậy phương án $D$ đúng; các phương án còn lại trái với kiến thức hoặc điều kiện áp dụng.
}
\end{ex}

% ===== Câu 271 =====
% ID: L12C3B15-63-TL
% Mức: VD
\dangbt{Cân bằng và chuyển động của dây dẫn mang dòng điện trong từ trường}
\begin{bt}
% ID: L12C3B15-63-TL
% Mức: VD
Phân tích sai lầm trong nhận định: “Dây cân bằng chỉ cần lực từ khác không, không cần xét trọng lực hay lực căng.”
\loigiai{
Cần nêu được: Dây cân bằng khi tổng vectơ các lực, gồm lực từ và các lực cơ học, bằng vectơ không. Khi vận dụng phải xác định đúng dữ kiện, phương chiều nếu có, điều kiện áp dụng, hệ đơn vị và kiểm tra tính hợp lí. Nhận định “Dây cân bằng chỉ cần lực từ khác không, không cần xét trọng lực hay lực căng.” sai vì trái với kiến thức nêu trên.
}
\end{bt}

% ===== Câu 272 =====
% ID: L12C3B15-63-TLN
% Mức: TH
\begin{ex}
% ID: L12C3B15-63-TLN
% Mức: TH
Trong bốn phát biểu về Cân bằng và chuyển động của dây dẫn mang dòng điện trong từ trường, có bao nhiêu phát biểu đúng? (Chỉ ghi một số nguyên.) (1) Dây cân bằng chỉ cần lực từ khác không, không cần xét trọng lực hay lực căng. (2) Dây cân bằng khi tổng vectơ các lực, gồm lực từ và các lực cơ học, bằng vectơ không. (3) Dây cân bằng chỉ cần lực từ khác không, không cần xét trọng lực hay lực căng. (4) Có thể bỏ qua đơn vị của kết quả.
\shortans{1}
\loigiai{
Đối chiếu từng phát biểu với kiến thức: Dây cân bằng khi tổng vectơ các lực, gồm lực từ và các lực cơ học, bằng vectơ không. Số phát biểu đúng là 1.
}
\end{ex}

% ===== Câu 273 =====
% ID: L12C3B15-63-TN
% Mức: NB
\dangbt{Nhận biết lực từ và đặc trưng của vectơ cảm ứng từ}
\begin{ex}
% ID: L12C3B15-63-TN
% Mức: NB
Trong nội dung Nhận biết lực từ và đặc trưng của vectơ cảm ứng từ?
\choice
{\True Vectơ cảm ứng từ tại một điểm có hướng trùng với hướng của từ trường tại điểm đó; đơn vị là tesla (T).}
{Cảm ứng từ là đại lượng vô hướng và có đơn vị niutơn.}
{Có thể bỏ qua phương, chiều và đơn vị của các đại lượng trong mọi trường hợp.}
{Kết quả không phụ thuộc dữ kiện và điều kiện áp dụng.}
\loigiai{
Vectơ cảm ứng từ tại một điểm có hướng trùng với hướng của từ trường tại điểm đó; đơn vị là tesla (T). Vì vậy phương án $A$ đúng; các phương án còn lại trái với kiến thức hoặc điều kiện áp dụng.
}
\end{ex}

% ===== Câu 274 =====
% ID: L12C3B15-64-TL
% Mức: VD
\dangbt{Cân bằng và chuyển động của dây dẫn mang dòng điện trong từ trường}
\begin{bt}
% ID: L12C3B15-64-TL
% Mức: VD
Đề xuất các bước giải một bài toán thuộc dạng Cân bằng và chuyển động của dây dẫn mang dòng điện trong từ trường.
\loigiai{
Cần nêu được: Dây cân bằng khi tổng vectơ các lực, gồm lực từ và các lực cơ học, bằng vectơ không. Khi vận dụng phải xác định đúng dữ kiện, phương chiều nếu có, điều kiện áp dụng, hệ đơn vị và kiểm tra tính hợp lí. Nhận định “Dây cân bằng chỉ cần lực từ khác không, không cần xét trọng lực hay lực căng.” sai vì trái với kiến thức nêu trên.
}
\end{bt}

% ===== Câu 275 =====
% ID: L12C3B15-64-TLN
% Mức: TH
\begin{ex}
% ID: L12C3B15-64-TLN
% Mức: TH
Trong bốn phát biểu về Cân bằng và chuyển động của dây dẫn mang dòng điện trong từ trường, có bao nhiêu phát biểu đúng? (Chỉ ghi một số nguyên.) (1) Dây cân bằng khi tổng vectơ các lực, gồm lực từ và các lực cơ học, bằng vectơ không. (2) Cần kiểm tra điều kiện áp dụng. (3) Dây cân bằng chỉ cần lực từ khác không, không cần xét trọng lực hay lực căng. (4) Có thể bỏ qua đơn vị của kết quả.
\shortans{2}
\loigiai{
Đối chiếu từng phát biểu với kiến thức: Dây cân bằng khi tổng vectơ các lực, gồm lực từ và các lực cơ học, bằng vectơ không. Số phát biểu đúng là 2.
}
\end{ex}

% ===== Câu 276 =====
% ID: L12C3B15-64-TN
% Mức: NB
\dangbt{Nhận biết lực từ và đặc trưng của vectơ cảm ứng từ}
\begin{ex}
% ID: L12C3B15-64-TN
% Mức: NB
Tình huống 3: chọn kết luận chính xác về Nhận biết lực từ và đặc trưng của vectơ cảm ứng từ?
\choice
{Kết quả không phụ thuộc dữ kiện và điều kiện áp dụng.}
{\True Vectơ cảm ứng từ tại một điểm có hướng trùng với hướng của từ trường tại điểm đó; đơn vị là tesla (T).}
{Cảm ứng từ là đại lượng vô hướng và có đơn vị niutơn.}
{Có thể bỏ qua phương, chiều và đơn vị của các đại lượng trong mọi trường hợp.}
\loigiai{
Vectơ cảm ứng từ tại một điểm có hướng trùng với hướng của từ trường tại điểm đó; đơn vị là tesla (T). Vì vậy phương án $B$ đúng; các phương án còn lại trái với kiến thức hoặc điều kiện áp dụng.
}
\end{ex}

% ===== Câu 277 =====
% ID: L12C3B15-65-TL
% Mức: VD
\dangbt{Cân bằng và chuyển động của dây dẫn mang dòng điện trong từ trường}
\begin{bt}
% ID: L12C3B15-65-TL
% Mức: VD
Nêu một tình huống thực tiễn liên quan đến Cân bằng và chuyển động của dây dẫn mang dòng điện trong từ trường và giải thích bằng kiến thức Vật lí.
\loigiai{
Cần nêu được: Dây cân bằng khi tổng vectơ các lực, gồm lực từ và các lực cơ học, bằng vectơ không. Khi vận dụng phải xác định đúng dữ kiện, phương chiều nếu có, điều kiện áp dụng, hệ đơn vị và kiểm tra tính hợp lí. Nhận định “Dây cân bằng chỉ cần lực từ khác không, không cần xét trọng lực hay lực căng.” sai vì trái với kiến thức nêu trên.
}
\end{bt}

% ===== Câu 278 =====
% ID: L12C3B15-65-TLN
% Mức: VD
\begin{ex}
% ID: L12C3B15-65-TLN
% Mức: VD
Trong bốn phát biểu về Cân bằng và chuyển động của dây dẫn mang dòng điện trong từ trường, có bao nhiêu phát biểu đúng? (Chỉ ghi một số nguyên.) (1) Dây cân bằng chỉ cần lực từ khác không, không cần xét trọng lực hay lực căng. (2) Dây cân bằng khi tổng vectơ các lực, gồm lực từ và các lực cơ học, bằng vectơ không. (3) Dây cân bằng chỉ cần lực từ khác không, không cần xét trọng lực hay lực căng. (4) Có thể bỏ qua đơn vị của kết quả.
\shortans{2}
\loigiai{
Đối chiếu từng phát biểu với kiến thức: Dây cân bằng khi tổng vectơ các lực, gồm lực từ và các lực cơ học, bằng vectơ không. Số phát biểu đúng là 2.
}
\end{ex}

% ===== Câu 279 =====
% ID: L12C3B15-65-TN
% Mức: TH
\dangbt{Nhận biết lực từ và đặc trưng của vectơ cảm ứng từ}
\begin{ex}
% ID: L12C3B15-65-TN
% Mức: TH
Nội dung nào mô tả đúng nhất Nhận biết lực từ và đặc trưng của vectơ cảm ứng từ?
\choice
{Có thể bỏ qua phương, chiều và đơn vị của các đại lượng trong mọi trường hợp.}
{Kết quả không phụ thuộc dữ kiện và điều kiện áp dụng.}
{\True Vectơ cảm ứng từ tại một điểm có hướng trùng với hướng của từ trường tại điểm đó; đơn vị là tesla (T).}
{Cảm ứng từ là đại lượng vô hướng và có đơn vị niutơn.}
\loigiai{
Vectơ cảm ứng từ tại một điểm có hướng trùng với hướng của từ trường tại điểm đó; đơn vị là tesla (T). Vì vậy phương án $C$ đúng; các phương án còn lại trái với kiến thức hoặc điều kiện áp dụng.
}
\end{ex}

% ===== Câu 280 =====
% ID: L12C3B15-66-TL
% Mức: VDC
\dangbt{Cân bằng và chuyển động của dây dẫn mang dòng điện trong từ trường}
\begin{bt}
% ID: L12C3B15-66-TL
% Mức: VDC
So sánh nhận định đúng “Dây cân bằng khi tổng vectơ các lực, gồm lực từ và các lực cơ học, bằng vectơ không.” với nhận định sai “Dây cân bằng chỉ cần lực từ khác không, không cần xét trọng lực hay lực căng.”; chỉ rõ căn cứ Vật lí.
\loigiai{
Cần nêu được: Dây cân bằng khi tổng vectơ các lực, gồm lực từ và các lực cơ học, bằng vectơ không. Khi vận dụng phải xác định đúng dữ kiện, phương chiều nếu có, điều kiện áp dụng, hệ đơn vị và kiểm tra tính hợp lí. Nhận định “Dây cân bằng chỉ cần lực từ khác không, không cần xét trọng lực hay lực căng.” sai vì trái với kiến thức nêu trên.
}
\end{bt}

% ===== Câu 281 =====
% ID: L12C3B15-66-TLN
% Mức: VD
\begin{ex}
% ID: L12C3B15-66-TLN
% Mức: VD
Trong bốn phát biểu về Cân bằng và chuyển động của dây dẫn mang dòng điện trong từ trường, có bao nhiêu phát biểu đúng? (Chỉ ghi một số nguyên.) (1) Dây cân bằng khi tổng vectơ các lực, gồm lực từ và các lực cơ học, bằng vectơ không. (2) Cần kiểm tra điều kiện áp dụng. (3) Dây cân bằng chỉ cần lực từ khác không, không cần xét trọng lực hay lực căng. (4) Có thể bỏ qua đơn vị của kết quả.
\shortans{2}
\loigiai{
Đối chiếu từng phát biểu với kiến thức: Dây cân bằng khi tổng vectơ các lực, gồm lực từ và các lực cơ học, bằng vectơ không. Số phát biểu đúng là 2.
}
\end{ex}

% ===== Câu 282 =====
% ID: L12C3B15-66-TN
% Mức: TH
\dangbt{Nhận biết lực từ và đặc trưng của vectơ cảm ứng từ}
\begin{ex}
% ID: L12C3B15-66-TN
% Mức: TH
Khi phân tích Nhận biết lực từ và đặc trưng của vectơ cảm ứng từ?
\choice
{Cảm ứng từ là đại lượng vô hướng và có đơn vị niutơn.}
{Có thể bỏ qua phương, chiều và đơn vị của các đại lượng trong mọi trường hợp.}
{Kết quả không phụ thuộc dữ kiện và điều kiện áp dụng.}
{\True Vectơ cảm ứng từ tại một điểm có hướng trùng với hướng của từ trường tại điểm đó; đơn vị là tesla (T).}
\loigiai{
Vectơ cảm ứng từ tại một điểm có hướng trùng với hướng của từ trường tại điểm đó; đơn vị là tesla (T). Vì vậy phương án $D$ đúng; các phương án còn lại trái với kiến thức hoặc điều kiện áp dụng.
}
\end{ex}

% ===== Câu 283 =====
% ID: L12C3B15-67-TL
% Mức: TH
\dangbt{Bài toán tổng hợp lực từ gắn với đồ thị, thí nghiệm và thực tiễn}
\begin{bt}
% ID: L12C3B15-67-TL
% Mức: TH
Trình bày kiến thức cốt lõi của Bài toán tổng hợp lực từ gắn với đồ thị, thí nghiệm và thực tiễn và nêu điều kiện áp dụng.
\loigiai{
Cần nêu được: Với $B$ và l không đổi, đồ thị $F$ theo $I$ là đường thẳng qua gốc có hệ số góc Bl sin(alpha). Khi vận dụng phải xác định đúng dữ kiện, phương chiều nếu có, điều kiện áp dụng, hệ đơn vị và kiểm tra tính hợp lí. Nhận định “Đồ thị $F$ theo $I$ luôn là đường cong và không phụ thuộc B.” sai vì trái với kiến thức nêu trên.
}
\end{bt}

% ===== Câu 284 =====
% ID: L12C3B15-67-TLN
% Mức: VDC
\dangbt{Cân bằng và chuyển động của dây dẫn mang dòng điện trong từ trường}
\begin{ex}
% ID: L12C3B15-67-TLN
% Mức: VDC
Trong bốn phát biểu về Cân bằng và chuyển động của dây dẫn mang dòng điện trong từ trường, có bao nhiêu phát biểu đúng? (Chỉ ghi một số nguyên.) (1) Dây cân bằng chỉ cần lực từ khác không, không cần xét trọng lực hay lực căng. (2) Dây cân bằng khi tổng vectơ các lực, gồm lực từ và các lực cơ học, bằng vectơ không. (3) Phải dùng đúng định luật cho tình huống. (4) Dây cân bằng khi tổng vectơ các lực, gồm lực từ và các lực cơ học, bằng vectơ không.
\shortans{3}
\loigiai{
Đối chiếu từng phát biểu với kiến thức: Dây cân bằng khi tổng vectơ các lực, gồm lực từ và các lực cơ học, bằng vectơ không. Số phát biểu đúng là 3.
}
\end{ex}

% ===== Câu 285 =====
% ID: L12C3B15-67-TN
% Mức: TH
\dangbt{Nhận biết lực từ và đặc trưng của vectơ cảm ứng từ}
\begin{ex}
% ID: L12C3B15-67-TN
% Mức: TH
Một học sinh ôn tập Nhận biết lực từ và đặc trưng của vectơ cảm ứng từ?
\choice
{\True Vectơ cảm ứng từ tại một điểm có hướng trùng với hướng của từ trường tại điểm đó; đơn vị là tesla (T).}
{Cảm ứng từ là đại lượng vô hướng và có đơn vị niutơn.}
{Có thể bỏ qua phương, chiều và đơn vị của các đại lượng trong mọi trường hợp.}
{Kết quả không phụ thuộc dữ kiện và điều kiện áp dụng.}
\loigiai{
Vectơ cảm ứng từ tại một điểm có hướng trùng với hướng của từ trường tại điểm đó; đơn vị là tesla (T). Vì vậy phương án $A$ đúng; các phương án còn lại trái với kiến thức hoặc điều kiện áp dụng.
}
\end{ex}

% ===== Câu 286 =====
% ID: L12C3B15-68-TL
% Mức: TH
\dangbt{Bài toán tổng hợp lực từ gắn với đồ thị, thí nghiệm và thực tiễn}
\begin{bt}
% ID: L12C3B15-68-TL
% Mức: TH
Giải thích vì sao nhận định sau đúng: “Với $B$ và l không đổi, đồ thị $F$ theo $I$ là đường thẳng qua gốc có hệ số góc Bl sin(alpha).”
\loigiai{
Cần nêu được: Với $B$ và l không đổi, đồ thị $F$ theo $I$ là đường thẳng qua gốc có hệ số góc Bl sin(alpha). Khi vận dụng phải xác định đúng dữ kiện, phương chiều nếu có, điều kiện áp dụng, hệ đơn vị và kiểm tra tính hợp lí. Nhận định “Đồ thị $F$ theo $I$ luôn là đường cong và không phụ thuộc B.” sai vì trái với kiến thức nêu trên.
}
\end{bt}

% ===== Câu 287 =====
% ID: L12C3B15-68-TLN
% Mức: TH
\begin{ex}
% ID: L12C3B15-68-TLN
% Mức: TH
Trong bốn phát biểu về Bài toán tổng hợp lực từ gắn với đồ thị, thí nghiệm và thực tiễn, có bao nhiêu phát biểu đúng? (Chỉ ghi một số nguyên.) (1) Với $B$ và l không đổi, đồ thị $F$ theo $I$ là đường thẳng qua gốc có hệ số góc Bl sin(alpha). (2) Cần kiểm tra điều kiện áp dụng. (3) Đồ thị $F$ theo $I$ luôn là đường cong và không phụ thuộc B. (4) Có thể bỏ qua đơn vị của kết quả.
\shortans{2}
\loigiai{
Đối chiếu từng phát biểu với kiến thức: Với $B$ và l không đổi, đồ thị $F$ theo $I$ là đường thẳng qua gốc có hệ số góc Bl sin(alpha). Số phát biểu đúng là 2.
}
\end{ex}

% ===== Câu 288 =====
% ID: L12C3B15-68-TN
% Mức: TH
\dangbt{Nhận biết lực từ và đặc trưng của vectơ cảm ứng từ}
\begin{ex}
% ID: L12C3B15-68-TN
% Mức: TH
Chọn phát biểu không trái với kiến thức về Nhận biết lực từ và đặc trưng của vectơ cảm ứng từ?
\choice
{Kết quả không phụ thuộc dữ kiện và điều kiện áp dụng.}
{\True Vectơ cảm ứng từ tại một điểm có hướng trùng với hướng của từ trường tại điểm đó; đơn vị là tesla (T).}
{Cảm ứng từ là đại lượng vô hướng và có đơn vị niutơn.}
{Có thể bỏ qua phương, chiều và đơn vị của các đại lượng trong mọi trường hợp.}
\loigiai{
Vectơ cảm ứng từ tại một điểm có hướng trùng với hướng của từ trường tại điểm đó; đơn vị là tesla (T). Vì vậy phương án $B$ đúng; các phương án còn lại trái với kiến thức hoặc điều kiện áp dụng.
}
\end{ex}

% ===== Câu 289 =====
% ID: L12C3B15-69-TL
% Mức: VD
\dangbt{Bài toán tổng hợp lực từ gắn với đồ thị, thí nghiệm và thực tiễn}
\begin{bt}
% ID: L12C3B15-69-TL
% Mức: VD
Phân tích sai lầm trong nhận định: “Đồ thị $F$ theo $I$ luôn là đường cong và không phụ thuộc B.”
\loigiai{
Cần nêu được: Với $B$ và l không đổi, đồ thị $F$ theo $I$ là đường thẳng qua gốc có hệ số góc Bl sin(alpha). Khi vận dụng phải xác định đúng dữ kiện, phương chiều nếu có, điều kiện áp dụng, hệ đơn vị và kiểm tra tính hợp lí. Nhận định “Đồ thị $F$ theo $I$ luôn là đường cong và không phụ thuộc B.” sai vì trái với kiến thức nêu trên.
}
\end{bt}

% ===== Câu 290 =====
% ID: L12C3B15-69-TLN
% Mức: TH
\begin{ex}
% ID: L12C3B15-69-TLN
% Mức: TH
Trong bốn phát biểu về Bài toán tổng hợp lực từ gắn với đồ thị, thí nghiệm và thực tiễn, có bao nhiêu phát biểu đúng? (Chỉ ghi một số nguyên.) (1) Đồ thị $F$ theo $I$ luôn là đường cong và không phụ thuộc B. (2) Với $B$ và l không đổi, đồ thị $F$ theo $I$ là đường thẳng qua gốc có hệ số góc Bl sin(alpha). (3) Đồ thị $F$ theo $I$ luôn là đường cong và không phụ thuộc B. (4) Có thể bỏ qua đơn vị của kết quả.
\shortans{1}
\loigiai{
Đối chiếu từng phát biểu với kiến thức: Với $B$ và l không đổi, đồ thị $F$ theo $I$ là đường thẳng qua gốc có hệ số góc Bl sin(alpha). Số phát biểu đúng là 1.
}
\end{ex}

% ===== Câu 291 =====
% ID: L12C3B15-69-TN
% Mức: VD
\dangbt{Nhận biết lực từ và đặc trưng của vectơ cảm ứng từ}
\begin{ex}
% ID: L12C3B15-69-TN
% Mức: VD
Cơ sở Vật lí đúng dùng để giải Nhận biết lực từ và đặc trưng của vectơ cảm ứng từ?
\choice
{Có thể bỏ qua phương, chiều và đơn vị của các đại lượng trong mọi trường hợp.}
{Kết quả không phụ thuộc dữ kiện và điều kiện áp dụng.}
{\True Vectơ cảm ứng từ tại một điểm có hướng trùng với hướng của từ trường tại điểm đó; đơn vị là tesla (T).}
{Cảm ứng từ là đại lượng vô hướng và có đơn vị niutơn.}
\loigiai{
Vectơ cảm ứng từ tại một điểm có hướng trùng với hướng của từ trường tại điểm đó; đơn vị là tesla (T). Vì vậy phương án $C$ đúng; các phương án còn lại trái với kiến thức hoặc điều kiện áp dụng.
}
\end{ex}

% ===== Câu 292 =====
% ID: L12C3B15-70-TL
% Mức: VD
\dangbt{Bài toán tổng hợp lực từ gắn với đồ thị, thí nghiệm và thực tiễn}
\begin{bt}
% ID: L12C3B15-70-TL
% Mức: VD
Đề xuất các bước giải một bài toán thuộc dạng Bài toán tổng hợp lực từ gắn với đồ thị, thí nghiệm và thực tiễn.
\loigiai{
Cần nêu được: Với $B$ và l không đổi, đồ thị $F$ theo $I$ là đường thẳng qua gốc có hệ số góc Bl sin(alpha). Khi vận dụng phải xác định đúng dữ kiện, phương chiều nếu có, điều kiện áp dụng, hệ đơn vị và kiểm tra tính hợp lí. Nhận định “Đồ thị $F$ theo $I$ luôn là đường cong và không phụ thuộc B.” sai vì trái với kiến thức nêu trên.
}
\end{bt}

% ===== Câu 293 =====
% ID: L12C3B15-70-TLN
% Mức: TH
\begin{ex}
% ID: L12C3B15-70-TLN
% Mức: TH
Trong bốn phát biểu về Bài toán tổng hợp lực từ gắn với đồ thị, thí nghiệm và thực tiễn, có bao nhiêu phát biểu đúng? (Chỉ ghi một số nguyên.) (1) Với $B$ và l không đổi, đồ thị $F$ theo $I$ là đường thẳng qua gốc có hệ số góc Bl sin(alpha). (2) Cần kiểm tra điều kiện áp dụng. (3) Đồ thị $F$ theo $I$ luôn là đường cong và không phụ thuộc B. (4) Có thể bỏ qua đơn vị của kết quả.
\shortans{2}
\loigiai{
Đối chiếu từng phát biểu với kiến thức: Với $B$ và l không đổi, đồ thị $F$ theo $I$ là đường thẳng qua gốc có hệ số góc Bl sin(alpha). Số phát biểu đúng là 2.
}
\end{ex}

% ===== Câu 294 =====
% ID: L12C3B15-70-TN
% Mức: VD
\dangbt{Nhận biết lực từ và đặc trưng của vectơ cảm ứng từ}
\begin{ex}
% ID: L12C3B15-70-TN
% Mức: VD
Trong một bài thực hành về Nhận biết lực từ và đặc trưng của vectơ cảm ứng từ?
\choice
{Cảm ứng từ là đại lượng vô hướng và có đơn vị niutơn.}
{Có thể bỏ qua phương, chiều và đơn vị của các đại lượng trong mọi trường hợp.}
{Kết quả không phụ thuộc dữ kiện và điều kiện áp dụng.}
{\True Vectơ cảm ứng từ tại một điểm có hướng trùng với hướng của từ trường tại điểm đó; đơn vị là tesla (T).}
\loigiai{
Vectơ cảm ứng từ tại một điểm có hướng trùng với hướng của từ trường tại điểm đó; đơn vị là tesla (T). Vì vậy phương án $D$ đúng; các phương án còn lại trái với kiến thức hoặc điều kiện áp dụng.
}
\end{ex}

% ===== Câu 295 =====
% ID: L12C3B15-71-TL
% Mức: VD
\dangbt{Bài toán tổng hợp lực từ gắn với đồ thị, thí nghiệm và thực tiễn}
\begin{bt}
% ID: L12C3B15-71-TL
% Mức: VD
Nêu một tình huống thực tiễn liên quan đến Bài toán tổng hợp lực từ gắn với đồ thị, thí nghiệm và thực tiễn và giải thích bằng kiến thức Vật lí.
\loigiai{
Cần nêu được: Với $B$ và l không đổi, đồ thị $F$ theo $I$ là đường thẳng qua gốc có hệ số góc Bl sin(alpha). Khi vận dụng phải xác định đúng dữ kiện, phương chiều nếu có, điều kiện áp dụng, hệ đơn vị và kiểm tra tính hợp lí. Nhận định “Đồ thị $F$ theo $I$ luôn là đường cong và không phụ thuộc B.” sai vì trái với kiến thức nêu trên.
}
\end{bt}

% ===== Câu 296 =====
% ID: L12C3B15-71-TLN
% Mức: VD
\begin{ex}
% ID: L12C3B15-71-TLN
% Mức: VD
Trong bốn phát biểu về Bài toán tổng hợp lực từ gắn với đồ thị, thí nghiệm và thực tiễn, có bao nhiêu phát biểu đúng? (Chỉ ghi một số nguyên.) (1) Đồ thị $F$ theo $I$ luôn là đường cong và không phụ thuộc B. (2) Với $B$ và l không đổi, đồ thị $F$ theo $I$ là đường thẳng qua gốc có hệ số góc Bl sin(alpha). (3) Đồ thị $F$ theo $I$ luôn là đường cong và không phụ thuộc B. (4) Có thể bỏ qua đơn vị của kết quả.
\shortans{2}
\loigiai{
Đối chiếu từng phát biểu với kiến thức: Với $B$ và l không đổi, đồ thị $F$ theo $I$ là đường thẳng qua gốc có hệ số góc Bl sin(alpha). Số phát biểu đúng là 2.
}
\end{ex}

% ===== Câu 297 =====
% ID: L12C3B15-71-TN
% Mức: VD
\dangbt{Nhận biết lực từ và đặc trưng của vectơ cảm ứng từ}
\begin{ex}
% ID: L12C3B15-71-TN
% Mức: VD
Kết luận đúng khi vận dụng Nhận biết lực từ và đặc trưng của vectơ cảm ứng từ?
\choice
{\True Vectơ cảm ứng từ tại một điểm có hướng trùng với hướng của từ trường tại điểm đó; đơn vị là tesla (T).}
{Cảm ứng từ là đại lượng vô hướng và có đơn vị niutơn.}
{Có thể bỏ qua phương, chiều và đơn vị của các đại lượng trong mọi trường hợp.}
{Kết quả không phụ thuộc dữ kiện và điều kiện áp dụng.}
\loigiai{
Vectơ cảm ứng từ tại một điểm có hướng trùng với hướng của từ trường tại điểm đó; đơn vị là tesla (T). Vì vậy phương án $A$ đúng; các phương án còn lại trái với kiến thức hoặc điều kiện áp dụng.
}
\end{ex}

% ===== Câu 298 =====
% ID: L12C3B15-72-TL
% Mức: VDC
\dangbt{Bài toán tổng hợp lực từ gắn với đồ thị, thí nghiệm và thực tiễn}
\begin{bt}
% ID: L12C3B15-72-TL
% Mức: VDC
So sánh nhận định đúng “Với $B$ và l không đổi, đồ thị $F$ theo $I$ là đường thẳng qua gốc có hệ số góc Bl sin(alpha).” với nhận định sai “Đồ thị $F$ theo $I$ luôn là đường cong và không phụ thuộc B.”; chỉ rõ căn cứ Vật lí.
\loigiai{
Cần nêu được: Với $B$ và l không đổi, đồ thị $F$ theo $I$ là đường thẳng qua gốc có hệ số góc Bl sin(alpha). Khi vận dụng phải xác định đúng dữ kiện, phương chiều nếu có, điều kiện áp dụng, hệ đơn vị và kiểm tra tính hợp lí. Nhận định “Đồ thị $F$ theo $I$ luôn là đường cong và không phụ thuộc B.” sai vì trái với kiến thức nêu trên.
}
\end{bt}

% ===== Câu 299 =====
% ID: L12C3B15-72-TLN
% Mức: VD
\begin{ex}
% ID: L12C3B15-72-TLN
% Mức: VD
Trong bốn phát biểu về Bài toán tổng hợp lực từ gắn với đồ thị, thí nghiệm và thực tiễn, có bao nhiêu phát biểu đúng? (Chỉ ghi một số nguyên.) (1) Với $B$ và l không đổi, đồ thị $F$ theo $I$ là đường thẳng qua gốc có hệ số góc Bl sin(alpha). (2) Cần kiểm tra điều kiện áp dụng. (3) Đồ thị $F$ theo $I$ luôn là đường cong và không phụ thuộc B. (4) Có thể bỏ qua đơn vị của kết quả.
\shortans{2}
\loigiai{
Đối chiếu từng phát biểu với kiến thức: Với $B$ và l không đổi, đồ thị $F$ theo $I$ là đường thẳng qua gốc có hệ số góc Bl sin(alpha). Số phát biểu đúng là 2.
}
\end{ex}

% ===== Câu 300 =====
% ID: L12C3B15-72-TN
% Mức: NB
\dangbt{Xác định phương, chiều lực từ bằng quy tắc bàn tay trái}
\begin{ex}
% ID: L12C3B15-72-TN
% Mức: NB
Phát biểu nào sau đây đúng về Xác định phương, chiều lực từ bằng quy tắc bàn tay trái?
\choice
{Lực từ luôn cùng chiều dòng điện trong dây.}
{Có thể bỏ qua phương, chiều và đơn vị của các đại lượng trong mọi trường hợp.}
{Kết quả không phụ thuộc dữ kiện và điều kiện áp dụng.}
{\True Lực từ tác dụng lên đoạn dây có phương vuông góc với cả dòng điện và vectơ cảm ứng từ; chiều xác định bằng quy tắc bàn tay trái.}
\loigiai{
Lực từ tác dụng lên đoạn dây có phương vuông góc với cả dòng điện và vectơ cảm ứng từ; chiều xác định bằng quy tắc bàn tay trái. Vì vậy phương án $D$ đúng; các phương án còn lại trái với kiến thức hoặc điều kiện áp dụng.
}
\end{ex}

% ===== Câu 301 =====
% ID: L12C3B15-73-TLN
% Mức: VDC
\dangbt{Bài toán tổng hợp lực từ gắn với đồ thị, thí nghiệm và thực tiễn}
\begin{ex}
% ID: L12C3B15-73-TLN
% Mức: VDC
Trong bốn phát biểu về Bài toán tổng hợp lực từ gắn với đồ thị, thí nghiệm và thực tiễn, có bao nhiêu phát biểu đúng? (Chỉ ghi một số nguyên.) (1) Đồ thị $F$ theo $I$ luôn là đường cong và không phụ thuộc B. (2) Với $B$ và l không đổi, đồ thị $F$ theo $I$ là đường thẳng qua gốc có hệ số góc Bl sin(alpha). (3) Phải dùng đúng định luật cho tình huống. (4) Với $B$ và l không đổi, đồ thị $F$ theo $I$ là đường thẳng qua gốc có hệ số góc Bl sin(alpha).
\shortans{3}
\loigiai{
Đối chiếu từng phát biểu với kiến thức: Với $B$ và l không đổi, đồ thị $F$ theo $I$ là đường thẳng qua gốc có hệ số góc Bl sin(alpha). Số phát biểu đúng là 3.
}
\end{ex}

% ===== Câu 302 =====
% ID: L12C3B15-73-TN
% Mức: NB
\dangbt{Xác định phương, chiều lực từ bằng quy tắc bàn tay trái}
\begin{ex}
% ID: L12C3B15-73-TN
% Mức: NB
Trong nội dung Xác định phương, chiều lực từ bằng quy tắc bàn tay trái?
\choice
{\True Lực từ tác dụng lên đoạn dây có phương vuông góc với cả dòng điện và vectơ cảm ứng từ; chiều xác định bằng quy tắc bàn tay trái.}
{Lực từ luôn cùng chiều dòng điện trong dây.}
{Có thể bỏ qua phương, chiều và đơn vị của các đại lượng trong mọi trường hợp.}
{Kết quả không phụ thuộc dữ kiện và điều kiện áp dụng.}
\loigiai{
Lực từ tác dụng lên đoạn dây có phương vuông góc với cả dòng điện và vectơ cảm ứng từ; chiều xác định bằng quy tắc bàn tay trái. Vì vậy phương án $A$ đúng; các phương án còn lại trái với kiến thức hoặc điều kiện áp dụng.
}
\end{ex}

% ===== Câu 303 =====
% ID: L12C3B15-74-TN
% Mức: NB
\begin{ex}
% ID: L12C3B15-74-TN
% Mức: NB
Tình huống 3: chọn kết luận chính xác về Xác định phương, chiều lực từ bằng quy tắc bàn tay trái?
\choice
{Kết quả không phụ thuộc dữ kiện và điều kiện áp dụng.}
{\True Lực từ tác dụng lên đoạn dây có phương vuông góc với cả dòng điện và vectơ cảm ứng từ; chiều xác định bằng quy tắc bàn tay trái.}
{Lực từ luôn cùng chiều dòng điện trong dây.}
{Có thể bỏ qua phương, chiều và đơn vị của các đại lượng trong mọi trường hợp.}
\loigiai{
Lực từ tác dụng lên đoạn dây có phương vuông góc với cả dòng điện và vectơ cảm ứng từ; chiều xác định bằng quy tắc bàn tay trái. Vì vậy phương án $B$ đúng; các phương án còn lại trái với kiến thức hoặc điều kiện áp dụng.
}
\end{ex}

% ===== Câu 304 =====
% ID: L12C3B15-75-TN
% Mức: TH
\begin{ex}
% ID: L12C3B15-75-TN
% Mức: TH
Nội dung nào mô tả đúng nhất Xác định phương, chiều lực từ bằng quy tắc bàn tay trái?
\choice
{Có thể bỏ qua phương, chiều và đơn vị của các đại lượng trong mọi trường hợp.}
{Kết quả không phụ thuộc dữ kiện và điều kiện áp dụng.}
{\True Lực từ tác dụng lên đoạn dây có phương vuông góc với cả dòng điện và vectơ cảm ứng từ; chiều xác định bằng quy tắc bàn tay trái.}
{Lực từ luôn cùng chiều dòng điện trong dây.}
\loigiai{
Lực từ tác dụng lên đoạn dây có phương vuông góc với cả dòng điện và vectơ cảm ứng từ; chiều xác định bằng quy tắc bàn tay trái. Vì vậy phương án $C$ đúng; các phương án còn lại trái với kiến thức hoặc điều kiện áp dụng.
}
\end{ex}

% ===== Câu 305 =====
% ID: L12C3B15-76-TN
% Mức: TH
\begin{ex}
% ID: L12C3B15-76-TN
% Mức: TH
Khi phân tích Xác định phương, chiều lực từ bằng quy tắc bàn tay trái?
\choice
{Lực từ luôn cùng chiều dòng điện trong dây.}
{Có thể bỏ qua phương, chiều và đơn vị của các đại lượng trong mọi trường hợp.}
{Kết quả không phụ thuộc dữ kiện và điều kiện áp dụng.}
{\True Lực từ tác dụng lên đoạn dây có phương vuông góc với cả dòng điện và vectơ cảm ứng từ; chiều xác định bằng quy tắc bàn tay trái.}
\loigiai{
Lực từ tác dụng lên đoạn dây có phương vuông góc với cả dòng điện và vectơ cảm ứng từ; chiều xác định bằng quy tắc bàn tay trái. Vì vậy phương án $D$ đúng; các phương án còn lại trái với kiến thức hoặc điều kiện áp dụng.
}
\end{ex}

% ===== Câu 306 =====
% ID: L12C3B15-77-TN
% Mức: TH
\begin{ex}
% ID: L12C3B15-77-TN
% Mức: TH
Một học sinh ôn tập Xác định phương, chiều lực từ bằng quy tắc bàn tay trái?
\choice
{\True Lực từ tác dụng lên đoạn dây có phương vuông góc với cả dòng điện và vectơ cảm ứng từ; chiều xác định bằng quy tắc bàn tay trái.}
{Lực từ luôn cùng chiều dòng điện trong dây.}
{Có thể bỏ qua phương, chiều và đơn vị của các đại lượng trong mọi trường hợp.}
{Kết quả không phụ thuộc dữ kiện và điều kiện áp dụng.}
\loigiai{
Lực từ tác dụng lên đoạn dây có phương vuông góc với cả dòng điện và vectơ cảm ứng từ; chiều xác định bằng quy tắc bàn tay trái. Vì vậy phương án $A$ đúng; các phương án còn lại trái với kiến thức hoặc điều kiện áp dụng.
}
\end{ex}

% ===== Câu 307 =====
% ID: L12C3B15-78-TN
% Mức: TH
\begin{ex}
% ID: L12C3B15-78-TN
% Mức: TH
Chọn phát biểu không trái với kiến thức về Xác định phương, chiều lực từ bằng quy tắc bàn tay trái?
\choice
{Kết quả không phụ thuộc dữ kiện và điều kiện áp dụng.}
{\True Lực từ tác dụng lên đoạn dây có phương vuông góc với cả dòng điện và vectơ cảm ứng từ; chiều xác định bằng quy tắc bàn tay trái.}
{Lực từ luôn cùng chiều dòng điện trong dây.}
{Có thể bỏ qua phương, chiều và đơn vị của các đại lượng trong mọi trường hợp.}
\loigiai{
Lực từ tác dụng lên đoạn dây có phương vuông góc với cả dòng điện và vectơ cảm ứng từ; chiều xác định bằng quy tắc bàn tay trái. Vì vậy phương án $B$ đúng; các phương án còn lại trái với kiến thức hoặc điều kiện áp dụng.
}
\end{ex}

% ===== Câu 308 =====
% ID: L12C3B15-79-TN
% Mức: VD
\begin{ex}
% ID: L12C3B15-79-TN
% Mức: VD
Cơ sở Vật lí đúng dùng để giải Xác định phương, chiều lực từ bằng quy tắc bàn tay trái?
\choice
{Có thể bỏ qua phương, chiều và đơn vị của các đại lượng trong mọi trường hợp.}
{Kết quả không phụ thuộc dữ kiện và điều kiện áp dụng.}
{\True Lực từ tác dụng lên đoạn dây có phương vuông góc với cả dòng điện và vectơ cảm ứng từ; chiều xác định bằng quy tắc bàn tay trái.}
{Lực từ luôn cùng chiều dòng điện trong dây.}
\loigiai{
Lực từ tác dụng lên đoạn dây có phương vuông góc với cả dòng điện và vectơ cảm ứng từ; chiều xác định bằng quy tắc bàn tay trái. Vì vậy phương án $C$ đúng; các phương án còn lại trái với kiến thức hoặc điều kiện áp dụng.
}
\end{ex}

% ===== Câu 309 =====
% ID: L12C3B15-80-TN
% Mức: VD
\begin{ex}
% ID: L12C3B15-80-TN
% Mức: VD
Trong một bài thực hành về Xác định phương, chiều lực từ bằng quy tắc bàn tay trái?
\choice
{Lực từ luôn cùng chiều dòng điện trong dây.}
{Có thể bỏ qua phương, chiều và đơn vị của các đại lượng trong mọi trường hợp.}
{Kết quả không phụ thuộc dữ kiện và điều kiện áp dụng.}
{\True Lực từ tác dụng lên đoạn dây có phương vuông góc với cả dòng điện và vectơ cảm ứng từ; chiều xác định bằng quy tắc bàn tay trái.}
\loigiai{
Lực từ tác dụng lên đoạn dây có phương vuông góc với cả dòng điện và vectơ cảm ứng từ; chiều xác định bằng quy tắc bàn tay trái. Vì vậy phương án $D$ đúng; các phương án còn lại trái với kiến thức hoặc điều kiện áp dụng.
}
\end{ex}

% ===== Câu 310 =====
% ID: L12C3B15-81-TN
% Mức: VD
\begin{ex}
% ID: L12C3B15-81-TN
% Mức: VD
Kết luận đúng khi vận dụng Xác định phương, chiều lực từ bằng quy tắc bàn tay trái?
\choice
{\True Lực từ tác dụng lên đoạn dây có phương vuông góc với cả dòng điện và vectơ cảm ứng từ; chiều xác định bằng quy tắc bàn tay trái.}
{Lực từ luôn cùng chiều dòng điện trong dây.}
{Có thể bỏ qua phương, chiều và đơn vị của các đại lượng trong mọi trường hợp.}
{Kết quả không phụ thuộc dữ kiện và điều kiện áp dụng.}
\loigiai{
Lực từ tác dụng lên đoạn dây có phương vuông góc với cả dòng điện và vectơ cảm ứng từ; chiều xác định bằng quy tắc bàn tay trái. Vì vậy phương án $A$ đúng; các phương án còn lại trái với kiến thức hoặc điều kiện áp dụng.
}
\end{ex}

% ===== Câu 311 =====
% ID: L12C3B15-82-TN
% Mức: NB
\dangbt{Tính lực từ tác dụng lên đoạn dây dẫn mang dòng điện}
\begin{ex}
% ID: L12C3B15-82-TN
% Mức: NB
Phát biểu nào sau đây đúng về Tính lực từ tác dụng lên đoạn dây dẫn mang dòng điện?
\choice
{Lực từ được tính bằng $F$ = BI/l và không phụ thuộc góc.}
{Có thể bỏ qua phương, chiều và đơn vị của các đại lượng trong mọi trường hợp.}
{Kết quả không phụ thuộc dữ kiện và điều kiện áp dụng.}
{\True Độ lớn lực từ lên đoạn dây thẳng dài l mang dòng $I$ trong từ trường đều là $F$ = BIl sin(alpha).}
\loigiai{
Độ lớn lực từ lên đoạn dây thẳng dài l mang dòng $I$ trong từ trường đều là $F$ = BIl sin(alpha). Vì vậy phương án $D$ đúng; các phương án còn lại trái với kiến thức hoặc điều kiện áp dụng.
}
\end{ex}

% ===== Câu 312 =====
% ID: L12C3B15-83-TN
% Mức: NB
\begin{ex}
% ID: L12C3B15-83-TN
% Mức: NB
Trong nội dung Tính lực từ tác dụng lên đoạn dây dẫn mang dòng điện?
\choice
{\True Độ lớn lực từ lên đoạn dây thẳng dài l mang dòng $I$ trong từ trường đều là $F$ = BIl sin(alpha).}
{Lực từ được tính bằng $F$ = BI/l và không phụ thuộc góc.}
{Có thể bỏ qua phương, chiều và đơn vị của các đại lượng trong mọi trường hợp.}
{Kết quả không phụ thuộc dữ kiện và điều kiện áp dụng.}
\loigiai{
Độ lớn lực từ lên đoạn dây thẳng dài l mang dòng $I$ trong từ trường đều là $F$ = BIl sin(alpha). Vì vậy phương án $A$ đúng; các phương án còn lại trái với kiến thức hoặc điều kiện áp dụng.
}
\end{ex}

% ===== Câu 313 =====
% ID: L12C3B15-84-TN
% Mức: NB
\begin{ex}
% ID: L12C3B15-84-TN
% Mức: NB
Tình huống 3: chọn kết luận chính xác về Tính lực từ tác dụng lên đoạn dây dẫn mang dòng điện?
\choice
{Kết quả không phụ thuộc dữ kiện và điều kiện áp dụng.}
{\True Độ lớn lực từ lên đoạn dây thẳng dài l mang dòng $I$ trong từ trường đều là $F$ = BIl sin(alpha).}
{Lực từ được tính bằng $F$ = BI/l và không phụ thuộc góc.}
{Có thể bỏ qua phương, chiều và đơn vị của các đại lượng trong mọi trường hợp.}
\loigiai{
Độ lớn lực từ lên đoạn dây thẳng dài l mang dòng $I$ trong từ trường đều là $F$ = BIl sin(alpha). Vì vậy phương án $B$ đúng; các phương án còn lại trái với kiến thức hoặc điều kiện áp dụng.
}
\end{ex}

% ===== Câu 314 =====
% ID: L12C3B15-85-TN
% Mức: TH
\begin{ex}
% ID: L12C3B15-85-TN
% Mức: TH
Nội dung nào mô tả đúng nhất Tính lực từ tác dụng lên đoạn dây dẫn mang dòng điện?
\choice
{Có thể bỏ qua phương, chiều và đơn vị của các đại lượng trong mọi trường hợp.}
{Kết quả không phụ thuộc dữ kiện và điều kiện áp dụng.}
{\True Độ lớn lực từ lên đoạn dây thẳng dài l mang dòng $I$ trong từ trường đều là $F$ = BIl sin(alpha).}
{Lực từ được tính bằng $F$ = BI/l và không phụ thuộc góc.}
\loigiai{
Độ lớn lực từ lên đoạn dây thẳng dài l mang dòng $I$ trong từ trường đều là $F$ = BIl sin(alpha). Vì vậy phương án $C$ đúng; các phương án còn lại trái với kiến thức hoặc điều kiện áp dụng.
}
\end{ex}

% ===== Câu 315 =====
% ID: L12C3B15-86-TN
% Mức: TH
\begin{ex}
% ID: L12C3B15-86-TN
% Mức: TH
Khi phân tích Tính lực từ tác dụng lên đoạn dây dẫn mang dòng điện?
\choice
{Lực từ được tính bằng $F$ = BI/l và không phụ thuộc góc.}
{Có thể bỏ qua phương, chiều và đơn vị của các đại lượng trong mọi trường hợp.}
{Kết quả không phụ thuộc dữ kiện và điều kiện áp dụng.}
{\True Độ lớn lực từ lên đoạn dây thẳng dài l mang dòng $I$ trong từ trường đều là $F$ = BIl sin(alpha).}
\loigiai{
Độ lớn lực từ lên đoạn dây thẳng dài l mang dòng $I$ trong từ trường đều là $F$ = BIl sin(alpha). Vì vậy phương án $D$ đúng; các phương án còn lại trái với kiến thức hoặc điều kiện áp dụng.
}
\end{ex}

% ===== Câu 316 =====
% ID: L12C3B15-87-TN
% Mức: TH
\begin{ex}
% ID: L12C3B15-87-TN
% Mức: TH
Một học sinh ôn tập Tính lực từ tác dụng lên đoạn dây dẫn mang dòng điện?
\choice
{\True Độ lớn lực từ lên đoạn dây thẳng dài l mang dòng $I$ trong từ trường đều là $F$ = BIl sin(alpha).}
{Lực từ được tính bằng $F$ = BI/l và không phụ thuộc góc.}
{Có thể bỏ qua phương, chiều và đơn vị của các đại lượng trong mọi trường hợp.}
{Kết quả không phụ thuộc dữ kiện và điều kiện áp dụng.}
\loigiai{
Độ lớn lực từ lên đoạn dây thẳng dài l mang dòng $I$ trong từ trường đều là $F$ = BIl sin(alpha). Vì vậy phương án $A$ đúng; các phương án còn lại trái với kiến thức hoặc điều kiện áp dụng.
}
\end{ex}

% ===== Câu 317 =====
% ID: L12C3B15-88-TN
% Mức: TH
\begin{ex}
% ID: L12C3B15-88-TN
% Mức: TH
Chọn phát biểu không trái với kiến thức về Tính lực từ tác dụng lên đoạn dây dẫn mang dòng điện?
\choice
{Kết quả không phụ thuộc dữ kiện và điều kiện áp dụng.}
{\True Độ lớn lực từ lên đoạn dây thẳng dài l mang dòng $I$ trong từ trường đều là $F$ = BIl sin(alpha).}
{Lực từ được tính bằng $F$ = BI/l và không phụ thuộc góc.}
{Có thể bỏ qua phương, chiều và đơn vị của các đại lượng trong mọi trường hợp.}
\loigiai{
Độ lớn lực từ lên đoạn dây thẳng dài l mang dòng $I$ trong từ trường đều là $F$ = BIl sin(alpha). Vì vậy phương án $B$ đúng; các phương án còn lại trái với kiến thức hoặc điều kiện áp dụng.
}
\end{ex}

% ===== Câu 318 =====
% ID: L12C3B15-89-TN
% Mức: VD
\begin{ex}
% ID: L12C3B15-89-TN
% Mức: VD
Cơ sở Vật lí đúng dùng để giải Tính lực từ tác dụng lên đoạn dây dẫn mang dòng điện?
\choice
{Có thể bỏ qua phương, chiều và đơn vị của các đại lượng trong mọi trường hợp.}
{Kết quả không phụ thuộc dữ kiện và điều kiện áp dụng.}
{\True Độ lớn lực từ lên đoạn dây thẳng dài l mang dòng $I$ trong từ trường đều là $F$ = BIl sin(alpha).}
{Lực từ được tính bằng $F$ = BI/l và không phụ thuộc góc.}
\loigiai{
Độ lớn lực từ lên đoạn dây thẳng dài l mang dòng $I$ trong từ trường đều là $F$ = BIl sin(alpha). Vì vậy phương án $C$ đúng; các phương án còn lại trái với kiến thức hoặc điều kiện áp dụng.
}
\end{ex}

% ===== Câu 319 =====
% ID: L12C3B15-90-TN
% Mức: VD
\begin{ex}
% ID: L12C3B15-90-TN
% Mức: VD
Trong một bài thực hành về Tính lực từ tác dụng lên đoạn dây dẫn mang dòng điện?
\choice
{Lực từ được tính bằng $F$ = BI/l và không phụ thuộc góc.}
{Có thể bỏ qua phương, chiều và đơn vị của các đại lượng trong mọi trường hợp.}
{Kết quả không phụ thuộc dữ kiện và điều kiện áp dụng.}
{\True Độ lớn lực từ lên đoạn dây thẳng dài l mang dòng $I$ trong từ trường đều là $F$ = BIl sin(alpha).}
\loigiai{
Độ lớn lực từ lên đoạn dây thẳng dài l mang dòng $I$ trong từ trường đều là $F$ = BIl sin(alpha). Vì vậy phương án $D$ đúng; các phương án còn lại trái với kiến thức hoặc điều kiện áp dụng.
}
\end{ex}

% ===== Câu 320 =====
% ID: L12C3B15-91-TN
% Mức: VD
\begin{ex}
% ID: L12C3B15-91-TN
% Mức: VD
Kết luận đúng khi vận dụng Tính lực từ tác dụng lên đoạn dây dẫn mang dòng điện?
\choice
{\True Độ lớn lực từ lên đoạn dây thẳng dài l mang dòng $I$ trong từ trường đều là $F$ = BIl sin(alpha).}
{Lực từ được tính bằng $F$ = BI/l và không phụ thuộc góc.}
{Có thể bỏ qua phương, chiều và đơn vị của các đại lượng trong mọi trường hợp.}
{Kết quả không phụ thuộc dữ kiện và điều kiện áp dụng.}
\loigiai{
Độ lớn lực từ lên đoạn dây thẳng dài l mang dòng $I$ trong từ trường đều là $F$ = BIl sin(alpha). Vì vậy phương án $A$ đúng; các phương án còn lại trái với kiến thức hoặc điều kiện áp dụng.
}
\end{ex}

% ===== Câu 321 =====
% ID: L12C3B15-92-TN
% Mức: NB
\dangbt{Xác định cảm ứng từ từ phép đo lực từ}
\begin{ex}
% ID: L12C3B15-92-TN
% Mức: NB
Phát biểu nào sau đây đúng về Xác định cảm ứng từ từ phép đo lực từ?
\choice
{Khi dây vuông góc từ trường, cảm ứng từ bằng FIl.}
{Có thể bỏ qua phương, chiều và đơn vị của các đại lượng trong mọi trường hợp.}
{Kết quả không phụ thuộc dữ kiện và điều kiện áp dụng.}
{\True Khi dây vuông góc từ trường, cảm ứng từ được xác định bởi $B$ = F/(Il).}
\loigiai{
Khi dây vuông góc từ trường, cảm ứng từ được xác định bởi $B$ = F/(Il). Vì vậy phương án $D$ đúng; các phương án còn lại trái với kiến thức hoặc điều kiện áp dụng.
}
\end{ex}

% ===== Câu 322 =====
% ID: L12C3B15-93-TN
% Mức: NB
\begin{ex}
% ID: L12C3B15-93-TN
% Mức: NB
Trong nội dung Xác định cảm ứng từ từ phép đo lực từ?
\choice
{\True Khi dây vuông góc từ trường, cảm ứng từ được xác định bởi $B$ = F/(Il).}
{Khi dây vuông góc từ trường, cảm ứng từ bằng FIl.}
{Có thể bỏ qua phương, chiều và đơn vị của các đại lượng trong mọi trường hợp.}
{Kết quả không phụ thuộc dữ kiện và điều kiện áp dụng.}
\loigiai{
Khi dây vuông góc từ trường, cảm ứng từ được xác định bởi $B$ = F/(Il). Vì vậy phương án $A$ đúng; các phương án còn lại trái với kiến thức hoặc điều kiện áp dụng.
}
\end{ex}

% ===== Câu 323 =====
% ID: L12C3B15-94-TN
% Mức: NB
\begin{ex}
% ID: L12C3B15-94-TN
% Mức: NB
Tình huống 3: chọn kết luận chính xác về Xác định cảm ứng từ từ phép đo lực từ?
\choice
{Kết quả không phụ thuộc dữ kiện và điều kiện áp dụng.}
{\True Khi dây vuông góc từ trường, cảm ứng từ được xác định bởi $B$ = F/(Il).}
{Khi dây vuông góc từ trường, cảm ứng từ bằng FIl.}
{Có thể bỏ qua phương, chiều và đơn vị của các đại lượng trong mọi trường hợp.}
\loigiai{
Khi dây vuông góc từ trường, cảm ứng từ được xác định bởi $B$ = F/(Il). Vì vậy phương án $B$ đúng; các phương án còn lại trái với kiến thức hoặc điều kiện áp dụng.
}
\end{ex}

% ===== Câu 324 =====
% ID: L12C3B15-95-TN
% Mức: TH
\begin{ex}
% ID: L12C3B15-95-TN
% Mức: TH
Nội dung nào mô tả đúng nhất Xác định cảm ứng từ từ phép đo lực từ?
\choice
{Có thể bỏ qua phương, chiều và đơn vị của các đại lượng trong mọi trường hợp.}
{Kết quả không phụ thuộc dữ kiện và điều kiện áp dụng.}
{\True Khi dây vuông góc từ trường, cảm ứng từ được xác định bởi $B$ = F/(Il).}
{Khi dây vuông góc từ trường, cảm ứng từ bằng FIl.}
\loigiai{
Khi dây vuông góc từ trường, cảm ứng từ được xác định bởi $B$ = F/(Il). Vì vậy phương án $C$ đúng; các phương án còn lại trái với kiến thức hoặc điều kiện áp dụng.
}
\end{ex}

% ===== Câu 325 =====
% ID: L12C3B15-96-TN
% Mức: TH
\begin{ex}
% ID: L12C3B15-96-TN
% Mức: TH
Khi phân tích Xác định cảm ứng từ từ phép đo lực từ?
\choice
{Khi dây vuông góc từ trường, cảm ứng từ bằng FIl.}
{Có thể bỏ qua phương, chiều và đơn vị của các đại lượng trong mọi trường hợp.}
{Kết quả không phụ thuộc dữ kiện và điều kiện áp dụng.}
{\True Khi dây vuông góc từ trường, cảm ứng từ được xác định bởi $B$ = F/(Il).}
\loigiai{
Khi dây vuông góc từ trường, cảm ứng từ được xác định bởi $B$ = F/(Il). Vì vậy phương án $D$ đúng; các phương án còn lại trái với kiến thức hoặc điều kiện áp dụng.
}
\end{ex}

% ===== Câu 326 =====
% ID: L12C3B15-97-TN
% Mức: TH
\begin{ex}
% ID: L12C3B15-97-TN
% Mức: TH
Một học sinh ôn tập Xác định cảm ứng từ từ phép đo lực từ?
\choice
{\True Khi dây vuông góc từ trường, cảm ứng từ được xác định bởi $B$ = F/(Il).}
{Khi dây vuông góc từ trường, cảm ứng từ bằng FIl.}
{Có thể bỏ qua phương, chiều và đơn vị của các đại lượng trong mọi trường hợp.}
{Kết quả không phụ thuộc dữ kiện và điều kiện áp dụng.}
\loigiai{
Khi dây vuông góc từ trường, cảm ứng từ được xác định bởi $B$ = F/(Il). Vì vậy phương án $A$ đúng; các phương án còn lại trái với kiến thức hoặc điều kiện áp dụng.
}
\end{ex}

% ===== Câu 327 =====
% ID: L12C3B15-98-TN
% Mức: TH
\begin{ex}
% ID: L12C3B15-98-TN
% Mức: TH
Chọn phát biểu không trái với kiến thức về Xác định cảm ứng từ từ phép đo lực từ?
\choice
{Kết quả không phụ thuộc dữ kiện và điều kiện áp dụng.}
{\True Khi dây vuông góc từ trường, cảm ứng từ được xác định bởi $B$ = F/(Il).}
{Khi dây vuông góc từ trường, cảm ứng từ bằng FIl.}
{Có thể bỏ qua phương, chiều và đơn vị của các đại lượng trong mọi trường hợp.}
\loigiai{
Khi dây vuông góc từ trường, cảm ứng từ được xác định bởi $B$ = F/(Il). Vì vậy phương án $B$ đúng; các phương án còn lại trái với kiến thức hoặc điều kiện áp dụng.
}
\end{ex}

% ===== Câu 328 =====
% ID: L12C3B15-99-TN
% Mức: VD
\begin{ex}
% ID: L12C3B15-99-TN
% Mức: VD
Cơ sở Vật lí đúng dùng để giải Xác định cảm ứng từ từ phép đo lực từ?
\choice
{Có thể bỏ qua phương, chiều và đơn vị của các đại lượng trong mọi trường hợp.}
{Kết quả không phụ thuộc dữ kiện và điều kiện áp dụng.}
{\True Khi dây vuông góc từ trường, cảm ứng từ được xác định bởi $B$ = F/(Il).}
{Khi dây vuông góc từ trường, cảm ứng từ bằng FIl.}
\loigiai{
Khi dây vuông góc từ trường, cảm ứng từ được xác định bởi $B$ = F/(Il). Vì vậy phương án $C$ đúng; các phương án còn lại trái với kiến thức hoặc điều kiện áp dụng.
}
\end{ex}
